% Copyright 2021 by an anonymous marmot
%
% This file may be distributed and/or modified
%
% 1. under the LaTeX Project Public License and/or
% 2. under the GNU Public License.
%
% See the file doc/generic/pgf/licenses/LICENSE for more details.
\ProvidesFileRCS{tikzlibrary3dtools.code.tex}
\usetikzlibrary{3d,calc,decorations,fpu}% to do: ability to switch on and off
\makeatletter%
\ifnum\month>10\relax
\typeout{Hibernation season has just started.}%
\fi
\ifnum\month<5\relax
\typeout{Hibernation season has not yet ended.}%
\fi
\ifnum\month=2\relax
\ifnum\day=2
\typeout{Yay! Groundhog Day!}%
\fi
\fi
\def\tikz@td@extractpgfversionaux#1.#2.#3|#4#5#6;{\def#4{#1}\def#5{#2}\def#6{#3}}%
\def\tikz@td@checkversion{%
\edef\pgfutil@tmp{\noexpand\tikz@td@extractpgfversionaux\pgfversion|\noexpand\pgfutil@tmpa\noexpand\pgfutil@tmpb\noexpand\pgfutil@tmpc;}%
\pgfutil@tmp%\typeout{a=\pgfutil@tmpa,b=\pgfutil@tmpb,c=\pgfutil@tmpc}%
\pgfmathtruncatemacro{\pgfutil@tmpe}{ifthenelse(10*\pgfutil@tmpa+\pgfutil@tmpb<31,0,1)}%
\ifnum\pgfutil@tmpe=0%
\message{You are using a too old version of pgf (\pgfversion). You need at least
version 3.1.1 to use the features of this library.}%
\fi
\edef\pgfutil@tmpd{\expandafter\@firstoftwo\pgfutil@tmpc\empty}%
\pgfmathtruncatemacro{\pgfutil@tmpe}{ifthenelse(100*\pgfutil@tmpa+10*\pgfutil@tmpb+\pgfutil@tmpc<313,0,1)}%
\ifnum\pgfutil@tmpe=0%
\message{Some routines of 3dtools may not work because they rely on the pushing
and popping techniques that were added in version 3.1.3 (but the current version
is \pgfversion).^^J}%
\fi
\pgfmathtruncatemacro{\pgfutil@tmpe}{ifthenelse(100*\pgfutil@tmpa+10*\pgfutil@tmpb+\pgfutil@tmpc<316,0,1)}%
\ifnum\pgfutil@tmpe=0%
\def\pgfutil@pushmacro##1{%
    \xdef\pgfutil@pushmacro@string{\string##1}%
    \ifcsname pgfutil@pushedmacro@\pgfutil@pushmacro@string\endcsname\else
        % \newcount is \outer in Plain
        \csname newcount\expandafter\endcsname\csname pgfutil@pushedmacro@\pgfutil@pushmacro@string\endcsname
    \fi
    \global\advance\csname pgfutil@pushedmacro@\pgfutil@pushmacro@string\endcsname 1\relax
    \global\expandafter\let\csname\the\csname pgfutil@pushedmacro@\pgfutil@pushmacro@string\endcsname\pgfutil@pushmacro@string\endcsname##1%
}
\message{pgfutil@pushmacro fixed because current pgf version \pgfversion\space is smaller than 3.1.6.^^J}%
\fi
}%
\tikz@td@checkversion%
%  cos(#1)*cos(#2) 
%  & sin(#1)*cos(#2) 
%  & -sin(#2) \\
%  cos(#1)*sin(#2)*sin(#3)-sin(#1)*cos(#3) 
%  &  sin(#1)*sin(#2)*sin(#3)+cos(#1)*cos(#3) 
%  & cos(#2)*sin(#3) \\
%  cos(#1)*sin(#2)*cos(#3)+sin(#1)*sin(#3) 
%  & sin(#1)*sin(#2)*cos(#3)-cos(#1)*sin(#3) 
%  & cos(#2)*cos(#3) \\ %<-normal on screen
% 
\newcommand{\orthmat}[3]{% the entries of this matrix keep track
\pgfmathparse{cos(#1)*cos(#2)}% of the current transformation
\xdef\tikz@td@matAA{\pgfmathresult}%
\pgfmathparse{sin(#1)*cos(#2)}% 
\xdef\tikz@td@matAB{\pgfmathresult}% 
\pgfmathparse{-sin(#2)}% 
\xdef\tikz@td@matAC{\pgfmathresult}%
\pgfmathparse{cos(#1)*sin(#2)*sin(#3)-sin(#1)*cos(#3)}%
\xdef\tikz@td@matBA{\pgfmathresult}%  
\pgfmathparse{sin(#1)*sin(#2)*sin(#3)+cos(#1)*cos(#3)}% 
\xdef\tikz@td@matBB{\pgfmathresult}%
\pgfmathparse{cos(#2)*sin(#3)}%
\xdef\tikz@td@matBC{\pgfmathresult}%
\pgfmathparse{cos(#1)*sin(#2)*cos(#3)+sin(#1)*sin(#3) }%
\xdef\tikz@td@matCA{\pgfmathresult}% 
\pgfmathparse{sin(#1)*sin(#2)*cos(#3)-cos(#1)*sin(#3)}% 
\xdef\tikz@td@matCB{\pgfmathresult}%
\pgfmathparse{cos(#2)*cos(#3)}%
\xdef\tikz@td@matCC{\pgfmathresult}}%
\tikzset{3d/.cd,phi/.initial=0,psi/.initial=0,theta/.initial=0,
install view/.style={/utils/exec=\tikzset{3d/.cd,#1}%
\orthmat{\pgfkeysvalueof{/tikz/3d/phi}}{%
\pgfkeysvalueof{/tikz/3d/psi}}{\pgfkeysvalueof{/tikz/3d/theta}},%
/tikz/x={({\tikz@td@matAA*1cm},{\tikz@td@matBA*1cm})},%
/tikz/y={({\tikz@td@matAB*1cm},{\tikz@td@matBB*1cm})},%
/tikz/z={({\tikz@td@matAC*1cm},{\tikz@td@matBC*1cm})}}}%
\def\pgfmathparse@td@FPU#1{\begingroup%
\pgfkeys{/pgf/fpu,/pgf/fpu/output format=fixed}%
\pgfmathparse{#1}%
\pgfmathsmuggle\pgfmathresult\endgroup}%
\def\pgfmathsetmacro@td@FPU#1#2{\begingroup%
\pgfkeys{/pgf/fpu,/pgf/fpu/output format=fixed}%
\pgfmathsetmacro#1{#2}%
\pgfmathsmuggle#1\endgroup}%
%
% ---------------------------------------------------------------------------
% Numeric back end (expl3/l3fp, l3sort). Requires the LaTeX kernel >= 2020-02
% (expl3 preloaded); the library is LaTeX-only anyway (\newcommand, \newcounter,
% \PackageWarning, \begin{scope}, ...). Public pgfmath names, the TD
% mini-language and all keys are unchanged.
\ExplSyntaxOn
\scan_new:N \s__tdtools_stop
\scan_new:N \s__tdtools_mark
\fp_new:N \l__tdtools_ax_fp \fp_new:N \l__tdtools_ay_fp \fp_new:N \l__tdtools_az_fp
\fp_new:N \l__tdtools_bx_fp \fp_new:N \l__tdtools_by_fp \fp_new:N \l__tdtools_bz_fp
\fp_new:N \l__tdtools_cx_fp \fp_new:N \l__tdtools_cy_fp \fp_new:N \l__tdtools_cz_fp
\fp_new:N \l__tdtools_rx_fp \fp_new:N \l__tdtools_ry_fp \fp_new:N \l__tdtools_rz_fp
\fp_new:N \l__tdtools_rs_fp
\fp_new:N \l__tdtools_f_fp
\fp_new:N \l__tdtools_t_fp
\int_new:N \l__tdtools_type_int
\tl_new:N \l__tdtools_tmp_tl
\tl_new:N \l__tdtools_c_tl
\tl_new:N \l__tdtools_result_tl
\seq_new:N \l__tdtools_comp_seq
\clist_new:N \l__tdtools_tmp_clist
\fp_const:Nn \c__tdtools_cm_fp { \dim_to_fp:n { 1cm } }
\fp_const:Nn \c__tdtools_cmsq_fp { \c__tdtools_cm_fp * \c__tdtools_cm_fp }
\msg_new:nnn { 3dtools } { bad-coordinate }
  { Cannot~interpret~`(#1)'~as~a~3d~coordinate.~Using~(0,0,0). }
\msg_new:nnn { 3dtools } { syntax }
  { Cannot~parse~`#1'~in~TD~expression~(missing~`#2'). }
\msg_new:nnn { 3dtools } { trailing }
  { Ignoring~unparsed~material~`#1'~in~TD~expression. }
% ---------------------------------------------------------------------------
% Numbers. pgfmath hands declared functions either plain decimals or, when the
% fpu is active, its internal <flag>Y<mantissa>e<exp>] format.
\cs_new:Npn \__tdtools_num:n #1 { \__tdtools_num:w #1 Y \q_nil Y \q_stop }
\cs_new:Npn \__tdtools_num:w #1 Y #2 Y #3 \q_stop
  { \quark_if_nil:nTF {#2} {#1} { \__tdtools_num_float:nw {#1} #2 \q_stop } }
\cs_new:Npn \__tdtools_num_float:nw #1 #2 ] #3 \q_stop
  {
    \int_case:nnF {#1}
      { {0} {0} {1} {#2} {2} {-#2} {3} {nan} {4} {inf} {5} {-inf} } {nan}
  }
\cs_generate_variant:Nn \__tdtools_num:n { V }
% Evaluate a pgfmath expression into an fp variable. Material made only of
% digits, spaces and . + - * / ( ) (same syntax and precedence in pgfmath and
% l3fp) goes straight to l3fp; everything else (pgfmath functions in degrees,
% units, macros that do not expand to numbers, user-declared functions, ...)
% still goes through the pgf fpu, exactly as before.
\cs_new:Npn \__tdtools_class:N #1
  {
    \if_int_compare:w `#1 < 48 \exp_stop_f:
      \if_int_compare:w `#1 < 40 \exp_stop_f:
        \if_int_compare:w `#1 = 32 \exp_stop_f: \else: X \fi:
      \else:
        \if_int_compare:w `#1 = 44 \exp_stop_f: X \fi:
      \fi:
    \else:
      \if_int_compare:w `#1 > 57 \exp_stop_f: X \else: 1 \fi:
    \fi:
  }
% true iff the classified string has a digit and no X
\cs_new:Npn \__tdtools_if_arith:w #1 X #2 \q_stop
  { \tl_if_empty:nTF {#2} { \tl_if_empty:nTF {#1} { \use_ii:nn } { \use_i:nn } } { \use_ii:nn } }
\cs_new_protected:Npn \__tdtools_eval:Nn #1#2
  {
    \tl_set:Ne \l__tdtools_tmp_tl {#2}
    \exp_last_unbraced:Ne \__tdtools_if_arith:w
      { \exp_args:NV \str_map_function:nN \l__tdtools_tmp_tl \__tdtools_class:N }
      X \q_stop
      { \fp_set:Nn #1 { \l__tdtools_tmp_tl } }
      {
        \group_begin:
          \pgfkeys { /pgf/fpu , /pgf/fpu/output~format = float }
          \pgfmathparse { \l__tdtools_tmp_tl }
          \exp_args:NNNe
        \group_end:
        \fp_set:Nn #1 { \__tdtools_num:V \pgfmathresult }
      }
  }
% ---------------------------------------------------------------------------
% Coordinates: (x,y), (x,y,z) or (name). Names are resolved recursively
% (a name may stand for another name); (X) coordinate (X) resolves to (0,0,0);
% an unknown name gives a warning and (0,0,0). #1 is unused (kept for callers).
\int_new:N \l__tdtools_depth_int
\cs_new_protected:Npn \__tdtools_coord:nn #1#2
  {
    \tl_set:Ne \l__tdtools_c_tl {#2}
    \seq_set_split:NnV \l__tdtools_comp_seq { , } \l__tdtools_c_tl
    \int_case:nnF { \seq_count:N \l__tdtools_comp_seq }
      {
        { 1 }
          {
            \exp_args:Ne \__tdtools_coord_named:n
              { \seq_item:Nn \l__tdtools_comp_seq { 1 } }
          }
        { 2 }
          {
            \seq_pop_left:NN \l__tdtools_comp_seq \l__tdtools_tmp_tl
            \__tdtools_eval:NV \l__tdtools_cx_fp \l__tdtools_tmp_tl
            \seq_pop_left:NN \l__tdtools_comp_seq \l__tdtools_tmp_tl
            \__tdtools_eval:NV \l__tdtools_cy_fp \l__tdtools_tmp_tl
            \fp_zero:N \l__tdtools_cz_fp
          }
        { 3 }
          {
            \seq_pop_left:NN \l__tdtools_comp_seq \l__tdtools_tmp_tl
            \__tdtools_eval:NV \l__tdtools_cx_fp \l__tdtools_tmp_tl
            \seq_pop_left:NN \l__tdtools_comp_seq \l__tdtools_tmp_tl
            \__tdtools_eval:NV \l__tdtools_cy_fp \l__tdtools_tmp_tl
            \seq_pop_left:NN \l__tdtools_comp_seq \l__tdtools_tmp_tl
            \__tdtools_eval:NV \l__tdtools_cz_fp \l__tdtools_tmp_tl
          }
      }
      {
        \msg_error:nnn { 3dtools } { bad-coordinate } {#2}
        \__tdtools_coord_zero:
      }
  }
\cs_generate_variant:Nn \__tdtools_eval:Nn { NV }
\cs_generate_variant:Nn \seq_set_split:Nnn { NnV }
\cs_new_protected:Npn \__tdtools_coord_zero:
  { \fp_zero:N \l__tdtools_cx_fp \fp_zero:N \l__tdtools_cy_fp \fp_zero:N \l__tdtools_cz_fp }
\cs_new_protected:Npn \__tdtools_coord_unknown:n #1
  {
    \PackageWarning{3dtools}{Unknown~coordinate~`#1'.~Using~(0,0,0).}
    \__tdtools_coord_zero:
  }
\cs_new_protected:Npn \__tdtools_coord_named:n #1
  {
    \cs_if_exist:cTF { tikz@td@n@ #1 }
      {
        \tikz@td@lookup {#1}
        \int_case:nn { \tikz@td@lkmode }
          {
            { 0 } { \__tdtools_coord_text:n {#1} }
            { 1 }
              {
                \clist_set:NV \l__tdtools_tmp_clist \tikz@td@lkv
                \fp_set:Nn \l__tdtools_cx_fp { \clist_item:Nn \l__tdtools_tmp_clist { 1 } }
                \fp_set:Nn \l__tdtools_cy_fp { \clist_item:Nn \l__tdtools_tmp_clist { 2 } }
                \fp_set:Nn \l__tdtools_cz_fp { \clist_item:Nn \l__tdtools_tmp_clist { 3 } }
              }
            { 2 } { \tikz@td@unknownwarning {#1} \__tdtools_coord_text:n {#1} }
          }
      }
      { \__tdtools_coord_text:n {#1} }
  }
% the text a coordinate was defined with (\tikz@dcl@coord@<name>)
\cs_new_protected:Npn \__tdtools_coord_text:n #1
  {
    \cs_if_exist:cTF { tikz@dcl@coord@ #1 }
      {
        \exp_last_unbraced:Nv \__tdtools_coord_named:w { tikz@dcl@coord@ #1 }
          ( ) \s__tdtools_stop {#1}
      }
      { \__tdtools_coord_unknown:n {#1} }
  }
\cs_new_protected:Npn \__tdtools_coord_named:w #1 ( #2 ) #3 \s__tdtools_stop #4
  {
    \tl_if_blank:nTF {#1}
      {
        \tl_if_eq:nnTF {#2} {#4}
          { \__tdtools_coord_zero: } % (X) coordinate (X)
          {
            \int_compare:nNnTF { \l__tdtools_depth_int } > { 64 }
              { \__tdtools_coord_unknown:n {#4} } % cyclic definitions
              {
                \group_begin:
                  \int_incr:N \l__tdtools_depth_int
                  \__tdtools_coord:nn { } {#2}
                  \use:e
                    {
                \group_end:
                      \fp_set:Nn \exp_not:N \l__tdtools_cx_fp { \fp_use:N \l__tdtools_cx_fp }
                      \fp_set:Nn \exp_not:N \l__tdtools_cy_fp { \fp_use:N \l__tdtools_cy_fp }
                      \fp_set:Nn \exp_not:N \l__tdtools_cz_fp { \fp_use:N \l__tdtools_cz_fp }
                    }
              }
          }
      }
      { \__tdtools_coord_unknown:n {#4} }
  }
% ---------------------------------------------------------------------------
% TD mini-language: same grammar as \pgfmathtdparse (calc-like).
\cs_new_protected:Npn \__tdtools_parse:n #1
  {
    \fp_zero:N \l__tdtools_ax_fp \fp_zero:N \l__tdtools_ay_fp \fp_zero:N \l__tdtools_az_fp
    \int_zero:N \l__tdtools_type_int
    \fp_set:Nn \l__tdtools_f_fp { 1 }
    \__tdtools_dispatch:w + #1 \s__tdtools_mark \s__tdtools_stop
  }
% Delimited macros below never scan past the end of the TD string: check that
% the delimiter #1 occurs before \s__tdtools_stop, else report and stop.
\cs_new_protected:Npn \__tdtools_guard:nNw #1#2#3 \s__tdtools_stop
  {
    \tl_if_in:nnTF {#3} {#1}
      { #2 #3 \s__tdtools_stop }
      {
        \__tdtools_syntax_error:nw {#1} #3
        \__tdtools_end:w \s__tdtools_mark \s__tdtools_stop
      }
  }
\cs_new_protected:Npn \__tdtools_syntax_error:nw #1#2 \s__tdtools_mark
  { \msg_error:nnnn { 3dtools } { syntax } {#2} {#1} }
\cs_new_protected:Npn \__tdtools_dispatch:w
  {
    \pgfutil@ifnextchar + \__tdtools_add:w
      {
        \pgfutil@ifnextchar - \__tdtools_sub:w
          {
            \pgfutil@ifnextchar x \__tdtools_vecprod:w
              {
                \pgfutil@ifnextchar o \__tdtools_scalprod:w
                  {
                    \pgfutil@ifnextchar [ { \__tdtools_guard:nNw { ] } \__tdtools_bracket:w }
                      {
                        \pgfutil@ifnextchar > \__tdtools_gt:w \__tdtools_end:w
                      }
                  }
              }
          }
      }
  }
\cs_new_protected:Npn \__tdtools_factor_or_coord:
  {
    \pgfutil@ifnextchar ( { \__tdtools_guard:nNw { ) } \__tdtools_coordinate:w }
      {
        \pgfutil@ifnextchar [ { \__tdtools_guard:nNw { ] } \__tdtools_bracket:w }
          {
            \pgfutil@ifnextchar - \__tdtools_negate:w
              { \__tdtools_guard:nNw { *( } \__tdtools_factor:w }
          }
      }
  }
\cs_new_protected:Npn \__tdtools_negate:w - % unary minus (F11)
  { \fp_set:Nn \l__tdtools_f_fp { - \l__tdtools_f_fp } \__tdtools_factor_or_coord: }
\cs_new_protected:Npn \__tdtools_add:w +
  { \fp_set:Nn \l__tdtools_f_fp { 1 } \__tdtools_factor_or_coord: }
\cs_new_protected:Npn \__tdtools_sub:w -
  { \fp_set:Nn \l__tdtools_f_fp { -1 } \__tdtools_factor_or_coord: }
\cs_new_protected:Npn \__tdtools_bracket:w [ #1 ]
  {
    \__tdtools_eval:Nn \l__tdtools_t_fp {#1}
    \fp_set:Nn \l__tdtools_f_fp { \l__tdtools_f_fp * \l__tdtools_t_fp }
    \__tdtools_factor_or_coord:
  }
\cs_new_protected:Npn \__tdtools_factor:w #1 *(
  {
    \tl_if_blank:nF {#1} % `[expr]*(...)' as used in the manual
      {
        \__tdtools_eval:Nn \l__tdtools_t_fp {#1}
        \fp_set:Nn \l__tdtools_f_fp { \l__tdtools_t_fp * \l__tdtools_f_fp }
      }
    \__tdtools_coordinate:w (
  }
\cs_new_protected:Npn \__tdtools_coordinate:w ( #1 )
  {
    \__tdtools_coord:nn { \c_false_bool } {#1}
    \fp_set:Nn \l__tdtools_ax_fp { \l__tdtools_ax_fp + \l__tdtools_f_fp * \l__tdtools_cx_fp }
    \fp_set:Nn \l__tdtools_ay_fp { \l__tdtools_ay_fp + \l__tdtools_f_fp * \l__tdtools_cy_fp }
    \fp_set:Nn \l__tdtools_az_fp { \l__tdtools_az_fp + \l__tdtools_f_fp * \l__tdtools_cz_fp }
    \__tdtools_dispatch:w
  }
\cs_new_protected:Npn \__tdtools_shift:n #1
  {
    \int_set:Nn \l__tdtools_type_int {#1}
    \fp_set_eq:NN \l__tdtools_bx_fp \l__tdtools_ax_fp
    \fp_set_eq:NN \l__tdtools_by_fp \l__tdtools_ay_fp
    \fp_set_eq:NN \l__tdtools_bz_fp \l__tdtools_az_fp
    \fp_zero:N \l__tdtools_ax_fp \fp_zero:N \l__tdtools_ay_fp \fp_zero:N \l__tdtools_az_fp
    \fp_set:Nn \l__tdtools_f_fp { 1 }
    \__tdtools_factor_or_coord:
  }
\cs_new_protected:Npn \__tdtools_vecprod:w x { \__tdtools_shift:n { 1 } }
\cs_new_protected:Npn \__tdtools_scalprod:w o { \__tdtools_shift:n { 2 } }
\cs_new_protected:Npn \__tdtools_gt:w > { \__tdtools_end:w } % old end marker
\cs_new_protected:Npn \__tdtools_end:w #1 \s__tdtools_mark \s__tdtools_stop
  {
    \tl_if_blank:nF {#1} { \msg_warning:nnn { 3dtools } { trailing } {#1} }
    \int_case:nn { \l__tdtools_type_int }
      {
        { 0 }
          {
            \fp_set_eq:NN \l__tdtools_rx_fp \l__tdtools_ax_fp
            \fp_set_eq:NN \l__tdtools_ry_fp \l__tdtools_ay_fp
            \fp_set_eq:NN \l__tdtools_rz_fp \l__tdtools_az_fp
          }
        { 1 } % b x a, as \vpauxx etc.
          {
            \fp_set:Nn \l__tdtools_rx_fp
              { \l__tdtools_by_fp * \l__tdtools_az_fp - \l__tdtools_bz_fp * \l__tdtools_ay_fp }
            \fp_set:Nn \l__tdtools_ry_fp
              { \l__tdtools_ax_fp * \l__tdtools_bz_fp - \l__tdtools_bx_fp * \l__tdtools_az_fp }
            \fp_set:Nn \l__tdtools_rz_fp
              { \l__tdtools_bx_fp * \l__tdtools_ay_fp - \l__tdtools_by_fp * \l__tdtools_ax_fp }
          }
        { 2 }
          {
            \fp_set:Nn \l__tdtools_rs_fp
              {
                \l__tdtools_ax_fp * \l__tdtools_bx_fp
                + \l__tdtools_ay_fp * \l__tdtools_by_fp
                + \l__tdtools_az_fp * \l__tdtools_bz_fp
              }
          }
      }
  }
\cs_new:Npn \__tdtools_vec_tl:
  {
    \fp_to_decimal:N \l__tdtools_rx_fp ,
    \fp_to_decimal:N \l__tdtools_ry_fp ,
    \fp_to_decimal:N \l__tdtools_rz_fp
  }
\cs_new_protected:Npn \__tdtools_return:e #1
  {
    \tl_set:Ne \l__tdtools_result_tl {#1}
    \exp_args:NNNV
    \group_end:
    \tl_set:Nn \pgfmathresult \l__tdtools_result_tl
  }
% ---------------------------------------------------------------------------
% pgfmath-facing functions (same names/arity as the original library)
\pgfmathdeclarefunction{TD}{1}
  {
    \group_begin:
    \__tdtools_parse:n {#1}
    \__tdtools_return:e
      {
        \int_compare:nNnTF { \l__tdtools_type_int } = { 2 }
          { \fp_to_decimal:N \l__tdtools_rs_fp } { \__tdtools_vec_tl: }
      }
  }
\pgfmathdeclarefunction{TDunit}{1}
  {
    \group_begin:
    \__tdtools_parse:n {#1}
    \fp_set:Nn \l__tdtools_t_fp
      {
        sqrt ( \l__tdtools_rx_fp * \l__tdtools_rx_fp + \l__tdtools_ry_fp * \l__tdtools_ry_fp
               + \l__tdtools_rz_fp * \l__tdtools_rz_fp )
      }
    \fp_compare:nNnTF { min ( 1 , \l__tdtools_t_fp ) } < { 0.002 }
      { \PackageWarning{3dtools}{Vector~too~short,~cannot~normalize~it.} }
      {
        \fp_set:Nn \l__tdtools_rx_fp { \l__tdtools_rx_fp / \l__tdtools_t_fp }
        \fp_set:Nn \l__tdtools_ry_fp { \l__tdtools_ry_fp / \l__tdtools_t_fp }
        \fp_set:Nn \l__tdtools_rz_fp { \l__tdtools_rz_fp / \l__tdtools_t_fp }
      }
    \__tdtools_return:e { \__tdtools_vec_tl: }
  }
% components of a 3d expression; a bare name A means the coordinate (A)
\cs_new_protected:Npn \__tdtools_component:nN #1#2
  {
    \group_begin:
    \tl_if_in:nnTF {#1} { ( }
      { \__tdtools_parse:n {#1} }
      { \__tdtools_parse:n { (#1) } }
    \__tdtools_return:e { \fp_to_decimal:N #2 }
  }
\pgfmathdeclarefunction{TDx}{1} { \__tdtools_component:nN {#1} \l__tdtools_rx_fp }
\pgfmathdeclarefunction{TDy}{1} { \__tdtools_component:nN {#1} \l__tdtools_ry_fp }
\pgfmathdeclarefunction{TDz}{1} { \__tdtools_component:nN {#1} \l__tdtools_rz_fp }
\pgfmathdeclarefunction{tddistance}{2}
  {
    \group_begin:
    \__tdtools_parse:n { #1 - #2 o #1 - #2 }
    \__tdtools_return:e { \fp_to_decimal:n { sqrt ( \l__tdtools_rs_fp ) } }
  }
% screen normal from the current x/y/z unit vectors (as in the original)
\cs_new:Npn \__tdtools_nscreen:NNNN #1#2#3#4
  { ( #1 * #2 - #3 * #4 ) / \c__tdtools_cmsq_fp }
\cs_new:Npn \__tdtools_nx: { \__tdtools_nscreen:NNNN \pgf@yx \pgf@zy \pgf@yy \pgf@zx }
\cs_new:Npn \__tdtools_ny: { \__tdtools_nscreen:NNNN \pgf@zx \pgf@xy \pgf@xx \pgf@zy }
\cs_new:Npn \__tdtools_nz: { \__tdtools_nscreen:NNNN \pgf@xx \pgf@yy \pgf@yx \pgf@xy }
\pgfmathdeclarefunction{screendepth}{3}
  {
    \tl_set:Ne \pgfmathresult
      {
        \fp_to_decimal:n
          {
            \__tdtools_nx: * ( \__tdtools_num:n {#1} )
            + \__tdtools_ny: * ( \__tdtools_num:n {#2} )
            + \__tdtools_nz: * ( \__tdtools_num:n {#3} )
          }
      }
  }
\pgfmathdeclarefunction*{nscreenx}{0}
  { \tl_set:Ne \pgfmathresult { \fp_to_decimal:n { \__tdtools_nx: } } }
\pgfmathdeclarefunction*{nscreeny}{0}
  { \tl_set:Ne \pgfmathresult { \fp_to_decimal:n { \__tdtools_ny: } } }
\pgfmathdeclarefunction*{nscreenz}{0}
  { \tl_set:Ne \pgfmathresult { \fp_to_decimal:n { \__tdtools_nz: } } }
% internal entry points used elsewhere in the library (interfaces unchanged;
% \tikz@td@parse@coord touches no \pgfutil@tmp* macros)
\cs_new_protected:Npn \pgfmathtdparse #1
  {
    \group_begin:
    \__tdtools_parse:n {#1}
    \__tdtools_return:e
      {
        \int_compare:nNnTF { \l__tdtools_type_int } = { 2 }
          { \fp_to_decimal:N \l__tdtools_rs_fp } { \__tdtools_vec_tl: }
      }
  }
\cs_new_protected:Npn \tikz@td@parse@coord (#1)
  {
    \__tdtools_coord:nn { \c_false_bool } {#1}
    \tl_set:Ne \pgf@X { \fp_to_decimal:N \l__tdtools_cx_fp pt }
    \tl_set:Ne \pgf@Y { \fp_to_decimal:N \l__tdtools_cy_fp pt }
    \tl_set:Ne \pgf@Z { \fp_to_decimal:N \l__tdtools_cz_fp pt }
  }
% ---------------------------------------------------------------------------
% .sort numeric list: l3sort (merge sort, stable) instead of bubble sort.
\seq_new:N \l__tdtools_sort_seq
\int_new:N \l__tdtools_sort_int
\tl_new:N \l__tdtools_cmp_tl
\tl_new:N \l__tdtools_vals_tl
\tl_new:N \l__tdtools_perm_tl
\cs_new_protected:Npn \__tdtools_sort_numeric:NN #1#2
  {
    \group_begin:
    \seq_clear:N \l__tdtools_sort_seq
    \int_zero:N \l__tdtools_sort_int
    \tl_set:Ne \l__tdtools_cmp_tl { \pgfkeysvalueof { /list~management/sort~comparison } }
    \tl_set:Ne \l__tdtools_tmp_tl { \pgfkeysvalueof { \pgfkeyscurrentpath } }
    \seq_set_split:NnV \l__tdtools_comp_seq { , } \l__tdtools_tmp_tl
    \seq_map_inline:Nn \l__tdtools_comp_seq
      {
        \int_incr:N \l__tdtools_sort_int
        \pgfkeys { /list~management/sort~map = ##1 }
        \seq_put_right:Ne \l__tdtools_sort_seq
          { { \pgfmathresult } { \int_use:N \l__tdtools_sort_int } }
      }
    \seq_sort:Nn \l__tdtools_sort_seq { \__tdtools_sort_cmp:nnnn ##1 ##2 }
    \tl_set:Ne \l__tdtools_vals_tl
      { \seq_map_tokens:Nn \l__tdtools_sort_seq { \__tdtools_sort_first:n } }
    \tl_set:Ne \l__tdtools_perm_tl
      { \seq_map_tokens:Nn \l__tdtools_sort_seq { \__tdtools_sort_second:n } }
    \use:e
      {
        \group_end:
        \tl_set:Nn \exp_not:N #1 { \tl_tail:V \l__tdtools_vals_tl }
        \tl_set:Nn \exp_not:N #2 { \tl_tail:V \l__tdtools_perm_tl }
      }
  }
\cs_new:Npn \__tdtools_sort_first:n #1 { \__tdtools_sort_first:nn #1 }
\cs_new:Npn \__tdtools_sort_first:nn #1#2 { , \exp_not:n {#1} }
\cs_new:Npn \__tdtools_sort_second:n #1 { \__tdtools_sort_second:nn #1 }
\cs_new:Npn \__tdtools_sort_second:nn #1#2 { , \exp_not:n {#2} }
% bubble sort kept a pair (prev, next) iff `next <cmp> prev'
\cs_new_protected:Npn \__tdtools_sort_cmp:nnnn #1#2#3#4
  {
    \fp_compare:nTF { ( \__tdtools_num:n {#3} ) \l__tdtools_cmp_tl ( \__tdtools_num:n {#1} ) }
      { \sort_return_same: } { \sort_return_swapped: }
  }
\cs_new_eq:NN \tikz@td@sort@numeric \__tdtools_sort_numeric:NN
\ExplSyntaxOff%
\long\def\tikz@td@raw@coord(#1){\csname tikz@dcl@coord@#1\endcsname}% 
\long\def\tikz@td@makesenseof@coord#1{%
\expanded{\noexpand\pgfutil@in@{(}{#1}}% 
\ifpgfutil@in@% 
\tikz@td@parse@coord#1%
\else
\tikz@td@parse@coord(#1)%
\fi
}%
% projections 
\pgfmathdeclarefunction{xcomp3}{3}{% x component of a 3-vector 
\begingroup% 
\pgfmathparse@td@FPU{#1}% 
\pgfmathsmuggle\pgfmathresult\endgroup} 
\pgfmathdeclarefunction{ycomp3}{3}{% y component of a 3-vector 
\begingroup% 
\pgfmathparse@td@FPU{#2}% 
\pgfmathsmuggle\pgfmathresult\endgroup} 
\pgfmathdeclarefunction{zcomp3}{3}{% z component of a 3-vector 
\begingroup% 
\pgfmathparse@td@FPU{#3}% 
\pgfmathsmuggle\pgfmathresult\endgroup} 
% \pgfmathdeclarefunction{TDX}{1}{% x component of a 3-vector 
% \begingroup% 
% \edef\pgfmathresult{0}%
% % \typeout{#1}
% % \edef\temp{\noexpand\pgfmathparse@td@FPU{TDx#1}}%
% % \temp
% \typeout{\pgfmathresult}%
% \pgfmathsmuggle\pgfmathresult\endgroup}%
% \pgfmathdeclarefunction{TDY}{1}{% x component of a 3-vector 
% \begingroup% 
% \edef\mycoord{\tikz@td@raw@coord#1}%
% \pgfmathparse@td@FPU{ycomp3\mycoord}% 
% \pgfmathsmuggle\pgfmathresult\endgroup}%
% \pgfmathdeclarefunction{TDZ}{1}{% x component of a 3-vector 
% \begingroup% 
% \edef\mycoord{\tikz@td@raw@coord#1}%
% \pgfmathparse@td@FPU{zcomp3\mycoord}% 
% \pgfmathsmuggle\pgfmathresult\endgroup}%
%% ---------------------------------------------------------------------------
%% Main frame, local frames, depth tracking and physical components
%%
%% The main frame is given by the x, y and z vectors in force when no local
%% frame is installed (e.g. after 3d/install view). A local frame has an origin
%% O and basis vectors ex, ey, ez, all in main coordinates; the point with local
%% coordinates (a,b,c) sits at O+a*ex+b*ey+c*ez. Vectors are stored as comma
%% lists "x,y,z" without parentheses.
%%
%% Depth tracking. With r=(xx,yx,zx), s=(xy,yy,zy) (the rows of the main
%% projection, in cm) and n=r x s, a point v of the main frame has canvas
%% position (r.v,s.v) and depth n.v (both in cm). M=(r;s;n) is invertible for
%% any non-degenerate view, so canvas position plus depth determine v. Every
%% point TikZ computes carries its depth in \tikz@td@z next to \pgf@x/\pgf@y:
%% \pgfpointxyz, \pgfpointxy and \pgfpointpolarxy set it, node anchors take the
%% depth of the node, canvas points have depth 0 (they lie in the screen plane
%% through the origin of the current coordinate system), and the affine
%% constructions (++, +, calc's +, -, factors and !t!, pos= on straight
%% segments, shift={(...)}) propagate it. \tikz@td@sd is the depth of the origin
%% of the current coordinate system. Every named node or coordinate stores
%% {D}{cx}{cy}{px}{py}{z}{ok}{exact}{{picture}{main vectors}{frame}{context}}
%% in \tikz@td@n@<name>: depth D=z+\tikz@td@sd and canvas position (cx,cy),
%% the same before the transformation (px,py,z), whether the depth is known,
%% and whether the stored text of the coordinate is the point (then TD uses the
%% text, which is exact, as long as it is evaluated in the same context).
\newif\iftikz@td@lf@active
\newif\iftikz@td@lf@orthonormalize
\newif\iftikz@td@zok\tikz@td@zoktrue % global: depth of the last point is known
\newif\iftikz@td@simple % global: the last point was a plain xyz/xy literal
\newif\iftikz@td@untracked % local: depths cannot be tracked in this scope
\newif\iftikz@td@keeptext % local: 3d coordinate, TD keeps using the text
\newdimen\tikz@td@z \newdimen\tikz@td@za
\newdimen\tikz@td@dzx \newdimen\tikz@td@dzy \newdimen\tikz@td@dzz
\newdimen\tikz@td@sd
\newdimen\tikz@td@lastz \newdimen\tikz@td@lastzsaved \newdimen\tikz@td@tstartz
\newdimen\tikz@td@cza \newdimen\tikz@td@czb
\newcount\tikz@td@frameid
\def\tikz@td@lastok{1}\def\tikz@td@lastoksaved{1}\def\tikz@td@lastsimple{0}%
\def\tikz@td@tstartok{0}\let\tikz@td@tstart\relax\def\tikz@td@lf@id{0}%
\def\tikz@td@split#1,#2,#3\relax{\def\tikz@td@c@a{#1}\def\tikz@td@c@b{#2}\def\tikz@td@c@c{#3}}%
\def\tikz@td@pointxyz#1,#2,#3\relax{\pgfpointxyz{#1}{#2}{#3}}%
\ExplSyntaxOn
% expandable helpers: scalar product of two comma lists, and the test whether a
% coordinate text consists of numbers and + - * / ( ) . , only
\cs_new:Npn \__tdtools_dot:nn #1#2 { \exp_last_unbraced:Ne \__tdtools_dot:w { #1 , #2 } \q_stop }
\cs_new:Npn \__tdtools_dot:w #1 , #2 , #3 , #4 , #5 , #6 \q_stop
  { ( #1 ) * ( #4 ) + ( #2 ) * ( #5 ) + ( #3 ) * ( #6 ) }
\cs_new:Npn \__tdtools_pure:N #1
  { \if_int_compare:w `#1 = 44 \exp_stop_f: \else: \__tdtools_class:N #1 \fi: }
\cs_new_protected:Npn \tikz@td@ifpure #1#2#3
  {
    \tl_set:Ne \l__tdtools_tmp_tl { \str_map_function:nN {#1} \__tdtools_pure:N }
    \tl_if_in:NnTF \l__tdtools_tmp_tl { X } {#3} {#2}
  }
% \tikz@td@solve{q1,q2,q3}: \tikz@td@lkv = E^{-1} M^{-1} q, with the columns
% of M^{-1} in \tikz@td@m@da/db/dc and the rows of E^{-1} in \tikz@td@lf@da/db/dc
\cs_new_protected:Npn \tikz@td@solve #1
  {
    \clist_set:Ne \l__tdtools_tmp_clist {#1}
    \tl_set:Ne \l__tdtools_tmp_tl
      {
        \__tdtools_lin:n { 1 } , \__tdtools_lin:n { 2 } , \__tdtools_lin:n { 3 }
      }
    \tl_set:Ne \tikz@td@lkv
      {
        \fp_eval:n { round ( \__tdtools_dot:nn { \tikz@td@lf@da } { \l__tdtools_tmp_tl } , 5 ) } ,
        \fp_eval:n { round ( \__tdtools_dot:nn { \tikz@td@lf@db } { \l__tdtools_tmp_tl } , 5 ) } ,
        \fp_eval:n { round ( \__tdtools_dot:nn { \tikz@td@lf@dc } { \l__tdtools_tmp_tl } , 5 ) }
      }
  }
\cs_generate_variant:Nn \clist_set:Nn { Ne }
\cs_new:Npn \__tdtools_lin:n #1
  {
    \fp_eval:n
      {
        ( \clist_item:Nn \l__tdtools_tmp_clist { 1 } ) * ( \clist_item:Nn \tikz@td@m@da {#1} )
        + ( \clist_item:Nn \l__tdtools_tmp_clist { 2 } ) * ( \clist_item:Nn \tikz@td@m@db {#1} )
        + ( \clist_item:Nn \l__tdtools_tmp_clist { 3 } ) * ( \clist_item:Nn \tikz@td@m@dc {#1} )
      }
  }
\cs_new_eq:NN \tikz@td@dot \__tdtools_dot:nn
\ExplSyntaxOff
% \tikz@td@dualbasis{v1}{v2}{v3}: \tikz@td@da, \tikz@td@db, \tikz@td@dc with
% d_i.v_j = delta_ij, i.e. the rows of the inverse of the matrix with columns
% v1, v2, v3 (=columns of the inverse of the matrix with rows v1, v2, v3).
% \tikz@td@det is the determinant; it is empty if the vectors are degenerate.
\def\tikz@td@dualbasis#1#2#3{%
\pgfmathsetmacro\tikz@td@da{TD("(#2)x(#3)")}%
\pgfmathsetmacro\tikz@td@db{TD("(#3)x(#1)")}%
\pgfmathsetmacro\tikz@td@dc{TD("(#1)x(#2)")}%
\pgfmathsetmacro\tikz@td@det{TD("(#1)o(\tikz@td@da)")}%
\ifdim\fpeval{abs(\tikz@td@det)}pt<0.00001pt\relax
\let\tikz@td@det\pgfutil@empty
\else
\edef\tikz@td@inv{\fpeval{1/\tikz@td@det}}%
\pgfmathsetmacro\tikz@td@da{TD("\tikz@td@inv*(\tikz@td@da)")}%
\pgfmathsetmacro\tikz@td@db{TD("\tikz@td@inv*(\tikz@td@db)")}%
\pgfmathsetmacro\tikz@td@dc{TD("\tikz@td@inv*(\tikz@td@dc)")}%
\fi}%
% \tikz@td@minv{xx,xy,yx,yy,zx,zy} (dimensions): projection rows \tikz@td@m@r
% and \tikz@td@m@s, normal \tikz@td@m@n and the columns \tikz@td@m@da/db/dc of
% M^{-1}, which maps (canvas x, canvas y, depth) to coordinates. Cached.
\def\tikz@td@minv#1{%
\ifcsname tikz@td@minv@#1\endcsname\else\tikz@td@minv@compute#1\relax\fi
\csname tikz@td@minv@#1\endcsname}%
\def\tikz@td@minv@compute#1,#2,#3,#4,#5,#6\relax{%
\begingroup
\edef\tikz@td@m@r{\fpeval{#1/1cm},\fpeval{#3/1cm},\fpeval{#5/1cm}}%
\edef\tikz@td@m@s{\fpeval{#2/1cm},\fpeval{#4/1cm},\fpeval{#6/1cm}}%
\pgfmathsetmacro\tikz@td@m@n{TD("(\tikz@td@m@r)x(\tikz@td@m@s)")}%
\tikz@td@dualbasis\tikz@td@m@r\tikz@td@m@s\tikz@td@m@n
\ifx\tikz@td@det\pgfutil@empty
\PackageWarning{3dtools}{The view is degenerate (x, y and z vectors are
coplanar on the screen). Depths and physical components are meaningless.}%
\def\tikz@td@da{0,0,0}\let\tikz@td@db\tikz@td@da\let\tikz@td@dc\tikz@td@da
\fi
\expandafter\xdef\csname tikz@td@minv@#1,#2,#3,#4,#5,#6\endcsname{%
\def\noexpand\tikz@td@m@r{\tikz@td@m@r}\def\noexpand\tikz@td@m@s{\tikz@td@m@s}%
\def\noexpand\tikz@td@m@n{\tikz@td@m@n}\def\noexpand\tikz@td@m@da{\tikz@td@da}%
\def\noexpand\tikz@td@m@db{\tikz@td@db}\def\noexpand\tikz@td@m@dc{\tikz@td@dc}}%
\endgroup}%
\def\tikz@td@curvec{\the\pgf@xx,\the\pgf@xy,\the\pgf@yx,\the\pgf@yy,\the\pgf@zx,\the\pgf@zy}%
\def\tikz@td@mainvec{\iftikz@td@lf@active\tikz@td@m@vecstr\else\tikz@td@curvec\fi}%
% context of a TD evaluation (transformation in \tikz@td@T, set beforehand)
\def\tikz@td@ctxkey{\tikz@td@T|\tikz@td@curvec|\the\tikz@td@sd|\tikz@td@lf@id}%
% \tikz@td@lf@getmain: unless a local frame is active (then the data stored by
% the frame are kept), store the main frame: transformation \tikz@td@m@T, the
% state \tikz@td@m@state (vectors, depth vectors, depth of the origin) and the
% vectors \tikz@td@m@vecstr; then load \tikz@td@m@r etc. (\tikz@td@minv).
\def\tikz@td@lf@getmain{%
\iftikz@td@lf@active\else
\pgfgettransform\tikz@td@m@T
\edef\tikz@td@m@vecstr{\tikz@td@curvec}%
\edef\tikz@td@m@state{\pgf@xx=\the\pgf@xx\pgf@xy=\the\pgf@xy
 \pgf@yx=\the\pgf@yx\pgf@yy=\the\pgf@yy\pgf@zx=\the\pgf@zx\pgf@zy=\the\pgf@zy
 \tikz@td@dzx=\the\tikz@td@dzx\tikz@td@dzy=\the\tikz@td@dzy
 \tikz@td@dzz=\the\tikz@td@dzz\tikz@td@sd=\the\tikz@td@sd\relax}%
\fi
\expandafter\tikz@td@minv\expandafter{\tikz@td@m@vecstr}}%
\def\tikz@td@lf@restoremain{\pgfsettransform{\tikz@td@m@T}\tikz@td@m@state}%
% identity frame; \tikz@td@lf@g holds the frame installed last (global)
\def\tikz@td@lf@identity{\def\tikz@td@lf@o{0,0,0}\def\tikz@td@lf@a{1,0,0}%
\def\tikz@td@lf@b{0,1,0}\def\tikz@td@lf@c{0,0,1}\def\tikz@td@lf@da{1,0,0}%
\def\tikz@td@lf@db{0,1,0}\def\tikz@td@lf@dc{0,0,1}}%
\let\tikz@td@lf@g\tikz@td@lf@identity
\expandafter\let\csname tikz@td@lf@tab@0\endcsname\tikz@td@lf@identity
% frame used by TDtomain and TDtolocal: the active one, otherwise the last one
\def\tikz@td@lf@load{\iftikz@td@lf@active\else\tikz@td@lf@g\fi}%
% --- depth vectors: dz_i = n_i*1cm, recomputed whenever x, y or z change
\def\tikz@td@updatedz{%
\tikz@td@dzx=\dimexpr\pgf@yx*\pgf@zy/1864679-\pgf@zx*\pgf@yy/1864679\relax
\tikz@td@dzy=\dimexpr\pgf@zx*\pgf@xy/1864679-\pgf@xx*\pgf@zy/1864679\relax
\tikz@td@dzz=\dimexpr\pgf@xx*\pgf@yy/1864679-\pgf@yx*\pgf@xy/1864679\relax}%
\tikz@td@updatedz
\let\tikz@td@orig@pgfsetxvec\pgfsetxvec
\let\tikz@td@orig@pgfsetyvec\pgfsetyvec
\let\tikz@td@orig@pgfsetzvec\pgfsetzvec
\def\pgfsetxvec#1{\tikz@td@orig@pgfsetxvec{#1}\tikz@td@updatedz\ignorespaces}%
\def\pgfsetyvec#1{\tikz@td@orig@pgfsetyvec{#1}\tikz@td@updatedz\ignorespaces}%
\def\pgfsetzvec#1{\tikz@td@orig@pgfsetzvec{#1}\tikz@td@updatedz\ignorespaces}%
% --- producers
\def\tikz@td@zreset{\global\tikz@td@z=0pt\global\tikz@td@zoktrue\global\tikz@td@simplefalse}%
\protected\def\tikz@td@zunknown{\global\tikz@td@zokfalse\global\tikz@td@simplefalse}%
\protected\def\tikz@td@setz#1#2{\global\tikz@td@z=#1\relax
\ifnum#2=0 \global\tikz@td@zokfalse\fi\global\tikz@td@simplefalse}%
\let\tikz@td@orig@pgfpointxyz\pgfpointxyz
\def\pgfpointxyz#1#2#3{\tikz@td@orig@pgfpointxyz{#1}{#2}{#3}%
\global\tikz@td@z=\pgftemp@x\tikz@td@dzx
\global\advance\tikz@td@z by\pgftemp@y\tikz@td@dzy
\global\advance\tikz@td@z by\pgftemp@z\tikz@td@dzz\global\tikz@td@simpletrue}%
\let\tikz@td@orig@pgfpointxy\pgfpointxy
\def\pgfpointxy#1#2{\tikz@td@orig@pgfpointxy{#1}{#2}%
\global\tikz@td@z=\pgftemp@x\tikz@td@dzx
\global\advance\tikz@td@z by\pgftemp@y\tikz@td@dzy\global\tikz@td@simpletrue}%
\let\tikz@td@orig@pgfpointpolarxy\pgfpointpolarxy
\def\pgfpointpolarxy#1#2{\tikz@td@orig@pgfpointpolarxy{#1}{#2}%
\global\tikz@td@z=\pgf@sys@tonumber\pgf@xa\tikz@td@dzx
\global\advance\tikz@td@z by\pgf@sys@tonumber\pgf@ya\tikz@td@dzy\global\tikz@td@simplefalse}%
% named coordinates (A) and anchors (A.north) get the depth of the node (nodes
% are flat on the screen); see \tikz@parse@node below
\def\tikz@td@anchorz#1{%
\ifcsname tikz@td@n@#1\endcsname
\expandafter\expandafter\expandafter\tikz@td@anchorz@\csname tikz@td@n@#1\endcsname
\else\tikz@td@zunknown\fi}%
\def\tikz@td@anchorz@#1#2#3#4#5#6#7#8#9{\tikz@td@anchorz@@{#1}{#7}#9}%
\def\tikz@td@anchorz@@#1#2#3#4#5#6{% D ok picture vectors frame context
\global\tikz@td@z=\dimexpr#1-\tikz@td@sd\relax\global\tikz@td@simplefalse
\ifnum#2=0 \global\tikz@td@zokfalse\fi
\def\tikz@td@tmp{#3}\ifx\tikz@td@tmp\pgfpictureid\else\global\tikz@td@zokfalse\fi}%
% --- the current point of a path carries its depth
\def\tikz@last@position{\pgfqpoint{\tikz@lastx}{\tikz@lasty}%
\tikz@td@setz\tikz@td@lastz\tikz@td@lastok}%
\def\tikz@last@position@saved{\pgfqpoint{\tikz@lastxsaved}{\tikz@lastysaved}%
\tikz@td@setz\tikz@td@lastzsaved\tikz@td@lastoksaved}%
\let\tikz@td@orig@makelast\tikz@make@last@position
\def\tikz@make@last@position#1{%
\tikz@td@zreset
\iftikz@updatecurrent\let\tikz@td@next\tikz@td@savelast\else\let\tikz@td@next\relax\fi
\tikz@td@orig@makelast{#1}%
\tikz@td@lastz=\tikz@td@z
\edef\tikz@td@lastok{\iftikz@td@zok1\else0\fi}%
\edef\tikz@td@lastsimple{\iftikz@td@simple1\else0\fi}%
\let\tikz@td@lastspec\tikz@scan@point@coordinate
\tikz@td@next}%
\def\tikz@td@savelast{\tikz@td@lastzsaved=\tikz@td@lastz\let\tikz@td@lastoksaved\tikz@td@lastok}%
% named coordinates (A) are evaluated while scanning
\let\tikz@td@orig@parsenode\tikz@parse@node
\def\tikz@parse@node#1(#2){\def\tikz@td@cont{#1}\tikz@td@nodename#2.\tikz@stop
\tikz@td@orig@parsenode\tikz@td@nodecont(#2)}%
\def\tikz@td@nodename#1.#2\tikz@stop{\def\tikz@td@tmpb{#1}}%
\def\tikz@td@nodecont#1{\tikz@td@zreset
\ifcsname tikz@td@n@\tikz@pp@name{\tikz@td@tmpb}\endcsname
\tikz@td@anchorz{\tikz@pp@name{\tikz@td@tmpb}}\else\tikz@td@anchorz{\tikz@td@tmpb}\fi
\edef\tikz@td@tmp{\noexpand\tikz@td@cont{\unexpanded{#1}%
\noexpand\tikz@td@setz{\the\tikz@td@z}{\iftikz@td@zok1\else0\fi}}}\tikz@td@tmp}%
% ++(...) and +(...)
\def\tikz@add#1{\tikz@doafter{\tikz@td@pointadd{#1}{\tikz@last@position@saved}}}%
\def\tikz@td@pointadd#1#2{%
\tikz@td@zreset\pgf@process{#1}\pgf@xa=\pgf@x\pgf@ya=\pgf@y\tikz@td@za=\tikz@td@z
\pgf@process{#2}%
\global\advance\pgf@x by\pgf@xa\global\advance\pgf@y by\pgf@ya
\global\advance\tikz@td@z by\tikz@td@za\global\tikz@td@simplefalse}%
% start of a straight segment, for pos=
\let\tikz@td@orig@@lineto\tikz@@lineto
\def\tikz@@lineto#1{%
\edef\tikz@td@tstart{\noexpand\pgfqpoint{\the\tikz@lastx}{\the\tikz@lasty}}%
\tikz@td@tstartz=\tikz@td@lastz\let\tikz@td@tstartok\tikz@td@lastok
\tikz@td@orig@@lineto{#1}}%
% shift={(...)}: the origin moves in 3d
\tikzset{shift/.code={\tikz@addtransform{\tikz@scan@one@point\tikz@td@shift#1\relax}}}%
\def\tikz@td@shift#1{\tikz@td@zreset\pgf@process{#1}%
\pgftransformshift{\pgfqpoint{\pgf@x}{\pgf@y}}%
\iftikz@td@zok\advance\tikz@td@sd by\tikz@td@z\else\tikz@td@untrackedtrue\fi}%
% ([options]...): only shifts are tracked
\def\tikz@scan@handle@options#1{%
  {%
    \pgftransformreset
    \let\tikz@transform=\pgfutil@empty
    \tikz@td@za=\tikz@td@sd
    \expandafter\tikzset\expandafter{\tikz@scan@point@options}%
    \tikz@transform
    \advance\tikz@td@za by-\tikz@td@sd
    \tikz@td@zreset\pgf@process{#1}%
    \global\advance\tikz@td@z by-\tikz@td@za
    \edef\tikz@td@tmp{\iftikz@td@zok\ifpgf@pt@identity\iftikz@td@untracked0\else1\fi
     \else0\fi\else0\fi}%
    \pgf@process{\pgfpointtransformed{\pgfqpoint{\pgf@x}{\pgf@y}}}%
    \xdef\tikz@marshal{\expandafter\noexpand\tikz@scan@point@recall{%
      \noexpand\pgfqpoint{\the\pgf@x}{\the\pgf@y}%
      \noexpand\tikz@td@setz{\the\tikz@td@z}{\tikz@td@tmp}}}%
  }%
  \tikz@marshal
}%
% perpendicular coordinates (A|-B) have no 3d meaning
\def\tikz@parse@vhdone#1(#2|-#3){%
  {%
    \tikz@@@scan@@absolute\tikz@parse@vh@mid(#2)%
    \tikz@@@scan@@absolute\tikz@parse@vh@end(#3)%
    \xdef\tikz@marshal{\noexpand#1{\noexpand\pgfqpoint{\the\pgf@xa}{\the\pgf@ya}%
      \noexpand\tikz@td@zunknown}}%
  }%
  \tikz@marshal
}%
% scopes in which depths are lost
\let\tikz@td@orig@canvasplane\tikz@canvas@is@plane
\def\tikz@canvas@is@plane{\tikz@td@untrackedtrue\tikz@td@orig@canvasplane}%
\tikzset{reset cm/.append code={\tikz@td@untrackedtrue}}%
% calc: propagate the depth through +, -, factors and !t!; everything else
% (distances, projections, rotations) makes it unknown
\def\tikz@td@patchcalc{%
\ifdefined\tikz@cc@after@num\ifx\tikz@td@calcdone\relax\else
\global\let\tikz@td@calcdone\relax
\gdef\tikz@parse@calculator##1(${%$
  \def\tikz@cc@command{##1}%
  \begingroup
    \pgf@xa=0pt\pgf@ya=0pt\tikz@td@cza=0pt\def\tikz@td@cok{1}\def\tikz@td@cpend{0}%
    \tikz@cc@parse+}%
\global\let\tikz@td@orig@ccparse\tikz@cc@parse
\gdef\tikz@cc@parse{%
  \ifnum\tikz@td@cpend=1 \advance\tikz@td@cza by\tikz@cc@factor\tikz@td@czb
  \def\tikz@td@cpend{0}\fi
  \tikz@td@orig@ccparse}%
\gdef\tikz@cc@after@coordinate##1{%
  \tikz@td@zreset\pgf@process{##1}%
  \pgf@xb=\pgf@x\pgf@yb=\pgf@y\tikz@td@czb=\tikz@td@z
  \iftikz@td@zok\else\def\tikz@td@cok{0}\fi\def\tikz@td@cpend{1}%
  \tikz@cc@mid@checks}%
\gdef\tikz@cc@after@num##1{%
  \tikz@td@zreset\pgf@process{##1}%
  \iftikz@td@zok\else\def\tikz@td@cok{0}\fi
  \tikz@td@czb=\tikz@cc@mid@factor@one\tikz@td@czb
  \advance\tikz@td@czb by\tikz@cc@mid@factor\tikz@td@z
  \pgf@xb=\tikz@cc@mid@factor@one\pgf@xb
  \pgf@yb=\tikz@cc@mid@factor@one\pgf@yb
  \advance\pgf@xb by\tikz@cc@mid@factor\pgf@x
  \advance\pgf@yb by\tikz@cc@mid@factor\pgf@y
  \tikz@cc@mid@checks}%
\global\let\tikz@td@orig@ccunit\tikz@cc@after@unit
\gdef\tikz@cc@after@unit{\def\tikz@td@cok{0}\tikz@td@orig@ccunit}%
\global\let\tikz@td@orig@ccproject\tikz@cc@after@project
\gdef\tikz@cc@after@project{\def\tikz@td@cok{0}\tikz@td@orig@ccproject}%
\global\let\tikz@td@orig@ccrot\tikz@cc@handle@rot
\gdef\tikz@cc@handle@rot{\def\tikz@td@cok{0}\tikz@td@orig@ccrot}%
\gdef\tikz@cc@end$##1){%$
  \xdef\tikz@marshal{\noexpand\pgfqpoint{\the\pgf@xa}{\the\pgf@ya}%
    \noexpand\tikz@td@setz{\the\tikz@td@cza}{\tikz@td@cok}}%
  \endgroup
  \expandafter\tikz@cc@command\expandafter{\tikz@marshal}}%
\fi\fi}%
\tikz@td@patchcalc
\AtBeginDocument{\tikz@td@patchcalc}%
% --- recording named nodes and coordinates
\def\tikz@td@defaultat{\pgfqpoint{\the\tikz@lastx}{\the\tikz@lasty}}%
\let\tikz@td@orig@nodetrans\tikz@node@transformations
\def\tikz@node@transformations{\tikz@td@record\tikz@td@orig@nodetrans}%
\def\tikz@td@record{%
\def\tikz@td@exact{0}%
\ifx\tikz@time\pgfutil@empty
 \ifx\tikz@node@at\tikz@td@defaultat
  \edef\tikz@td@p{{\the\tikz@lastx}{\the\tikz@lasty}}%
  \tikz@td@za=\tikz@td@lastz\let\tikz@td@rok\tikz@td@lastok
  \ifnum\tikz@td@lastsimple=1 \tikz@td@checkexact\fi
 \else
  \tikz@td@zreset\pgf@process{\tikz@node@at}%
  \edef\tikz@td@p{{\the\pgf@x}{\the\pgf@y}}%
  \edef\tikz@node@at{\noexpand\pgfqpoint\tikz@td@p}%
  \tikz@td@za=\tikz@td@z\edef\tikz@td@rok{\iftikz@td@zok1\else0\fi}%
  \ifx\tikz@td@pend\relax\else\expandafter\tikz@td@usepend\tikz@td@pend\fi
 \fi
\else
 \def\tikz@td@rok{0}\def\tikz@td@p{{0pt}{0pt}}\tikz@td@za=0pt
 \ifx\tikz@timer\tikz@timer@line\ifx\tikz@timer@start\tikz@td@tstart
  \pgf@process{\pgfpointlineattime{\tikz@time}{\tikz@timer@start}{\tikz@timer@end}}%
  \edef\tikz@td@p{{\the\pgf@x}{\the\pgf@y}}%
  \tikz@td@za=\tikz@td@lastz\advance\tikz@td@za by-\tikz@td@tstartz
  \tikz@td@za=\tikz@time\tikz@td@za\advance\tikz@td@za by\tikz@td@tstartz
  \ifnum\tikz@td@lastok\tikz@td@tstartok=11 \def\tikz@td@rok{1}\fi
 \fi\fi
\fi
\iftikz@td@untracked\def\tikz@td@rok{0}\fi
\ifx\tikz@fig@name\pgfutil@empty\else
 \iftikz@td@keeptext
  \global\expandafter\let\csname tikz@td@n@\tikz@fig@name\endcsname\tikz@td@undefined
 \else
  \tikz@td@store
 \fi
\fi
% the node (or pic) has its own coordinate system with origin at the node
\ifnum\tikz@td@rok=1 \advance\tikz@td@sd by\tikz@td@za\else\tikz@td@untrackedtrue\fi}%
% coordinate (P) at (...) hands TikZ the canvas position as text; keep its depth
\global\let\tikz@td@pend\relax
\def\tikz@@coordinate@at@math#1{%
  \tikz@td@zreset\pgf@process{#1}%
  \xdef\tikz@td@pend{{\the\pgf@x}{\the\pgf@y}{\the\tikz@td@z}{\iftikz@td@zok1\else0\fi}}%
  \edef\tikz@temp{(\the\pgf@x,\the\pgf@y)}%
  \expandafter\tikz@coordinate@caller\tikz@temp{}%
}%
\def\tikz@td@usepend#1#2#3#4{%
\global\let\tikz@td@pend\relax
\edef\tikz@td@tmp{{#1}{#2}}\ifx\tikz@td@tmp\tikz@td@p\tikz@td@za=#3\def\tikz@td@rok{#4}\fi}%
\def\tikz@td@checkexact{%
\expandafter\ifx\csname tikz@dcl@coord@\tikz@fig@name\endcsname\tikz@td@lastspec
\def\tikz@td@exact{1}\fi}%
\def\tikz@td@store{%
\pgf@process{\pgfpointtransformed{\expandafter\pgfqpoint\tikz@td@p}}%
\ifnum\tikz@td@exact=1 \pgfgettransform\tikz@td@T\fi
\expandafter\xdef\csname tikz@td@n@\tikz@fig@name\endcsname{%
{\the\dimexpr\tikz@td@za+\tikz@td@sd\relax}{\the\pgf@x}{\the\pgf@y}\tikz@td@p
{\the\tikz@td@za}{\tikz@td@rok}{\tikz@td@exact}%
{{\pgfpictureid}{\tikz@td@mainvec}{\tikz@td@lf@id}{\ifnum\tikz@td@exact=1 \tikz@td@ctxkey\fi}}}}%
\tikzset{3d/keep text/.code={\tikz@td@keeptexttrue}}%
% --- looking a record up: \tikz@td@lkmode is 0 (use the text), 1 (the point is
% \tikz@td@lkv) or 2 (depth unknown). Inside the picture and under the main view
% the point was defined with, the result refers to the coordinate system in
% force (shifts and local frames included); elsewhere (other views, node texts,
% outside pictures), to the coordinate system the point was defined in.
\def\tikz@td@lookup#1{\def\tikz@td@lkname{#1}%
\expandafter\expandafter\expandafter\tikz@td@lookup@\csname tikz@td@n@#1\endcsname}%
% the stored text is exact if made of numbers and + - * / ( ) . , only
\def\tikz@td@ifpuretext#1{\expandafter\expandafter\expandafter\tikz@td@ifpure
\expandafter\expandafter\expandafter{\csname tikz@dcl@coord@\tikz@td@lkname\endcsname}{#1}{}}%
\def\tikz@td@lookup@#1#2#3#4#5#6#7#8#9{%
\ifnum#7=0 \def\tikz@td@lkmode{2}\else\tikz@td@lookup@@{#1}{#2}{#3}{#4}{#5}{#6}{#8}#9\fi}%
\def\tikz@td@lookup@@#1#2#3#4#5#6#7#8#9{% D cx cy px py z exact picture vectors
\def\tikz@td@tmp{#8}%
\def\tikz@td@next{\tikz@td@lookup@rec{#4}{#5}{#6}{#7}{#9}}%
\ifpgfpicture\ifx\tikz@td@tmp\pgfpictureid
\edef\tikz@td@tmp{\tikz@td@mainvec}\def\tikz@td@tmpb{#9}%
\ifx\tikz@td@tmp\tikz@td@tmpb
\def\tikz@td@next{\tikz@td@lookup@cur{#1}{#2}{#3}{#7}}%
\fi\fi\fi
\tikz@td@next}%
\def\tikz@td@lookup@cur#1#2#3#4#5#6{% D cx cy exact frame context
\def\tikz@td@lkmode{1}%
\ifnum#4=1
\pgfgettransform\tikz@td@T\edef\tikz@td@tmp{\tikz@td@ctxkey}\def\tikz@td@tmpb{#6}%
\ifx\tikz@td@tmp\tikz@td@tmpb\tikz@td@ifpuretext{\def\tikz@td@lkmode{0}}\fi
\fi
\ifnum\tikz@td@lkmode=1 \tikz@td@reconstruct{#2}{#3}{#1}\fi}%
\def\tikz@td@lookup@rec#1#2#3#4#5#6#7{% px py z exact vectors frame context
\def\tikz@td@lkmode{1}%
\ifnum#4=1 \tikz@td@ifpuretext{\def\tikz@td@lkmode{0}}\fi
\ifnum\tikz@td@lkmode=1
\begingroup
\tikz@td@minv{#5}%
\csname tikz@td@lf@tab@#6\endcsname
\tikz@td@solve{\fpeval{#1/1cm},\fpeval{#2/1cm},\fpeval{#3/1cm}}%
\xdef\tikz@td@lkv{\tikz@td@lkv}%
\endgroup
\fi}%
% \tikz@td@reconstruct{cx}{cy}{D}: coordinates, in the coordinate system in
% force, of the point with canvas position (cx,cy) and depth D
\def\tikz@td@reconstruct#1#2#3{%
\begingroup
\tikz@td@lf@getmain
\iftikz@td@lf@active\else\tikz@td@lf@identity\fi
\ifpgfpicture\else\tikz@td@sd=0pt\fi
% invert the transformation: c = A p + s, A = (aa ba; ab bb)
\edef\tikz@td@tmp{\fpeval{\pgf@pt@aa*\pgf@pt@bb-\pgf@pt@ba*\pgf@pt@ab}}%
\edef\tikz@td@q{%
\fpeval{(\pgf@pt@bb*(#1-\pgf@pt@x)-\pgf@pt@ba*(#2-\pgf@pt@y))/\tikz@td@tmp/1cm},%
\fpeval{(\pgf@pt@aa*(#2-\pgf@pt@y)-\pgf@pt@ab*(#1-\pgf@pt@x))/\tikz@td@tmp/1cm},%
\fpeval{(#3-\tikz@td@sd)/1cm}}%
\tikz@td@solve{\tikz@td@q}%
\xdef\tikz@td@lkv{\tikz@td@lkv}%
\endgroup}%
\def\tikz@td@unknownwarning#1{%
\ifcsname tikz@td@warned@#1\endcsname\else
\global\expandafter\let\csname tikz@td@warned@#1\endcsname\relax
\PackageWarning{3dtools}{The 3d position of `#1' is not known (it was defined
with a curve, an intersection, a perpendicular or non-affine calc coordinate,
or in a canvas is ... plane scope). Using the text it was defined with}%
\fi}%
% --- local frames
\def\tikz@td@lf@install#1{%
\begingroup
\iftikz@td@lf@active\tikz@td@lf@restoremain\tikz@td@lf@activefalse\fi
\def\tikz@td@lf@id{0}%
\pgfqkeys{/tikz/3d/local frame}{#1}%
\edef\pgfutil@tmpo{\pgfkeysvalueof{/tikz/3d/local frame/origin}}%
\edef\pgfutil@tmpa{\pgfkeysvalueof{/tikz/3d/local frame/ex}}%
\edef\pgfutil@tmpb{\pgfkeysvalueof{/tikz/3d/local frame/ey}}%
\edef\pgfutil@tmpc{\pgfkeysvalueof{/tikz/3d/local frame/ez}}%
\pgfmathsetmacro\pgfutil@tmpo{TD("\pgfutil@tmpo")}%
\pgfmathsetmacro\pgfutil@tmpa{TD("\pgfutil@tmpa")}%
\pgfmathsetmacro\pgfutil@tmpb{TD("\pgfutil@tmpb")}%
\iftikz@td@lf@orthonormalize
\pgfmathsetmacro\pgfutil@tmpa{TDunit("(\pgfutil@tmpa)")}%
\pgfmathsetmacro\pgfutil@tmpc{TDunit("(\pgfutil@tmpa)x(\pgfutil@tmpb)")}%
\pgfmathsetmacro\pgfutil@tmpb{TD("(\pgfutil@tmpc)x(\pgfutil@tmpa)")}%
\else
\ifx\pgfutil@tmpc\pgfutil@empty
\pgfmathsetmacro\pgfutil@tmpc{TD("(\pgfutil@tmpa)x(\pgfutil@tmpb)")}%
\else
\pgfmathsetmacro\pgfutil@tmpc{TD("\pgfutil@tmpc")}%
\fi
\fi
\tikz@td@dualbasis\pgfutil@tmpa\pgfutil@tmpb\pgfutil@tmpc
\ifx\tikz@td@det\pgfutil@empty
\PackageWarning{3dtools}{The vectors ex=(\pgfutil@tmpa), ey=(\pgfutil@tmpb)
and ez=(\pgfutil@tmpc) are (almost) coplanar. Local frame not installed.}%
\global\let\tikz@td@lf@new\relax
\else
\xdef\tikz@td@lf@g{\def\noexpand\tikz@td@lf@o{\pgfutil@tmpo}%
\def\noexpand\tikz@td@lf@a{\pgfutil@tmpa}\def\noexpand\tikz@td@lf@b{\pgfutil@tmpb}%
\def\noexpand\tikz@td@lf@c{\pgfutil@tmpc}\def\noexpand\tikz@td@lf@da{\tikz@td@da}%
\def\noexpand\tikz@td@lf@db{\tikz@td@db}\def\noexpand\tikz@td@lf@dc{\tikz@td@dc}}%
\global\let\tikz@td@lf@new\tikz@td@lf@g
\edef\pgfutil@tmpn{\pgfkeysvalueof{/tikz/3d/local frame/name}}%
\ifx\pgfutil@tmpn\pgfutil@empty\else
\expandafter\xdef\csname tikz@td@lf@named@\pgfutil@tmpn\endcsname{%
origin={(\pgfutil@tmpo)},ex={(\pgfutil@tmpa)},ey={(\pgfutil@tmpb)},ez={(\pgfutil@tmpc)}}%
\fi
\fi
\endgroup
\ifx\tikz@td@lf@new\relax\else
\iftikz@td@lf@active\tikz@td@lf@restoremain\else\tikz@td@lf@getmain\fi
\tikz@td@lf@new
\expandafter\tikz@td@split\tikz@td@lf@a\relax
\pgf@process{\pgfpointxyz{\tikz@td@c@a}{\tikz@td@c@b}{\tikz@td@c@c}}%
\edef\pgfutil@tmpa{{\the\pgf@x}{\the\pgf@y}}%
\expandafter\tikz@td@split\tikz@td@lf@b\relax
\pgf@process{\pgfpointxyz{\tikz@td@c@a}{\tikz@td@c@b}{\tikz@td@c@c}}%
\edef\pgfutil@tmpb{{\the\pgf@x}{\the\pgf@y}}%
\expandafter\tikz@td@split\tikz@td@lf@c\relax
\pgf@process{\pgfpointxyz{\tikz@td@c@a}{\tikz@td@c@b}{\tikz@td@c@c}}%
\edef\pgfutil@tmpc{{\the\pgf@x}{\the\pgf@y}}%
\tikz@td@zreset\pgf@process{\expandafter\tikz@td@pointxyz\tikz@td@lf@o\relax}%
\pgftransformshift{\pgfqpoint{\pgf@x}{\pgf@y}}%
\advance\tikz@td@sd by\tikz@td@z
\pgfsetxvec{\expandafter\pgfqpoint\pgfutil@tmpa}%
\pgfsetyvec{\expandafter\pgfqpoint\pgfutil@tmpb}%
\pgfsetzvec{\expandafter\pgfqpoint\pgfutil@tmpc}%
\tikz@td@dzx=\fpeval{\tikz@td@dot{\tikz@td@m@n}{\tikz@td@lf@a}}cm
\tikz@td@dzy=\fpeval{\tikz@td@dot{\tikz@td@m@n}{\tikz@td@lf@b}}cm
\tikz@td@dzz=\fpeval{\tikz@td@dot{\tikz@td@m@n}{\tikz@td@lf@c}}cm
\global\advance\tikz@td@frameid by1
\edef\tikz@td@lf@id{\the\tikz@td@frameid}%
\expandafter\xdef\csname tikz@td@lf@tab@\tikz@td@lf@id\endcsname{%
\def\noexpand\tikz@td@lf@da{\tikz@td@lf@da}\def\noexpand\tikz@td@lf@db{\tikz@td@lf@db}%
\def\noexpand\tikz@td@lf@dc{\tikz@td@lf@dc}}%
\tikz@td@lf@activetrue
\fi}%
\tikzset{3d/install local frame/.code={\tikz@addtransform{\tikz@td@lf@install{#1}}},
3d/use local frame/.code={%
\ifcsname tikz@td@lf@named@#1\endcsname
\edef\pgfutil@tmp{\noexpand\tikz@addtransform{\noexpand\tikz@td@lf@install{%
\csname tikz@td@lf@named@#1\endcsname}}}\pgfutil@tmp
\else
\PackageWarning{3dtools}{Unknown local frame `#1'.}%
\fi},
3d/main frame/.code={\tikz@addtransform{\iftikz@td@lf@active
\tikz@td@lf@restoremain\tikz@td@lf@activefalse\def\tikz@td@lf@id{0}\fi}},
3d/local frame/.cd,origin/.initial={(0,0,0)},ex/.initial={(1,0,0)},
ey/.initial={(0,1,0)},ez/.initial={},name/.initial={},
orthonormalize/.is if=tikz@td@lf@orthonormalize,orthonormalize/.default=true,
from dreibein/.style={origin=\pgfkeysvalueof{/tikz/3d/aux keys/A},
 ex=\pgfkeysvalueof{/tikz/3d/aux keys/ex},ey=\pgfkeysvalueof{/tikz/3d/aux keys/ey},
 ez=\pgfkeysvalueof{/tikz/3d/aux keys/ez}}}%
% local -> main and main -> local, for points (the origin is taken into account)
\pgfmathdeclarefunction{TDtomain}{1}{%
\begingroup
\tikz@td@lf@load
\pgfmathsetmacro\pgfutil@tmpv{TD("#1")}%
\expandafter\tikz@td@split\pgfutil@tmpv\relax
\pgfmathparse{TD("(\tikz@td@lf@o)+\tikz@td@c@a*(\tikz@td@lf@a)%
+\tikz@td@c@b*(\tikz@td@lf@b)+\tikz@td@c@c*(\tikz@td@lf@c)")}%
\pgfmathsmuggle\pgfmathresult\endgroup}%
\pgfmathdeclarefunction{TDtolocal}{1}{%
\begingroup
\tikz@td@lf@load
\pgfmathsetmacro\pgfutil@tmpv{TD("#1")}%
\pgfmathsetmacro\pgfutil@tmpa{TD("(\pgfutil@tmpv)-(\tikz@td@lf@o)o(\tikz@td@lf@da)")}%
\pgfmathsetmacro\pgfutil@tmpb{TD("(\pgfutil@tmpv)-(\tikz@td@lf@o)o(\tikz@td@lf@db)")}%
\pgfmathsetmacro\pgfutil@tmpc{TD("(\pgfutil@tmpv)-(\tikz@td@lf@o)o(\tikz@td@lf@dc)")}%
\edef\pgfmathresult{\pgfutil@tmpa,\pgfutil@tmpb,\pgfutil@tmpc}%
\pgfmathsmuggle\pgfmathresult\endgroup}%
% --- physical components: every named node and coordinate is recorded, the
% keys are kept for compatibility
\tikzset{3d/record physical components/.code={},3d/store components/.code={}}%
\def\tikz@td@phys#1#2{%
\ifcsname tikz@td@n@#2\endcsname
\expandafter\expandafter\expandafter\tikz@td@phys@\csname tikz@td@n@#2\endcsname#1%
\else
\PackageWarning{3dtools}{Physical #1 component of #2 not stored. Setting it to 0.}%
\def\pgfmathresult{0}%
\fi}%
\def\tikz@td@phys@#1#2#3#4#5#6#7#8#9{\tikz@td@phys@@{#1}{#2}{#3}{#7}}%
\def\tikz@td@phys@@#1#2#3#4#5{% D cx cy ok component
\if z#5\ifnum#4=0 \PackageWarning{3dtools}{The depth of a coordinate is not
known. Setting it to 0.}\def\pgfmathresult{0}\else
\edef\pgfmathresult{\fpeval{#1/1cm}}\fi\fi
\if x#5\edef\pgfmathresult{\fpeval{#2/1cm}}\fi
\if y#5\edef\pgfmathresult{\fpeval{#3/1cm}}\fi}%
\pgfmathdeclarefunction{TDphysx}{1}{\begingroup\tikz@td@phys{x}{#1}%
\pgfmathsmuggle\pgfmathresult\endgroup}%
\pgfmathdeclarefunction{TDphysy}{1}{\begingroup\tikz@td@phys{y}{#1}%
\pgfmathsmuggle\pgfmathresult\endgroup}%
\pgfmathdeclarefunction{TDphysz}{1}{\begingroup\tikz@td@phys{z}{#1}%
\pgfmathsmuggle\pgfmathresult\endgroup}%
\pgfmathdeclarefunction{TDphys}{1}{%
\begingroup%
\pgfmathsetmacro{\pgfutil@tmpx}{TDphysx("#1")}%
\pgfmathsetmacro{\pgfutil@tmpy}{TDphysy("#1")}%
\pgfmathsetmacro{\pgfutil@tmpz}{TDphysz("#1")}%
\edef\pgfmathresult{\pgfutil@tmpx,\pgfutil@tmpy,\pgfutil@tmpz}%
\pgfmathsmuggle\pgfmathresult\endgroup%
}%
% 3d/install dual basis: coordinates dual ex/ey/ez with TDlocalx(P)=dual ex.
% (phys(P)-phys(origin)) etc., i.e. the rows of the inverse of the matrix that
% maps coordinates of the current frame to physical components. For orthonormal
% frames and views they coincide with the images of the basis vectors.
\tikzset{3d/.cd,local ex/.initial={localex},local ey/.initial={localey},
local ez/.initial={localez},dual ex/.initial={dualex},dual ey/.initial={dualey},
dual ez/.initial={dualez},install dual basis/.code={%
\begingroup
\tikz@td@lf@getmain
\iftikz@td@lf@active\else\tikz@td@lf@identity\fi
\pgfutil@for\pgf@tmp:={a,b,c}\do{%
\edef\pgfutil@tmpe{\csname tikz@td@lf@\pgf@tmp\endcsname}%
\pgfmathsetmacro\pgfutil@tmpx{TD("(\tikz@td@m@r)o(\pgfutil@tmpe)")}%
\pgfmathsetmacro\pgfutil@tmpy{TD("(\tikz@td@m@s)o(\pgfutil@tmpe)")}%
\pgfmathsetmacro\pgfutil@tmpz{TD("(\tikz@td@m@n)o(\pgfutil@tmpe)")}%
\expandafter\xdef\csname pgfutil@tmpcol\pgf@tmp\endcsname{\pgfutil@tmpx,\pgfutil@tmpy,\pgfutil@tmpz}}%
\tikz@td@dualbasis\pgfutil@tmpcola\pgfutil@tmpcolb\pgfutil@tmpcolc
\xdef\pgfutil@tmp{\noexpand\path (\tikz@td@da) coordinate[/tikz/3d/keep text] (\pgfkeysvalueof{/tikz/3d/dual ex})
 (\tikz@td@db) coordinate[/tikz/3d/keep text] (\pgfkeysvalueof{/tikz/3d/dual ey})
 (\tikz@td@dc) coordinate[/tikz/3d/keep text] (\pgfkeysvalueof{/tikz/3d/dual ez});}%
\endgroup
\pgfutil@tmp}}
% TDlocal("P"): coordinates of P in the coordinate system in force (main frame
% or active local frame; in node texts, outside the shifts of the scope).
\pgfmathdeclarefunction{TDlocal}{1}{%
\begingroup%
\ifcsname tikz@td@n@#1\endcsname
\expandafter\expandafter\expandafter\tikz@td@local@\csname tikz@td@n@#1\endcsname
\let\pgfmathresult\tikz@td@lkv
\else
\PackageWarning{3dtools}{Physical components of #1 not stored. Setting result to 0.}%
\edef\pgfmathresult{0,0,0}%
\fi
\pgfmathsmuggle\pgfmathresult\endgroup%
}%
\def\tikz@td@local@#1#2#3#4#5#6#7#8#9{\tikz@td@reconstruct{#2}{#3}{#1}}%
\pgfmathdeclarefunction{TDlocalx}{1}{\begingroup\pgfmathparse{TDlocal("#1")}%
\expandafter\tikz@td@split\pgfmathresult\relax\let\pgfmathresult\tikz@td@c@a
\pgfmathsmuggle\pgfmathresult\endgroup}%
\pgfmathdeclarefunction{TDlocaly}{1}{\begingroup\pgfmathparse{TDlocal("#1")}%
\expandafter\tikz@td@split\pgfmathresult\relax\let\pgfmathresult\tikz@td@c@b
\pgfmathsmuggle\pgfmathresult\endgroup}%
\pgfmathdeclarefunction{TDlocalz}{1}{\begingroup\pgfmathparse{TDlocal("#1")}%
\expandafter\tikz@td@split\pgfmathresult\relax\let\pgfmathresult\tikz@td@c@c
\pgfmathsmuggle\pgfmathresult\endgroup}%
% \def\tizk@td@scalprod#1=#2.#3;{% 
% \edef\coordA{\tikz@td@raw@coord#2}% 
% \edef\coordB{\tikz@td@raw@coord#3}% 
% \pgfmathsetmacro\pgfutil@tmpa{scalarproduct({\coordA},{\coordB})}% 
% \edef#1{\pgfutil@tmpa}}% 
\def\tizk@td@spaux#1#2#3#4#5#6{(#1)*(#4)+(#2)*(#5)+(#3)*(#6)}% 
\pgfmathdeclarefunction{scalarproduct}{2}{% scalar product of two 3-vectors 
\begingroup% 
\pgfmathparse@td@FPU{\tizk@td@spaux#1#2}% 
\pgfmathsmuggle\pgfmathresult\endgroup} 
% vector product 
% vector product auxiliary functions 
\def\vpauxx#1#2#3#4#5#6{(#2)*(#6)-(#3)*(#5)}% 
\def\vpauxy#1#2#3#4#5#6{(#4)*(#3)-(#1)*(#6)}% 
\def\vpauxz#1#2#3#4#5#6{(#1)*(#5)-(#2)*(#4)}% 
% vector product pgf functions 
\pgfmathdeclarefunction{vpx}{2}{% x component of vector product 
\begingroup% 
\pgfmathparse@td@FPU{\vpauxx#1#2}% 
\pgfmathsmuggle\pgfmathresult\endgroup} 
\pgfmathdeclarefunction{vpy}{2}{% y component of vector product 
\begingroup% 
\pgfmathparse@td@FPU{\vpauxy#1#2}% 
\pgfmathsmuggle\pgfmathresult\endgroup} 
\pgfmathdeclarefunction{vpz}{2}{% z component of vector product 
\begingroup% 
\pgfmathparse@td@FPU{\vpauxz#1#2}% 
\pgfmathsmuggle\pgfmathresult\endgroup} 
%
% axisangles, usage \pgfmathsetmacro{\myaxisangles}{axisangles("(myn)")}
\pgfmathdeclarefunction{axisangles}{1}{% 
\begingroup% 
\pgfmathsetmacro{\pgfutil@tmpa}{sqrt(TD("#1o#1"))}%
\edef\pgfutil@tmpb{\tikz@td@raw@coord#1}%
\pgfmathsetmacro{\pgftemp@x}{xcomp3(\pgfutil@tmpb)/\pgfutil@tmpa}% 
\pgfmathsetmacro{\pgftemp@y}{ycomp3(\pgfutil@tmpb)/\pgfutil@tmpa}% 
\pgfmathsetmacro{\pgftemp@z}{zcomp3(\pgfutil@tmpb)/\pgfutil@tmpa}% 
\pgfmathtruncatemacro{\itest}{ifthenelse(abs(\pgftemp@z)==1,0,1)}%
\ifnum\itest=0%
	\pgfmathtruncatemacro{\jtest}{sign(\pgftemp@z)}%
	\ifnum\jtest=1%
		\edef\pgfutil@tmpc{0,0}%
	\else
		\edef\pgfutil@tmpc{0,180}%
	\fi	
\else
	\pgfmathsetmacro{\pgfutil@tmpb}{acos(\pgftemp@z)}% polar angle in [0,180]
	\pgfmathsetmacro{\pgfutil@tmpa}{atan2(\pgftemp@y,\pgftemp@x)}% azimuth
	\edef\pgfutil@tmpc{\pgfutil@tmpa,\pgfutil@tmpb}%
\fi
\edef\pgfmathresult{\pgfutil@tmpc}%
\pgfmathsmuggle\pgfmathresult\endgroup} 
% returns a vector that is orthogonal to the input vector
\pgfmathdeclarefunction{normalcandidate}{3}{%
\begingroup%
\pgfmathtruncatemacro{\icase}{(abs(#1)>0)+2*(abs(#2)>0)+4*(abs(#3)>0)}%
\ifcase\icase
\message{All three coordinates are zero. No normal can be computed.^^J}%
\edef\pgfmathresult{0,0,0}%
\or
\edef\pgfmathresult{0,1,0}%
\or
\edef\pgfmathresult{0,0,1}%
\or
\edef\pgfmathresult{0,0,1}%
\or
\edef\pgfmathresult{1,0,0}%
\or
\edef\pgfmathresult{0,1,0}%
\or
\edef\pgfmathresult{1,0,0}%
\or
\pgfmathsetmacro{\normalization}{(#1)*(#1)+(#2)*(#2)}%
\pgfmathsetmacro{\nx}{-1*(#2)/sqrt(\normalization)}%
\pgfmathsetmacro{\ny}{#1/sqrt(\normalization)}%
\edef\pgfmathresult{\nx,\ny,0}%
\fi
\pgfmathsmuggle\pgfmathresult\endgroup%
}%
\pgfmathdeclarefunction{randompermutation}{1}{%
\begingroup%	      
\c@pgf@counta=0%
\edef\pgfutil@tmpa{0}%
\loop
\advance\c@pgf@counta by1%
\edef\pgfutil@tmpa{\pgfutil@tmpa,\the\c@pgf@counta}%
\ifnum\c@pgf@counta<#1\relax
\repeat
\loop
\advance\c@pgf@counta by-1%
\pgfmathtruncatemacro{\pgfutil@tmpb}{min(1+rnd*(\c@pgf@counta+1),\c@pgf@counta+1)}%
\pgfmathtruncatemacro{\pgfutil@tmpc}{{\pgfutil@tmpa}[\pgfutil@tmpb]}%
\ifnum\c@pgf@counta=\numexpr#1-1\relax
\edef\pgfmathresult{\pgfutil@tmpc}%
\else
\edef\pgfmathresult{\pgfmathresult,\pgfutil@tmpc}%
\fi
\edef\pgfutil@tmpd{\pgfutil@tmpa}%
\edef\pgfutil@tmpa{0}%
\pgfutil@for\my@item:={\pgfutil@tmpd}\do{%
\unless\ifnum\pgfutil@tmpc=\my@item
\unless\ifnum\my@item=0\relax
\edef\pgfutil@tmpa{\pgfutil@tmpa,\my@item}%
\fi\fi}%
\ifnum\c@pgf@counta>0\relax
\repeat
\pgfmath@smuggleone\pgfmathresult%
\endgroup}
\pgfmathdeclarefunction{take}{2}{%
\begingroup%
\c@pgf@counta=0%
\pgfutil@for\my@item:={#1}\do{%
\advance\c@pgf@counta by1%
\ifnum\c@pgf@counta<#2\relax
\ifnum\c@pgf@counta=1\relax
\edef\pgfmathresult{\my@item}%
\else
\edef\pgfmathresult{\pgfmathresult,\my@item}%
\fi\fi}%
\pgfmath@smuggleone\pgfmathresult%
\endgroup}
%
%
% barycenter of a list of coordinates
\pgfmathdeclarefunction{barycenter}{1}{% barycentric coordinate 
\begingroup%
\edef\pgfutil@tmpb{0}%
\tikzset{3d/tmp/.code={\edef\pgfutil@tmpb{\the\numexpr\pgfutil@tmpb+1}%
\ifnum\pgfutil@tmpb=1
 \edef\pgfutil@tmpc{##1}%
\else
 \edef\pgfutil@tmpc{\pgfutil@tmpc+##1}%
\fi},3d/tmp/.list={#1}}%
\pgfmathsetmacro{\pgfutil@tmpc}{TD("\pgfutil@tmpc")}%
\pgfmathsetmacro{\pgfutil@tmpa}{1/\pgfutil@tmpb}%
\pgfmathparse{TD("\pgfutil@tmpa*(\pgfutil@tmpc)")}%   
\pgfmathsmuggle\pgfmathresult\endgroup}%
%
% position of a real in a list of reals
\pgfmathdeclarefunction*{pos}{2}{%
\begingroup%
\c@pgf@counta=1%
\pgfutil@for\my@item:={#2}\do{%
\ifdim\my@item pt<#1pt%
\advance\c@pgf@counta by1%
\fi}%
\edef\pgfmathresult{\the\c@pgf@counta}%
\pgfmathsmuggle\pgfmathresult\endgroup%
}
\pgfmathdeclarefunction*{x2d}{1}{%
\begingroup%
\pgf@process{\pgfpointanchor{#1}{center}}%
\pgfmathparse{\pgf@x}%
\pgfmathsmuggle\pgfmathresult\endgroup%
}%
\pgfmathdeclarefunction*{y2d}{1}{%
\begingroup%
\pgf@process{\pgfpointanchor{#1}{center}}%
\pgfmathparse{\pgf@y}%
\pgfmathsmuggle\pgfmathresult\endgroup%
}%
%
\tikzset{convex hull of/.code={% expects a list of points
\edef\pgfutil@tmpl{#1}% store the list for manipulations
\edef\pgfutil@tmpp{}% path to be constructed
\c@pgf@counta1%
\pgfutil@for\pgfutil@tmpa:={\pgfutil@tmpl}\do{%
\ifnum\c@pgf@counta=1\relax
\pgfmathsetmacro{\pgfutil@tmpb}{x2d("\pgfutil@tmpa")}% min
\edef\pgfutil@tmpd{\pgfutil@tmpa}% leftmost point
\else
\pgfmathsetmacro{\pgfutil@tmpc}{x2d("\pgfutil@tmpa")}% x coordinate
\ifdim\pgfutil@tmpc pt<\pgfutil@tmpb pt\relax
\edef\pgfutil@tmpb{\pgfutil@tmpc}% new min
\edef\pgfutil@tmpd{\pgfutil@tmpa}% new leftmost point
\fi
\fi
\advance\c@pgf@counta by1%
}
\edef\pgfutil@tmpp{(\pgfutil@tmpd)}%
\edef\pgfutil@tmpj{\pgfutil@tmpd}%
\edef\pgfutil@tmpe{90}% slope correction
\edef\pgfutil@tmpk{1}% first round flag
\loop
%\typeout{Removing  \pgfutil@tmpd}%
\edef\pgfutil@tmph{\noexpand\@removeelement{\pgfutil@tmpd}{\pgfutil@tmpl}{\noexpand\pgfutil@tmpl}}%
\pgfutil@tmph
%\typeout{List is now \pgfutil@tmpl}%
\c@pgf@countb1%
\pgfutil@for\pgfutil@tmpa:={\pgfutil@tmpl}\do{%
\pgfmathsetmacro{\pgfutil@tmpf}{Mod(atan2(y2d("\pgfutil@tmpa")-y2d("\pgfutil@tmpd"),
x2d("\pgfutil@tmpa")-x2d("\pgfutil@tmpd"))+\pgfutil@tmpe,360)}%
%\typeout{\the\c@pgf@countb: \pgfutil@tmpa\space has angle \pgfutil@tmpf.}%
\ifnum\c@pgf@countb=1\relax
\edef\pgfutil@tmpg{\pgfutil@tmpf}%
\edef\pgfutil@tmpi{\pgfutil@tmpa}%
\else
%\unless\ifdim\pgfutil@tmpf pt>180pt\relax
\ifdim\pgfutil@tmpf pt>\pgfutil@tmpg pt\relax
\edef\pgfutil@tmpg{\pgfutil@tmpf}%
\edef\pgfutil@tmpi{\pgfutil@tmpa}%
\fi
%\fi
\fi
\advance\c@pgf@countb by1%
}
%
\ifx\pgfutil@tmpi\pgfutil@tmpj\relax
\c@pgf@counta=2\relax
\else
\edef\pgfutil@tmpd{\pgfutil@tmpi}%
\edef\pgfutil@tmpp{\pgfutil@tmpp--(\pgfutil@tmpd)}%
\pgfmathsetmacro{\pgfutil@tmpe}{180+\pgfutil@tmpe-\pgfutil@tmpg}% angle of the selected point
%\typeout{offset angle=\pgfutil@tmpe}%
\advance\c@pgf@counta by-1%
\fi
\ifnum\pgfutil@tmpk=1\relax
\edef\pgfutil@tmpl{\pgfutil@tmpl,\pgfutil@tmpj}% add the first point again
\edef\pgfutil@tmpk{0}%
\fi
\ifnum\c@pgf@counta>2\relax
\repeat
\edef\pgfutil@tmpp{\pgfutil@tmpp--cycle}%
%\typeout{hull path=\pgfutil@tmpp}%
\tikzset{insert path=\pgfutil@tmpp}%
}}
%
\tikzset{generous outside path/.style={overlay,insert path={[reset cm]
			(-16383.99999pt,-16383.99999pt) |-
			(16383.99999pt,16383.99999pt)|- cycle}},
		 generous outside path'/.style={overlay,insert path={[reset cm]
			(-16383.99999pt,-16383.99999pt) -|
			(16383.99999pt,16383.99999pt)-| cycle}}}
%			
\tikzset{even odd clip/.code={\pgfseteorule}}			
%
\tikzset{%
	save named path/.code={\tikz@addmode{\pgfsyssoftpath@getcurrentpath\pgfutil@tmpf
		\global\expandafter\let\csname tikz@td@named@path@#1\endcsname=\pgfutil@tmpf}},
	use named path/.code={%
		\expandafter\pgfsyssoftpath@setcurrentpath\expandafter{\csname tikz@td@named@path@#1\endcsname}}
}
%
\pgfkeys{% https://tex.stackexchange.com/a/20426
  /tikz/on layer/.code={
    \pgfonlayer{#1}\begingroup
    \aftergroup\endpgfonlayer
    \aftergroup\endgroup
  },
  /tikz/node on layer/.code={
    \gdef\node@@on@layer{%
      \setbox\tikz@tempbox=\hbox\bgroup\pgfonlayer{#1}\unhbox\tikz@tempbox\endpgfonlayer\egroup}
    \aftergroup\node@on@layer
  },
  /tikz/end node on layer/.code={
    \endpgfonlayer\endgroup\endgroup
  }
}
%
\def\node@on@layer{\aftergroup\node@@on@layer}
%
\tikzset{declare function={torusx(\u,\v,\R,\r)=cos(\u)*(\R + \r*cos(\v)); 
torusy(\u,\v,\R,\r)=(\R + \r*cos(\v))*sin(\u);
torusz(\u,\v,\R,\r)=\r*sin(\v);
torusvcrit1(\u,\th)=atan(tan(\th)*sin(\u));% first critical v value
torusvcrit2(\u,\th)=180+atan(tan(\th)*sin(\u));% second critical v value
torusvtest(\u,\v,\az,\el)=sin(-torusvcrit1(\u-\az,\el)+\v);
torusdisc(\th,\R,\r)=((pow(\r,2)-pow(\R,2))*pow(cot(\th),2)+% 
pow(\r,2)*(2+pow(tan(\th),2)))/pow(\R,2);% discriminant
torusumax(\th,\R,\r)=ifthenelse(torusdisc(\th,\R,\r)>0,asin(sqrt(abs(torusdisc(\th,\R,\r)))),0);
}}%
%
\tikzset{3d parse/.style={%/utils/exec=\pgfmathtdparse{#1},%
insert path={\pgfextra{\pgfmathtdparse{#1}}(\pgfmathresult)}},
3d coordinate/.style args={#1=#2}{%
%/utils/exec=\pgfmathtdparse{#2},%
insert path={\pgfextra{\pgfmathtdparse{#2}}(\pgfmathresult) coordinate[/tikz/3d/keep text] #1}}}%
\newcommand\pgfmathprintvector[2][]{%
\begingroup
\pgfkeys{#1}%
\pgfutil@tempcnta0%
\pgfutil@for\pgf@tmp:={#2}\do{\advance\pgfutil@tempcnta by1}%
\pgfutil@tempcntb1%
\pgfutil@for\pgf@tmp:={#2}\do{\advance\pgfutil@tempcntb by1%
\ifnum\pgfutil@tempcntb<\pgfutil@tempcnta
\pgfmathparse{\pgf@tmp}%
\pgfmathprintnumber{\pgfmathresult},%
\else
\pgfmathparse{\pgf@tmp}%
\pgfmathprintnumber{\pgfmathresult}%
\fi}%
\endgroup
}%
% 
% projection of a point on a line
\tikzset{3d/projection precomputations/.code n args={3}{%
\pgfmathsetmacro{\pgfutil@tmpa}{TD("#2-#3o#1-#3")/TD("#2-#3o#2-#3")}%
\pgfmathsetmacro{\pgfutil@tmpb}{TD("#3+\pgfutil@tmpa*#2-\pgfutil@tmpa*#3")}%
},3d/projection of/.style args={#1 on #2--#3}{%
3d/projection precomputations={#1}{#2}{#3},%
insert path={[3d parse={(\pgfutil@tmpb)}]}
}}
%
% define extended objects and project point on line or plane
\tikzset{3d/.cd,
plane through/.code args={#1 and #2 and #3 named #4}{%
%\typeout{#2-#1x#3-#1}%
\pgfmathsetmacro{\pgfutil@tmpa}{TD("#2-#1x#3-#1")}%
\pgfmathsetmacro{\pgfutil@tmpb}{TD("(\pgfutil@tmpa)o(\pgfutil@tmpa)")}%
\ifnum\fpeval{\pgfutil@tmpb<0.01}=1\relax
\PackageWarning{3dtools}{The points #1 and #2 and #3 are either too close to each other or almost
on one line, so not suited to define a plane.}%
\else
\pgfmathsetmacro{\pgfutil@tmpd}{TDunit("(\pgfutil@tmpa)")}%
\pgfmathsetmacro{\pgfutil@tmpe}{TD("#1")}%
%\typeout{normal=(\pgfutil@tmpd),point=(\pgfutil@tmpe),name=#4}%
\expandafter\xdef\csname tikz@td@plane@normal@#4\endcsname{\pgfutil@tmpd}%
\expandafter\xdef\csname tikz@td@plane@point@#4\endcsname{\pgfutil@tmpe}%
\expandafter\xdef\csname tikz@td@dim@#4\endcsname{2}%
\fi
},%
plane with normal/.code args={#1 through #2 named #3}{%
\pgfmathsetmacro{\pgfutil@tmpn}{TD("#1o#1")}%
\ifnum\fpeval{\pgfutil@tmpn<0.0004}=1\relax
\PackageWarning{3dtools}{Normal #1 is too short.}%
\else
\pgfmathsetmacro{\pgfutil@tmpd}{TDunit("#1")}%
\pgfmathsetmacro{\pgfutil@tmpe}{TD("#2")}%
\expandafter\xdef\csname tikz@td@plane@normal@#3\endcsname{\pgfutil@tmpd}%
\expandafter\xdef\csname tikz@td@plane@point@#3\endcsname{\pgfutil@tmpe}%
\expandafter\xdef\csname tikz@td@dim@#3\endcsname{2}%
\fi
},%
line through/.code args={#1 and #2 named #3}{%
\pgfmathsetmacro{\pgfutil@tmpd}{TD("#1")}%
\pgfmathsetmacro{\pgfutil@tmpe}{TD("#2")}%
\pgfmathsetmacro{\pgfutil@tmpf}{TD("#1-#2o#1-#2")}%
\ifnum\fpeval{\pgfutil@tmpf<0.001}=1\relax
\PackageWarning{3dtools}{The points #1 and #2 are too close to each other. No line defined.}%
\else
\expandafter\xdef\csname tikz@td@line@A@#3\endcsname{\pgfutil@tmpd}%
\expandafter\xdef\csname tikz@td@line@B@#3\endcsname{\pgfutil@tmpe}%
\expandafter\xdef\csname tikz@td@dim@#3\endcsname{1}%
\fi
},%
line with direction/.code args={#1 through #2 named #3}{%
\pgfmathsetmacro{\pgfutil@tmpd}{TD("#2")}%
\pgfmathsetmacro{\pgfutil@tmpe}{TD("#1+#2")}%
\pgfmathsetmacro{\pgfutil@tmpf}{TD("#1o#1")}%
\ifnum\fpeval{\pgfutil@tmpf<0.001}=1\relax
\PackageWarning{3dtools}{The points #1 and #2 are too close to each other. No line defined.}%
\else
\expandafter\xdef\csname tikz@td@line@A@#3\endcsname{\pgfutil@tmpd}%
\expandafter\xdef\csname tikz@td@line@B@#3\endcsname{\pgfutil@tmpe}%
\expandafter\xdef\csname tikz@td@dim@#3\endcsname{1}%
\fi
},%
sphere with center/.code args={#1 and radius #2 named #3}{%
\pgfmathsetmacro{\pgfutil@tmpd}{TD("#1")}%
\expandafter\xdef\csname tikz@td@sphere@C@#3\endcsname{\pgfutil@tmpd}%
\expandafter\xdef\csname tikz@td@sphere@r@#3\endcsname{#2}%
\expandafter\xdef\csname tikz@td@dim@#3\endcsname{3}% reserve 3 for sphere...
},
project/.code args={#1 on #2}{%
\expandafter\edef\expandafter\pgfutil@tmpd{\csname tikz@td@dim@#2\endcsname}%
\ifcase\pgfutil@tmpd
\PackageWarning{3dtools}{It is not particularly meaningful to project a point on a point.}
\or% line
\expandafter\edef\expandafter\pgfutil@tmpa{\csname tikz@td@line@A@#2\endcsname}%
\expandafter\edef\expandafter\pgfutil@tmpb{\csname tikz@td@line@B@#2\endcsname}%
\pgfmathsetmacro{\pgfutil@tmpd}{TD("(\pgfutil@tmpa)-(\pgfutil@tmpb)o#1-(\pgfutil@tmpb)")/TD("(\pgfutil@tmpa)-(\pgfutil@tmpb)o(\pgfutil@tmpa)-(\pgfutil@tmpb)")}%
\pgfmathsetmacro{\pgfutil@tmpc}{TD("(\pgfutil@tmpb)+\pgfutil@tmpd*(\pgfutil@tmpa)-\pgfutil@tmpd*(\pgfutil@tmpb)")}%
%\typeout{alpha=\pgfutil@tmpd,projection=(\pgfutil@tmpc)}%
\tikzset{insert path={(\pgfutil@tmpc)}}%
\or% plane
\expandafter\edef\expandafter\pgfutil@tmpn{\csname tikz@td@plane@normal@#2\endcsname}%
\expandafter\edef\expandafter\pgfutil@tmpa{\csname tikz@td@plane@point@#2\endcsname}%
\pgfmathsetmacro{\pgfutil@tmpb}{TD("(\pgfutil@tmpn)o(\pgfutil@tmpa)-#1")}%
\pgfmathsetmacro{\pgfutil@tmpc}{TD("#1+\pgfutil@tmpb*(\pgfutil@tmpn)")}%
%\typeout{alpha=\pgfutil@tmpb,projection=(\pgfutil@tmpc)}%
\tikzset{insert path={(\pgfutil@tmpc)}}%
\fi
}}
%%
%% intersections in 3d (so far: intersection of line with line, plane, or sphere)
\tikzset{3d/intersection of/.code args={#1 with #2}{%
\expandafter\edef\expandafter\pgfutil@tmpd{\csname tikz@td@dim@#1\endcsname}%
\expandafter\edef\expandafter\pgfutil@tmpe{\csname tikz@td@dim@#2\endcsname}%
\edef\pgfutil@tmpde{\pgfutil@tmpd\pgfutil@tmpe}%
\edef\pgfutil@tmpf{0}% flag
\ifnum\pgfutil@tmpde=11\relax% line with line
 \expandafter\edef\expandafter\pgfutil@tmpa{\csname tikz@td@line@B@#1\endcsname}%
 \expandafter\edef\expandafter\pgfutil@tmpb{\csname tikz@td@line@B@#2\endcsname}%
 \expandafter\edef\expandafter\pgfutil@tmpc{\csname tikz@td@line@A@#1\endcsname}%
 \expandafter\edef\expandafter\pgfutil@tmpd{\csname tikz@td@line@A@#2\endcsname}%
\edef\pgfutil@tmpf{3}% 
\else
\ifnum\pgfutil@tmpde=21\relax% plane with line
 \expandafter\edef\expandafter\pgfutil@tmpa{\csname tikz@td@line@A@#2\endcsname}%
 \expandafter\edef\expandafter\pgfutil@tmpb{\csname tikz@td@line@B@#2\endcsname}%
 \expandafter\edef\expandafter\pgfutil@tmpn{\csname tikz@td@plane@normal@#1\endcsname}%
 \expandafter\edef\expandafter\pgfutil@tmpc{\csname tikz@td@plane@point@#1\endcsname}%
 \edef\pgfutil@tmpf{1}%
 \else
  \ifnum\pgfutil@tmpde=12\relax% line with plane
   \expandafter\edef\expandafter\pgfutil@tmpa{\csname tikz@td@line@A@#1\endcsname}%
   \expandafter\edef\expandafter\pgfutil@tmpb{\csname tikz@td@line@B@#1\endcsname}%
   \expandafter\edef\expandafter\pgfutil@tmpn{\csname tikz@td@plane@normal@#2\endcsname}%
   \expandafter\edef\expandafter\pgfutil@tmpc{\csname tikz@td@plane@point@#2\endcsname}%
   \edef\pgfutil@tmpf{1}% flag
  \else
   \ifnum\pgfutil@tmpde=13\relax% line with sphere
	\expandafter\edef\expandafter\pgfutil@tmpa{\csname tikz@td@line@A@#1\endcsname}%
	\expandafter\edef\expandafter\pgfutil@tmpb{\csname tikz@td@line@B@#1\endcsname}%
	\expandafter\edef\expandafter\pgfutil@tmpc{\csname tikz@td@sphere@C@#2\endcsname}%
	\expandafter\edef\expandafter\pgfutil@tmpr{\csname tikz@td@sphere@r@#2\endcsname}%
	\edef\pgfutil@tmpf{2}%
	%\typeout{line with sphere}%
   \else
	\ifnum\pgfutil@tmpde=31\relax% sphere with line
	\expandafter\edef\expandafter\pgfutil@tmpa{\csname tikz@td@line@A@#2\endcsname}%
	\expandafter\edef\expandafter\pgfutil@tmpb{\csname tikz@td@line@B@#2\endcsname}%
	\expandafter\edef\expandafter\pgfutil@tmpc{\csname tikz@td@sphere@C@#1\endcsname}%
	\expandafter\edef\expandafter\pgfutil@tmpr{\csname tikz@td@sphere@r@#1\endcsname}%
	\edef\pgfutil@tmpf{2}%
    %\typeout{sphere with line}%
	\else
	 \PackageWarning{3dtools}{Currently, one intersection object has to be a plane and the other 
  	   one a line. However, #1 has dimension \pgfutil@tmpd\space and 
	   #2 has dimension \pgfutil@tmpe.}%
	\fi
   \fi	
  \fi
 \fi
\fi
\ifnum\pgfutil@tmpf=1\relax% reasonable input, stored in a, b, c and n
%\typeout{debug: A=(\pgfutil@tmpa), B=(\pgfutil@tmpb), C=(\pgfutil@tmpc), n=(\pgfutil@tmpn)}%
\pgfmathsetmacro{\pgfutil@tmpg}{TD("(\pgfutil@tmpn)o(\pgfutil@tmpb)-(\pgfutil@tmpa)")}%
\pgfmathtruncatemacro{\pgfutil@tmph}{abs(\pgfutil@tmpg)>0.01?1:0}% 
\ifnum\pgfutil@tmph=1\relax%
\pgfmathsetmacro{\pgfutil@tmpi}{TD("(\pgfutil@tmpn)o(\pgfutil@tmpc)-(\pgfutil@tmpa)")/\pgfutil@tmpg}%
\pgfmathsetmacro{\pgfutil@tmpj}{TD("(\pgfutil@tmpa)+\pgfutil@tmpi*(\pgfutil@tmpb)-\pgfutil@tmpi*(\pgfutil@tmpa)")}%
\tikzset{insert path={(\pgfutil@tmpj)}}%
%\typeout{debug: intersection is at (\pgfutil@tmpj)}%
\else
\PackageWarning{3dtools}{The plane and line do not intersect within the uncertainties. Returning (0,0,0) instead.}%
\tikzset{insert path={(0,0,0)}}%
\fi
\fi
\ifnum\pgfutil@tmpf=2\relax% line and sphere
\pgfmathsetmacro@td@FPU{\pgfutil@tmpg}{TD("(\pgfutil@tmpb)-(\pgfutil@tmpa)o(\pgfutil@tmpb)-(\pgfutil@tmpa)")}%
\pgfmathsetmacro@td@FPU{\pgfutil@tmph}{2*TD("(\pgfutil@tmpa)-(\pgfutil@tmpc)o(\pgfutil@tmpb)-(\pgfutil@tmpa)")}%
\pgfmathsetmacro@td@FPU{\pgfutil@tmpi}{TD("(\pgfutil@tmpa)-(\pgfutil@tmpc)o(\pgfutil@tmpa)-(\pgfutil@tmpc)")-\pgfutil@tmpr*\pgfutil@tmpr}%
\pgfmathsetmacro@td@FPU{\pgfutil@tmpj}{sign(\pgfutil@tmph*\pgfutil@tmph-4*\pgfutil@tmpg*\pgfutil@tmpi)}%
\ifnum\fpeval{\pgfutil@tmpj<0}=1\relax
\PackageWarning{3dtools}{The sphere and line do not intersect within the uncertainties. Returning (0,0,0) instead.}%
\tikzset{insert path={(0,0,0) coordinate (\pgfkeysvalueof{/tikz/3d/aux keys/intersection 1})
  (0,0,0) coordinate (\pgfkeysvalueof{/tikz/3d/aux keys/intersection 2})}}%
\else
\pgfmathsetmacro@td@FPU{\pgfutil@tmpk}{-1*(\pgfutil@tmph+sqrt(\pgfutil@tmph*\pgfutil@tmph-4*\pgfutil@tmpg*\pgfutil@tmpi))/2/\pgfutil@tmpg}%
%\typeout{k=\pgfutil@tmpk}
\pgfmathsetmacro{\pgfutil@tmpn}{TD("(\pgfutil@tmpa)+\pgfutil@tmpk*(\pgfutil@tmpb)-\pgfutil@tmpk*(\pgfutil@tmpa)")}%
\tikzset{insert path={(\pgfutil@tmpn) coordinate (\pgfkeysvalueof{/tikz/3d/aux keys/intersection 1})}}%
\pgfmathsetmacro@td@FPU{\pgfutil@tmpk}{-1*(\pgfutil@tmph-sqrt(\pgfutil@tmph*\pgfutil@tmph-4*\pgfutil@tmpg*\pgfutil@tmpi))/2/\pgfutil@tmpg}%
%\typeout{k=\pgfutil@tmpk}
\pgfmathsetmacro{\pgfutil@tmpn}{TD("(\pgfutil@tmpa)+\pgfutil@tmpk*(\pgfutil@tmpb)-\pgfutil@tmpk*(\pgfutil@tmpa)")}%
\tikzset{insert path={(\pgfutil@tmpn) coordinate (\pgfkeysvalueof{/tikz/3d/aux keys/intersection 2})}}%
% cont
\fi
\fi
\ifnum\pgfutil@tmpf=3\relax% intersection between two lines, follow https://math.stackexchange.com/a/271366
 \pgfmathsetmacro{\pgfutil@tmpC}{TD("(\pgfutil@tmpc)")}%
 \pgfmathsetmacro{\pgfutil@tmpD}{TD("(\pgfutil@tmpd)")}%
 \pgfmathsetmacro{\pgfutil@tmpi}{sqrt(TD("(\pgfutil@tmpC)-(\pgfutil@tmpD)o(\pgfutil@tmpC)-(\pgfutil@tmpD)"))}%
 \ifnum\fpeval{\pgfutil@tmpi<0.05}=1\relax
  \tikzset{insert path={(\pgfutil@tmpC) coordinate (\pgfkeysvalueof{/tikz/3d/aux keys/intersection 1})}}%
 \else
  \pgfmathsetmacro{\pgfutil@tmpE}{TD("(\pgfutil@tmpc)-(\pgfutil@tmpa)")}%
  \pgfmathsetmacro{\pgfutil@tmpF}{TD("(\pgfutil@tmpd)-(\pgfutil@tmpb)")}%
  \pgfmathsetmacro{\pgfutil@tmpG}{TD("(\pgfutil@tmpD)-(\pgfutil@tmpC)")}%
  \pgfmathsetmacro{\pgfutil@tmpH}{TD("(\pgfutil@tmpF)x(\pgfutil@tmpG)")}%
  \pgfmathsetmacro{\pgfutil@tmpK}{TD("(\pgfutil@tmpF)x(\pgfutil@tmpE)")}%
  \pgfmathsetmacro@td@FPU{\pgfutil@tmpk}{sqrt(TD("(\pgfutil@tmpK)o(\pgfutil@tmpK)"))}%
  \ifnum\fpeval{\pgfutil@tmpk<0.01}=1\relax
   \PackageWarning{3dtools}{The lines are (almost) parallel. No intersection inserted.}%
  \else
   % distance of the lines = |G.K|/|K|
   \pgfmathsetmacro@td@FPU{\pgfutil@tmph}{abs(TD("(\pgfutil@tmpG)o(\pgfutil@tmpK)"))/\pgfutil@tmpk}%
   \ifnum\fpeval{\pgfutil@tmph>0.01}=1\relax
    \PackageWarning{3dtools}{The lines are skew (distance \pgfutil@tmph). Inserting the point of #1 closest to #2.}%
   \fi
   \pgfmathsetmacro@td@FPU{\pgfutil@tmpl}{TD("(\pgfutil@tmpH)o(\pgfutil@tmpK)")/(\pgfutil@tmpk*\pgfutil@tmpk)}%
   \pgfmathsetmacro{\pgfutil@tmpM}{TD("(\pgfutil@tmpC)+\pgfutil@tmpl*(\pgfutil@tmpE)")}%
   \tikzset{insert path={(\pgfutil@tmpM) coordinate (\pgfkeysvalueof{/tikz/3d/aux keys/intersection 1})}}%
  \fi
 \fi
\fi
}}
%%
%% circumsphere center
\tikzset{3d/circumsphere center/.code={%
\tikzset{3d/aux keys/.cd,#1}%
\edef\pgfutil@tmpA{\pgfkeysvalueof{/tikz/3d/aux keys/A}}%
\edef\pgfutil@tmpB{\pgfkeysvalueof{/tikz/3d/aux keys/B}}%
\edef\pgfutil@tmpC{\pgfkeysvalueof{/tikz/3d/aux keys/C}}%
\edef\pgfutil@tmpD{\pgfkeysvalueof{/tikz/3d/aux keys/D}}%
% shifting to A=0 by redefining B, C and D
\pgfmathsetmacro{\pgfutil@tmpB}{TD("\pgfutil@tmpB-\pgfutil@tmpA")}%
\pgfmathsetmacro{\pgfutil@tmpC}{TD("\pgfutil@tmpC-\pgfutil@tmpA")}%
\pgfmathsetmacro{\pgfutil@tmpD}{TD("\pgfutil@tmpD-\pgfutil@tmpA")}%
% 
\pgfmathsetmacro{\pgfutil@tmpBxC}{TD("(\pgfutil@tmpB)x(\pgfutil@tmpC)")}%
\pgfmathsetmacro{\pgfutil@tmpCxD}{TD("(\pgfutil@tmpC)x(\pgfutil@tmpD)")}%
\pgfmathsetmacro{\pgfutil@tmpDxB}{TD("(\pgfutil@tmpD)x(\pgfutil@tmpB)")}%
% determinant (times 2)
\pgfmathsetmacro{\pgfutil@tmpd}{2*TD("(\pgfutil@tmpB)o(\pgfutil@tmpCxD)")}%
%\typeout{Determinant is \pgfutil@tmpd.}
\pgfmathtruncatemacro{\pgfutil@tmpi}{abs(\pgfutil@tmpd)>0.02?1:0}%
\ifnum\pgfutil@tmpi=0\relax
\PackageWarning{3dtools}{The points approximately sit on a common plane. The circumsphere cannot
be determined.}%
\else
\pgfmathsetmacro{\pgfutil@tmpnB}{TD("(\pgfutil@tmpB)o(\pgfutil@tmpB)")/\pgfutil@tmpd}%
\pgfmathsetmacro{\pgfutil@tmpnC}{TD("(\pgfutil@tmpC)o(\pgfutil@tmpC)")/\pgfutil@tmpd}%
\pgfmathsetmacro{\pgfutil@tmpnD}{TD("(\pgfutil@tmpD)o(\pgfutil@tmpD)")/\pgfutil@tmpd}%
\pgfmathsetmacro{\pgfutil@tmpI}{TD("\pgfutil@tmpA+\pgfutil@tmpnB*(\pgfutil@tmpCxD)+\pgfutil@tmpnC*(\pgfutil@tmpDxB)+\pgfutil@tmpnD*(\pgfutil@tmpBxC)")}%
\tikzset{insert path={(\pgfutil@tmpI)}}%
\fi
}}%
\tikzset{3d/circumcircle center/.code={%
\tikzset{3d/aux keys/.cd,#1}%
\edef\pgfutil@tmpA{\pgfkeysvalueof{/tikz/3d/aux keys/A}}%
\edef\pgfutil@tmpB{\pgfkeysvalueof{/tikz/3d/aux keys/B}}%
\edef\pgfutil@tmpC{\pgfkeysvalueof{/tikz/3d/aux keys/C}}%
% shifting to A=0 by redefining B and C
\pgfmathsetmacro{\pgfutil@tmpB}{TD("\pgfutil@tmpB-\pgfutil@tmpA")}%
\pgfmathsetmacro{\pgfutil@tmpC}{TD("\pgfutil@tmpC-\pgfutil@tmpA")}%
%\typeout{\pgfutil@tmpB,\pgfutil@tmpC}%
% 
\pgfmathsetmacro{\pgfutil@tmpBB}{TD("(\pgfutil@tmpB)o(\pgfutil@tmpB)")}%
\pgfmathsetmacro{\pgfutil@tmpBC}{TD("(\pgfutil@tmpB)o(\pgfutil@tmpC)")}%
\pgfmathsetmacro{\pgfutil@tmpCC}{TD("(\pgfutil@tmpC)o(\pgfutil@tmpC)")}%
%\typeout{\pgfutil@tmpBB,\pgfutil@tmpBC,\pgfutil@tmpCC}%
% determinant (times 2)
\pgfmathparse@td@FPU{2*(\pgfutil@tmpBB*\pgfutil@tmpCC-\pgfutil@tmpBC*\pgfutil@tmpBC)}%
\let\pgfutil@tmpd\pgfmathresult%
%\typeout{debug: determinant is \pgfutil@tmpd.}%
%\pgfmathtruncatemacro{\pgfutil@tmpi}{abs(\pgfutil@tmpd)>0.02?1:0}%
\pgfmathparse@td@FPU{int(abs(\pgfutil@tmpd)>0.02?1:0)}%
\pgfmathtruncatemacro{\pgfutil@tmpi}{\pgfmathresult}%
%\typeout{debug: i is \pgfutil@tmpi.}%
\ifnum\pgfutil@tmpi=0\relax
\PackageWarning{3dtools}{The points approximately sit on a common line. The circumcircle cannot
be determined.}%
\else
%\pgfmathsetmacro{\pgfutil@tmpb}{(\pgfutil@tmpBB-\pgfutil@tmpBC)*\pgfutil@tmpCC/\pgfutil@tmpd}%
\pgfmathparse@td@FPU{(\pgfutil@tmpBB-\pgfutil@tmpBC)*\pgfutil@tmpCC/\pgfutil@tmpd}%
\let\pgfutil@tmpb\pgfmathresult
%\typeout{debug: b is \pgfutil@tmpb.}%
%\pgfmathsetmacro{\pgfutil@tmpc}{(\pgfutil@tmpCC-\pgfutil@tmpBC)*\pgfutil@tmpBB/\pgfutil@tmpd}%
\pgfmathparse@td@FPU{(\pgfutil@tmpCC-\pgfutil@tmpBC)*\pgfutil@tmpBB/\pgfutil@tmpd}%
\let\pgfutil@tmpc\pgfmathresult
%\typeout{debug: c is \pgfutil@tmpc.}%
%\typeout{debug: b=\pgfutil@tmpb,c=\pgfutil@tmpc,(B)=(\pgfutil@tmpB),(C)=(\pgfutil@tmpC)}%
\pgfmathsetmacro{\pgfutil@tmpI}{TD("\pgfutil@tmpA+\pgfutil@tmpb*(\pgfutil@tmpB)+\pgfutil@tmpc*(\pgfutil@tmpC)")}%
%\typeout{debug: (I)=((\pgfutil@tmpI))}%
\tikzset{insert path={(\pgfutil@tmpI)}}%
\fi
}}%
\tikzset{3d/define orthonormal dreibein/.code={%
\def\tikz@td@dreibein@ok{0}%
\tikzset{3d/aux keys/.cd,#1}%
\edef\pgfutil@tmpA{\pgfkeysvalueof{/tikz/3d/aux keys/A}}%
\edef\pgfutil@tmpB{\pgfkeysvalueof{/tikz/3d/aux keys/B}}%
\edef\pgfutil@tmpC{\pgfkeysvalueof{/tikz/3d/aux keys/C}}%
\pgfmathsetmacro{\pgfutil@tmpB}{TD("\pgfutil@tmpB-\pgfutil@tmpA")}%
\pgfmathsetmacro{\pgfutil@tmpC}{TD("\pgfutil@tmpC-\pgfutil@tmpA")}%
\pgfmathsetmacro{\pgfutil@tmpD}{TD("(\pgfutil@tmpB)x(\pgfutil@tmpC)")}%
\pgfmathsetmacro@td@FPU{\pgfutil@tmpd}{min(1,TD("(\pgfutil@tmpD)o(\pgfutil@tmpD)"))}% changed May 23 2023
\ifnum\fpeval{\pgfutil@tmpd<0.01}=1\relax
\PackageWarning{3dtools}{The vectors \pgfutil@tmpA, \pgfutil@tmpB\space and 
\pgfutil@tmpC are too collinear. Cannot determine appropriate plane.}%
\else
\def\tikz@td@dreibein@ok{1}%
\pgfmathsetmacro{\pgfutil@tmpB}{TDunit("(\pgfutil@tmpB)")}% changed May 23 2023
\pgfmathsetmacro{\pgfutil@tmpD}{TDunit("(\pgfutil@tmpD)")}% changed May 23 2023
\pgfmathsetmacro{\pgfutil@tmpC}{TD("(\pgfutil@tmpD)x(\pgfutil@tmpB)")}%
% \typeout{3d coordinate={\pgfkeysvalueof{/tikz/3d/aux keys/ex}=(\pgfutil@tmpB)},
% 	3d coordinate={\pgfkeysvalueof{/tikz/3d/aux keys/ey}=(\pgfutil@tmpC)},
% 	3d coordinate={\pgfkeysvalueof{/tikz/3d/aux keys/ez}=(\pgfutil@tmpD)}}%
\expanded{\noexpand\path[overlay,
	3d coordinate={\pgfkeysvalueof{/tikz/3d/aux keys/ex}=(\pgfutil@tmpB)},
	3d coordinate={\pgfkeysvalueof{/tikz/3d/aux keys/ey}=(\pgfutil@tmpC)},
	3d coordinate={\pgfkeysvalueof{/tikz/3d/aux keys/ez}=(\pgfutil@tmpD)}];}%
\fi
},3d/aux keys/.cd,A/.initial={(A)},B/.initial={(B)},
C/.initial={(C)},D/.initial={(D)},
intersection 1/.initial={I-1},intersection 2/.initial={I-2},
ex/.initial={(ex)},
ey/.initial={(ey)},ez/.initial={(ez)},
i1/.initial=i1,i2/.initial=i2,
rA/.initial=1,rB/.initial=1,rC/.initial=1,
tmpcoord1/.initial=3dtoolstmp-1}
\tikzset{3d/screen coords/.style={x={(1cm,0cm)},y={(0cm,1cm)},z={(-1cm,-1cm)}}}%
%%
%% predefined pics
\makeatletter
\tikzset{pics/3d/dim line/.style={code={%
  \tikzset{3d/dim line/.cd,#1}%
  \pgfmathsetmacro{\pgfutil@tmpa}{tddistance("\pgfkeysvalueof{/tikz/3d/dim line/A}","\pgfkeysvalueof{/tikz/3d/dim line/B}")}%
  \ifnum\fpeval{\pgfutil@tmpa<0.02}=1\relax 
   \message{The points A=\pgfkeysvalueof{/tikz/3d/dim line/A} and B=\pgfkeysvalueof{/tikz/3d/dim line/B} are too close. Cannot draw dim line.}%
  \else 
    \pgfmathsetmacro{\pgfutil@tmpa}{TD("\pgfkeysvalueof{/tikz/3d/dim line/A}-\pgfkeysvalueof{/tikz/3d/dim line/B}x\pgfkeysvalueof{/tikz/3d/dim line/n}")}%
    \pgfmathsetmacro{\pgfutil@tmpb}{TD("(\pgfutil@tmpa)o(\pgfutil@tmpa)")}%
    \ifnum\fpeval{\pgfutil@tmpb<0.02}=1\relax
      \message{The normal n=\pgfkeysvalueof{/tikz/3d/dim line/n} is either too short or too collinear with A-B=\pgfkeysvalueof{/tikz/3d/dim line/A}-\pgfkeysvalueof{/tikz/3d/dim line/B}. Cannot draw dim line.}%
    \else 
      % \pgfmathsetmacro{\pgfutil@tmpC}{TD("\pgfkeysvalueof{/tikz/3d/dim line/B}-\pgfkeysvalueof{/tikz/3d/dim line/A}x\pgfkeysvalueof{/tikz/3d/dim line/n}")}%
      % \pgfmathsetmacro{\pgfutil@tmpC}{TD("\pgfkeysvalueof{/tikz/3d/dim line/B}+(\pgfutil@tmpC)")}%
      % %\typeout{C=(\pgfutil@tmpC)}%
      % \tikzset{3d/define orthonormal dreibein={%
      %   A=\pgfkeysvalueof{/tikz/3d/dim line/A},%
      %   B=\pgfkeysvalueof{/tikz/3d/dim line/B},%
      %   C={(\pgfutil@tmpC)},%
      %   ex={(\pgfkeysvalueof{/tikz/3d/dim line/auxiliary coordinate prefix}ex)},%
      %   ey={(\pgfkeysvalueof{/tikz/3d/dim line/auxiliary coordinate prefix}ey)},%
      %   ez={(\pgfkeysvalueof{/tikz/3d/dim line/auxiliary coordinate prefix}ez)}%
      % }}%
      \pgfmathsetmacro{\pgfutil@tmpa}{TDunit("(\pgfutil@tmpa)")}%
      \pgfmathsetmacro{\pgfutil@tmpb}{\pgfkeysvalueof{/tikz/3d/dim line/d}/1cm}%
      \pgfmathsetmacro{\pgfutil@tmpex}{TDunit("\pgfkeysvalueof{/tikz/3d/dim line/B}-\pgfkeysvalueof{/tikz/3d/dim line/A}")}%
      \expanded{\noexpand\path[overlay]
        (\pgfutil@tmpex) coordinate (\pgfkeysvalueof{/tikz/3d/dim line/auxiliary coordinate prefix}ex);}
      \pgfmathtruncatemacro{\pgfutil@tmpf}{ifthenelse(x2d("\pgfkeysvalueof{/tikz/3d/dim line/auxiliary coordinate prefix}ex")<0,1,0)}%
      \ifnum\pgfutil@tmpf=1\relax 
        \pgfmathsetmacro{\pgfutil@tmpex}{TDunit("\pgfkeysvalueof{/tikz/3d/dim line/A}-\pgfkeysvalueof{/tikz/3d/dim line/B}")}%
      \fi
      \pgfmathsetmacro{\pgfutil@tmpey}{TDunit("\pgfkeysvalueof{/tikz/3d/dim line/n}x(\pgfutil@tmpex)")}%
      \expanded{\noexpand\path[overlay]
        (\pgfutil@tmpey) coordinate (\pgfkeysvalueof{/tikz/3d/dim line/auxiliary coordinate prefix}ey);}
      \pgfmathtruncatemacro{\pgfutil@tmpf}{ifthenelse(y2d("\pgfkeysvalueof{/tikz/3d/dim line/auxiliary coordinate prefix}ey")<0,1,0)}%
      \ifnum\pgfutil@tmpf=1\relax 
        \pgfmathsetmacro{\pgfutil@tmpey}{TDunit("(\pgfutil@tmpex)x\pgfkeysvalueof{/tikz/3d/dim line/n}")}%
      \fi
      \pgfmathsetmacro{\pgfutil@tmpez}{TDunit("\pgfkeysvalueof{/tikz/3d/dim line/n}")}%
      \pgfmathsetmacro{\pgfutil@tmpA}{TD("\pgfkeysvalueof{/tikz/3d/dim line/A}+\pgfutil@tmpb*(\pgfutil@tmpa)")}%
      \pgfmathsetmacro{\pgfutil@tmpB}{TD("\pgfkeysvalueof{/tikz/3d/dim line/B}+\pgfutil@tmpb*(\pgfutil@tmpa)")}%
      \pgfmathsetmacro{\pgfutil@tmpC}{TD("0.5*\pgfkeysvalueof{/tikz/3d/dim line/A}+0.5*\pgfkeysvalueof{/tikz/3d/dim line/B}+\pgfutil@tmpb*(\pgfutil@tmpa)")}%
      \expanded{\noexpand\path[overlay]
      % (\pgfutil@tmpa) coordinate (\pgfkeysvalueof{/tikz/3d/dim line/auxiliary coordinate prefix}D)
      (\pgfutil@tmpex) coordinate (\pgfkeysvalueof{/tikz/3d/dim line/auxiliary coordinate prefix}ex)
      (\pgfutil@tmpey) coordinate (\pgfkeysvalueof{/tikz/3d/dim line/auxiliary coordinate prefix}ey)
      (\pgfutil@tmpez) coordinate (\pgfkeysvalueof{/tikz/3d/dim line/auxiliary coordinate prefix}ez)
      (\pgfutil@tmpA) coordinate (\pgfkeysvalueof{/tikz/3d/dim line/auxiliary coordinate prefix}A')
      (\pgfutil@tmpB) coordinate (\pgfkeysvalueof{/tikz/3d/dim line/auxiliary coordinate prefix}B')
      (\pgfutil@tmpC) coordinate (\pgfkeysvalueof{/tikz/3d/dim line/auxiliary coordinate prefix}AB')
      ;}%
      \begin{scope}[x={(\pgfkeysvalueof{/tikz/3d/dim line/auxiliary coordinate prefix}ex)},
          y={(\pgfkeysvalueof{/tikz/3d/dim line/auxiliary coordinate prefix}ey)},
          z={(\pgfkeysvalueof{/tikz/3d/dim line/auxiliary coordinate prefix}ez)},canvas is xy plane at z=0]          
          \path (\pgfkeysvalueof{/tikz/3d/dim line/auxiliary coordinate prefix}AB') node[3d/dim line/dim label](\pgfkeysvalueof{/tikz/3d/dim line/node name}){\tikz@td@dim@line@label};
          \fill let \p1=($(\pgfkeysvalueof{/tikz/3d/dim line/node name}.west)-(\pgfkeysvalueof{/tikz/3d/dim line/node name}.east)$),\n1={0.5*sqrt((\x1)*(\x1)+(\y1)*(\y1))},\p2=($(\pgfkeysvalueof{/tikz/3d/dim line/node name})-(\pgfkeysvalueof{/tikz/3d/dim line/auxiliary coordinate prefix}A')$),\n2={sqrt((\x2)*(\x2)+(\y2)*(\y2))-\n1-0.4pt} in (\pgfkeysvalueof{/tikz/3d/dim line/node name}) ++ (\n1,0) -- ++ (0,0.2pt) -- ++ (\n2-3pt,0pt)
          -- ++ (0pt,1.3pt) -- ++ (3pt,-1.5pt) -- ++ (-3pt,-1.5pt)  -- ++ (0pt,1.3pt) -- ++ (3pt-\n2,0pt) -- cycle
          (\pgfkeysvalueof{/tikz/3d/dim line/node name}) ++ (-\n1,0) -- ++ (0,-0.2pt) -- ++ (3pt-\n2,0pt)
          -- ++ (0pt,-1.3pt) -- ++ (-3pt,1.5pt) -- ++ (3pt,1.5pt)  -- ++ (0pt,-1.3pt) -- ++ (\n2-3pt,0pt) -- cycle
          (\pgfkeysvalueof{/tikz/3d/dim line/node name}) ++ (\n1+\n2,1.5pt) -- ++ (0.4pt,0pt) -- ++ (0pt,-3pt) -- ++ (-0.4pt,0pt) -- cycle
          (\pgfkeysvalueof{/tikz/3d/dim line/node name}) ++ (0pt-\n1-\n2,-1.5pt) -- ++ (-0.4pt,0pt) -- ++ (0pt,3pt) -- ++ (0.4pt,0pt) -- cycle;
      \end{scope}
    \fi 
  \fi  
}},3d/dim line/.cd,A/.initial={(0,0,0)},B/.initial={(1,0,0)},n/.initial={(0,0,1)},
l/.store in=\tikz@td@dim@line@label,
l=\empty,
auxiliary coordinate prefix/.initial=tmp,d/.initial=2mm,
dim label/.style={transform shape,circle,inner sep=1pt},
node name/.initial=dim label}
% based on https://en.wikipedia.org/wiki/Circumscribed_circle
\tikzset{pics/3d circle through 3 points/.style={code={%
 \tikzset{3d/circle through 3 points/.cd,#1}%
 \expanded{\noexpand\path[overlay,
3d coordinate={(\pgfkeysvalueof{/tikz/3d/circle through 3 points/auxiliary coordinate prefix}b)=\pgfkeysvalueof{/tikz/3d/circle through 3 points/A}-\pgfkeysvalueof{/tikz/3d/circle through 3 points/C}},
 3d coordinate={(\pgfkeysvalueof{/tikz/3d/circle through 3 points/auxiliary coordinate prefix}a)=\pgfkeysvalueof{/tikz/3d/circle through 3 points/B}-\pgfkeysvalueof{/tikz/3d/circle through 3 points/C}},
 3d coordinate={(\pgfkeysvalueof{/tikz/3d/circle through 3 points/auxiliary coordinate prefix}c)=\pgfkeysvalueof{/tikz/3d/circle through 3 points/A}-\pgfkeysvalueof{/tikz/3d/circle through 3 points/B}},
 3d coordinate={(\pgfkeysvalueof{/tikz/3d/circle through 3 points/auxiliary coordinate prefix}n)=(\pgfkeysvalueof{/tikz/3d/circle through 3 points/auxiliary coordinate prefix}a)x(\pgfkeysvalueof{/tikz/3d/circle through 3 points/auxiliary coordinate prefix}b)}
 ];}%
 \pgfmathsetmacro@td@FPU{\lengthn}{sqrt(TD("(\pgfkeysvalueof{/tikz/3d/circle through 3 points/auxiliary coordinate prefix}n)o(\pgfkeysvalueof{/tikz/3d/circle through 3 points/auxiliary coordinate prefix}n)"))}
 \ifnum\fpeval{\lengthn<0.02}=1\relax
   \message{The points are (almost) on a line. Circle cannot be determined.}
  \else
   \pgfmathsetmacro@td@FPU{\tmpradius}{sqrt(TD("(\pgfkeysvalueof{/tikz/3d/circle through 3 points/auxiliary coordinate prefix}a)o(\pgfkeysvalueof{/tikz/3d/circle through 3 points/auxiliary coordinate prefix}a)"))*%
   	sqrt(TD("(\pgfkeysvalueof{/tikz/3d/circle through 3 points/auxiliary coordinate prefix}b)o(\pgfkeysvalueof{/tikz/3d/circle through 3 points/auxiliary coordinate prefix}b)"))*sqrt(TD("(\pgfkeysvalueof{/tikz/3d/circle through 3 points/auxiliary coordinate prefix}c)o(\pgfkeysvalueof{/tikz/3d/circle through 3 points/auxiliary coordinate prefix}c)"))/%
		(2*\lengthn)}
   \pgfmathsetmacro@td@FPU{\coeffa}{-1*TD("(\pgfkeysvalueof{/tikz/3d/circle through 3 points/auxiliary coordinate prefix}b)o(\pgfkeysvalueof{/tikz/3d/circle through 3 points/auxiliary coordinate prefix}b)")/(2*TD("(\pgfkeysvalueof{/tikz/3d/circle through 3 points/auxiliary coordinate prefix}n)o(\pgfkeysvalueof{/tikz/3d/circle through 3 points/auxiliary coordinate prefix}n)"))}		
   \pgfmathsetmacro@td@FPU{\coeffb}{TD("(\pgfkeysvalueof{/tikz/3d/circle through 3 points/auxiliary coordinate prefix}a)o(\pgfkeysvalueof{/tikz/3d/circle through 3 points/auxiliary coordinate prefix}a)")/(2*TD("(\pgfkeysvalueof{/tikz/3d/circle through 3 points/auxiliary coordinate prefix}n)o(\pgfkeysvalueof{/tikz/3d/circle through 3 points/auxiliary coordinate prefix}n)"))}
   \expanded{%
   \noexpand\path[overlay,3d coordinate={(\pgfkeysvalueof{/tikz/3d/circle through 3 points/auxiliary coordinate prefix}u)=\coeffa*(\pgfkeysvalueof{/tikz/3d/circle through 3 points/auxiliary coordinate prefix}a)x(\pgfkeysvalueof{/tikz/3d/circle through 3 points/auxiliary coordinate prefix}n)},
     3d coordinate={(\pgfkeysvalueof{/tikz/3d/circle through 3 points/auxiliary coordinate prefix}v)=\coeffb*(\pgfkeysvalueof{/tikz/3d/circle through 3 points/auxiliary coordinate prefix}b)x(\pgfkeysvalueof{/tikz/3d/circle through 3 points/auxiliary coordinate prefix}n)}];
   \noexpand\path[3d coordinate={(\pgfkeysvalueof{/tikz/3d/circle through 3 points/center name})=\pgfkeysvalueof{/tikz/3d/circle through 3 points/C}+(\pgfkeysvalueof{/tikz/3d/circle through 3 points/auxiliary coordinate prefix}u)+(\pgfkeysvalueof{/tikz/3d/circle through 3 points/auxiliary coordinate prefix}v)}];
   }%
   \pgfmathsetmacro{\normalizationa}{1/sqrt(TD("(\pgfkeysvalueof{/tikz/3d/circle through 3 points/auxiliary coordinate prefix}a)o(\pgfkeysvalueof{/tikz/3d/circle through 3 points/auxiliary coordinate prefix}a)"))}
   \pgfmathsetmacro{\normalizationn}{1/sqrt(TD("(\pgfkeysvalueof{/tikz/3d/circle through 3 points/auxiliary coordinate prefix}n)o(\pgfkeysvalueof{/tikz/3d/circle through 3 points/auxiliary coordinate prefix}n)"))}
   \expanded{%
   \noexpand\path[overlay,3d coordinate={(\pgfkeysvalueof{/tikz/3d/circle through 3 points/auxiliary coordinate prefix}a)=\normalizationa*(\pgfkeysvalueof{/tikz/3d/circle through 3 points/auxiliary coordinate prefix}a)},
   3d coordinate={(\pgfkeysvalueof{/tikz/3d/circle through 3 points/auxiliary coordinate prefix}n)=\normalizationn*(\pgfkeysvalueof{/tikz/3d/circle through 3 points/auxiliary coordinate prefix}n)},
   3d coordinate={(\pgfkeysvalueof{/tikz/3d/circle through 3 points/auxiliary coordinate prefix}c)=(\pgfkeysvalueof{/tikz/3d/circle through 3 points/auxiliary coordinate prefix}a)x(\pgfkeysvalueof{/tikz/3d/circle through 3 points/auxiliary coordinate prefix}n)}];
   }%
   \expanded{%
	\noexpand\begin{scope}[plane x={(\pgfkeysvalueof{/tikz/3d/circle through 3 points/auxiliary coordinate prefix}a)},plane y={(\pgfkeysvalueof{/tikz/3d/circle through 3 points/auxiliary coordinate prefix}c)},canvas is plane]
	 \noexpand\draw[pic actions] (\pgfkeysvalueof{/tikz/3d/circle through 3 points/center name}) circle[radius=\tmpradius];
	 \noexpand\pgfkeys{/tikz/3d/circle through 3 points/extra}
	\noexpand\end{scope}}%
  \fi 
}},3d/circle through 3 points/.cd,A/.initial={(1,0,0)},B/.initial={(0,1,0)},
C/.initial={(0,0,1)},
auxiliary coordinate prefix/.initial=tmp,center name/.initial=M,
extra/.code={}}%
% 
% incircle
\tikzset{pics/3d incircle/.style={code={%
\tikzset{3d/incircle/.cd,#1}%
\pgfmathsetmacro{\pgfutil@tmpa}{tddistance("\pgfkeysvalueof{/tikz/3d/incircle/B}","\pgfkeysvalueof{/tikz/3d/incircle/C}")}%
\pgfmathsetmacro{\pgfutil@tmpb}{tddistance("\pgfkeysvalueof{/tikz/3d/incircle/A}","\pgfkeysvalueof{/tikz/3d/incircle/C}")}%
\pgfmathsetmacro{\pgfutil@tmpc}{tddistance("\pgfkeysvalueof{/tikz/3d/incircle/A}","\pgfkeysvalueof{/tikz/3d/incircle/B}")}%
\pgfmathsetmacro{\pgfutil@tmpd}{(\pgfutil@tmpa+\pgfutil@tmpb+\pgfutil@tmpc)/2}%
%\edef\pgfutil@tmpt{\noexpand\tikz@td@parse@coord\pgfkeysvalueof{/tikz/3d/incircle/A}}%
%\pgfutil@tmpt
\expanded{\noexpand\tikz@td@parse@coord\pgfkeysvalueof{/tikz/3d/incircle/A}}%
\pgfmathsetmacro{\pgfutil@tmp@Ax}{\pgf@X}%
\pgfmathsetmacro{\pgfutil@tmp@Ay}{\pgf@Y}%
\pgfmathsetmacro{\pgfutil@tmp@Az}{\pgf@Z}%
%\edef\pgfutil@tmpt{\noexpand\tikz@td@parse@coord\pgfkeysvalueof{/tikz/3d/incircle/B}}%
%\pgfutil@tmpt
\expanded{\noexpand\tikz@td@parse@coord\pgfkeysvalueof{/tikz/3d/incircle/B}}%
\pgfmathsetmacro{\pgfutil@tmp@Bx}{\pgf@X}%
\pgfmathsetmacro{\pgfutil@tmp@By}{\pgf@Y}%
\pgfmathsetmacro{\pgfutil@tmp@Bz}{\pgf@Z}%
%\edef\pgfutil@tmpt{\noexpand\tikz@td@parse@coord\pgfkeysvalueof{/tikz/3d/incircle/C}}%
%\pgfutil@tmpt
\expanded{\noexpand\tikz@td@parse@coord\pgfkeysvalueof{/tikz/3d/incircle/C}}%
\pgfmathsetmacro{\pgfutil@tmp@Cx}{\pgf@X}%
\pgfmathsetmacro{\pgfutil@tmp@Cy}{\pgf@Y}%
\pgfmathsetmacro{\pgfutil@tmp@Cz}{\pgf@Z}%
\pgfmathsetmacro{\pgfutil@tmpIx}{(\pgfutil@tmp@Ax*\pgfutil@tmpa+\pgfutil@tmp@Bx*\pgfutil@tmpb+\pgfutil@tmp@Cx*\pgfutil@tmpc)/\pgfutil@tmpd/2}%
\pgfmathsetmacro{\pgfutil@tmpIy}{(\pgfutil@tmp@Ay*\pgfutil@tmpa+\pgfutil@tmp@By*\pgfutil@tmpb+\pgfutil@tmp@Cy*\pgfutil@tmpc)/\pgfutil@tmpd/2}%
\pgfmathsetmacro{\pgfutil@tmpIz}{(\pgfutil@tmp@Az*\pgfutil@tmpa+\pgfutil@tmp@Bz*\pgfutil@tmpb+\pgfutil@tmp@Cz*\pgfutil@tmpc)/\pgfutil@tmpd/2}%
\path (\pgfutil@tmpIx,\pgfutil@tmpIy,\pgfutil@tmpIz) 
  coordinate (\pgfkeysvalueof{/tikz/3d/incircle/center name})
  [3d/projection of={(\pgfkeysvalueof{/tikz/3d/incircle/center name}) on 
	\pgfkeysvalueof{/tikz/3d/incircle/A}--\pgfkeysvalueof{/tikz/3d/incircle/B}}] 
 coordinate (\pgfkeysvalueof{/tikz/3d/incircle/projection prefix}c)
 [3d/projection of={(\pgfkeysvalueof{/tikz/3d/incircle/center name}) on 
	\pgfkeysvalueof{/tikz/3d/incircle/A}--\pgfkeysvalueof{/tikz/3d/incircle/C}}] 
 coordinate (\pgfkeysvalueof{/tikz/3d/incircle/projection prefix}b)
 [3d/projection of={(\pgfkeysvalueof{/tikz/3d/incircle/center name}) on 
	\pgfkeysvalueof{/tikz/3d/incircle/B}--\pgfkeysvalueof{/tikz/3d/incircle/C}}] 
 coordinate (\pgfkeysvalueof{/tikz/3d/incircle/projection prefix}a);
 \path[pic actions] pic{3d circle through 3 points={%
   	A={(\pgfkeysvalueof{/tikz/3d/incircle/projection prefix}c)},%
   	B={(\pgfkeysvalueof{/tikz/3d/incircle/projection prefix}b)},%
  	C={(\pgfkeysvalueof{/tikz/3d/incircle/projection prefix}a)},%
  	center name=\pgfkeysvalueof{/tikz/3d/incircle/center name}}}
	;
}},3d/incircle/.cd,
center name/.initial=I,
A/.initial={(1,0,0)},B/.initial={(0,1,0)},C/.initial={(0,0,1)},
projection prefix/.initial=tmpp}
%
% Circles and arcs in 3d, split where their visibility changes. Vectors are
% comma lists x,y,z in the coordinates of the current frame.
% \tikz@td@circle@frame{<normal>}: \tikz@td@kz is the unit normal, \tikz@td@kN
% the unit screen normal and \tikz@td@kw the length of its component in the
% plane of the circle; \tikz@td@kx is that component normalized (it points to
% the part of the circle closest to the viewer) or, if the circle faces the
% screen (\tikz@td@kw<0.04), any unit vector in the plane; \tikz@td@ky = kz x kx.
\ExplSyntaxOn
\cs_new:Npn \__tdtools_vcomb:nnnn #1#2#3#4
  { \exp_last_unbraced:Ne \__tdtools_vcomb:w { #2 , #4 } \q_stop {#1} {#3} }
\cs_new:Npn \__tdtools_vcomb:w #1 , #2 , #3 , #4 , #5 , #6 \q_stop #7#8
  {
    \fp_eval:n { (#7) * (#1) + (#8) * (#4) } ,
    \fp_eval:n { (#7) * (#2) + (#8) * (#5) } ,
    \fp_eval:n { (#7) * (#3) + (#8) * (#6) }
  }
\cs_new:Npn \__tdtools_vcross:nn #1#2
  { \exp_last_unbraced:Ne \__tdtools_vcross:w { #1 , #2 } \q_stop }
\cs_new:Npn \__tdtools_vcross:w #1 , #2 , #3 , #4 , #5 , #6 \q_stop
  {
    \fp_eval:n { (#2) * (#6) - (#3) * (#5) } ,
    \fp_eval:n { (#3) * (#4) - (#1) * (#6) } ,
    \fp_eval:n { (#1) * (#5) - (#2) * (#4) }
  }
\cs_new:Npn \__tdtools_vnorm:n #1 { sqrt ( \__tdtools_dot:nn {#1} {#1} ) }
\cs_new_protected:Npn \__tdtools_vunit:Nn #1#2
  { \tl_set:Ne #1 { \__tdtools_vcomb:nnnn { 1 / \__tdtools_vnorm:n {#2} } {#2} { 0 } {#2} } }
\cs_generate_variant:Nn \__tdtools_vunit:Nn { Ne }
\cs_new_protected:Npn \tikz@td@circle@frame #1
  {
    \__tdtools_vunit:Ne \tikz@td@kz {#1}
    \__tdtools_vunit:Ne \tikz@td@kN { \__tdtools_nx: , \__tdtools_ny: , \__tdtools_nz: }
    \tl_set:Ne \tikz@td@kx
      {
        \__tdtools_vcomb:nnnn { 1 } { \tikz@td@kN }
          { - ( \__tdtools_dot:nn { \tikz@td@kN } { \tikz@td@kz } ) } { \tikz@td@kz }
      }
    \tl_set:Ne \tikz@td@kw { \fp_eval:n { \__tdtools_vnorm:n { \tikz@td@kx } } }
    \fp_compare:nNnT { \tikz@td@kw } < { 0.04 }
      { \pgfmathsetmacro \tikz@td@kx { normalcandidate ( \tikz@td@kz ) } }
    \__tdtools_vunit:Ne \tikz@td@kx { \tikz@td@kx }
    \tl_set:Ne \tikz@td@ky { \__tdtools_vcross:nn { \tikz@td@kz } { \tikz@td@kx } }
  }
% \tikz@td@circle@arcs{<c>}{<start>}{<end>}{<u>}: after \tikz@td@circle@frame,
% calls \tikz@td@circle@seg{<a>}{<b>}{<visible>} for the parts of the arc
% from <start> to <end> (degrees from the direction <u> projected into the
% plane, towards kz x u) with the angle t from kx in [<a>,<b>], where <visible>
% is 1 if cos t > <c>. A full circle starts at a change of visibility.
\seq_new:N \l__tdtools_cuts_seq
\tl_new:N \l__tdtools_segs_tl
\tl_new:N \l__tdtools_prev_tl
\fp_new:N \l__tdtools_ka_fp
\fp_new:N \l__tdtools_kb_fp
\fp_new:N \l__tdtools_kf_fp
\cs_new_protected:Npn \tikz@td@circle@arcs #1#2#3#4
  {
    \fp_compare:nNnTF { abs (#1) } < { 1 }
      { \fp_set:Nn \l__tdtools_kf_fp { acosd (#1) } }
      { \fp_set:Nn \l__tdtools_kf_fp { 0 } }
    \fp_compare:nNnTF { (#3) - (#2) } < { 359.999 }
      {
        \tl_set:Ne \l__tdtools_tmp_tl
          { \__tdtools_vcomb:nnnn { 1 } {#4} { - ( \__tdtools_dot:nn {#4} { \tikz@td@kz } ) } { \tikz@td@kz } }
        \fp_compare:nNnT { \__tdtools_vnorm:n { \l__tdtools_tmp_tl } } < { 0.001 }
          { \tl_set_eq:NN \l__tdtools_tmp_tl \tikz@td@kx }
        \fp_set:Nn \l__tdtools_ka_fp
          {
            atand ( \__tdtools_dot:nn { \l__tdtools_tmp_tl } { \tikz@td@ky } ,
              \__tdtools_dot:nn { \l__tdtools_tmp_tl } { \tikz@td@kx } )
          }
        \fp_set:Nn \l__tdtools_kb_fp { \l__tdtools_ka_fp + (#3) }
        \fp_add:Nn \l__tdtools_ka_fp {#2}
      }
      {
        \fp_set_eq:NN \l__tdtools_ka_fp \l__tdtools_kf_fp
        \fp_set:Nn \l__tdtools_kb_fp { \l__tdtools_ka_fp + 360 }
      }
    \seq_clear:N \l__tdtools_cuts_seq
    \fp_compare:nNnT { 0 < \l__tdtools_kf_fp < 180 } = { 1 }
      {
        \int_step_inline:nnn
          { \fp_eval:n { floor ( ( \l__tdtools_ka_fp - 180 ) / 360 ) } }
          { \fp_eval:n { ceil ( ( \l__tdtools_kb_fp + 180 ) / 360 ) } }
          {
            \__tdtools_cut:n { 360 * ##1 - \l__tdtools_kf_fp }
            \__tdtools_cut:n { 360 * ##1 + \l__tdtools_kf_fp }
          }
      }
    \seq_sort:Nn \l__tdtools_cuts_seq
      { \fp_compare:nNnTF {##1} > {##2} { \sort_return_swapped: } { \sort_return_same: } }
    \seq_put_left:Ne \l__tdtools_cuts_seq { \fp_use:N \l__tdtools_ka_fp }
    \seq_put_right:Ne \l__tdtools_cuts_seq { \fp_use:N \l__tdtools_kb_fp }
    \tl_set:Ne \tikz@td@kf { \fp_use:N \l__tdtools_kf_fp }
    \tl_set:Ne \tikz@td@kga { \fp_use:N \l__tdtools_ka_fp }
    \tl_set:Ne \tikz@td@kgb { \fp_use:N \l__tdtools_kb_fp }
    \tl_clear:N \l__tdtools_segs_tl
    \tl_clear:N \l__tdtools_prev_tl
    \seq_map_inline:Nn \l__tdtools_cuts_seq
      {
        \tl_if_empty:NF \l__tdtools_prev_tl
          {
            \fp_compare:nNnT { ##1 - \l__tdtools_prev_tl } > { 0.001 }
              {
                \tl_put_right:Ne \l__tdtools_segs_tl
                  {
                    \exp_not:N \tikz@td@circle@seg { \l__tdtools_prev_tl } {##1}
                      { \fp_compare:nNnTF { cosd ( ( \l__tdtools_prev_tl + ##1 ) / 2 ) } > {#1} { 1 } { 0 } }
                  }
              }
          }
        \tl_set:Nn \l__tdtools_prev_tl {##1}
      }
    \tl_use:N \l__tdtools_segs_tl
  }
\cs_new_protected:Npn \__tdtools_cut:n #1
  {
    \fp_compare:nNnT { \l__tdtools_ka_fp + 0.001 < #1 < \l__tdtools_kb_fp - 0.001 } = { 1 }
      { \seq_put_right:Ne \l__tdtools_cuts_seq { \fp_eval:n {#1} } }
  }
\ExplSyntaxOff
% \tikz@td@circle@draw{<center>}{<r>}{<c>}{<start>}{<end>}{<u>}: after
% \tikz@td@circle@frame, draws the arc of radius <r> around <center> (see
% \tikz@td@circle@arcs) in the local frame kx, ky, kz. The disk is filled with
% \tikz@td@kfill first; visible parts use \tikz@td@kvis on layer
% \tikz@td@kvlay, the others \tikz@td@khid on \tikz@td@khlay. Every part is a
% piece of its own.
\def\tikz@td@circle@draw#1#2#3#4#5#6{%
  \pgfmathsetmacro\tikz@td@kr{#2}%
  \ifx\tikz@td@kfront\relax\let\tikz@td@kext\tikz@td@kr\else\def\tikz@td@kext{0}\fi
  % canvas images of the frame vectors (x={(<3d point>)} would include shifts)
  \tikz@scan@one@point\pgfutil@firstofone(\tikz@td@kx)\relax
  \edef\tikz@td@kcx{\the\pgf@x,\the\pgf@y}%
  \tikz@scan@one@point\pgfutil@firstofone(\tikz@td@ky)\relax
  \edef\tikz@td@kcy{\the\pgf@x,\the\pgf@y}%
  \tikz@scan@one@point\pgfutil@firstofone(\tikz@td@kz)\relax
  \edef\tikz@td@kcz{\the\pgf@x,\the\pgf@y}%
  \begin{scope}[shift={(#1)},x/.expanded={(\tikz@td@kcx)},
    y/.expanded={(\tikz@td@kcy)},z/.expanded={(\tikz@td@kcz)}]
    \tikz@td@piece{(0,0,0)}{\edef\pgfutil@tmpa{\noexpand\path[\tikz@td@kfill]}%
      \pgfutil@tmpa (0,0) circle[radius=\tikz@td@kr];}%
    \tikz@td@circle@arcs{#3}{#4}{#5}{#6}%
  \end{scope}}%
\def\tikz@td@circle@seg#1#2#3{%
  \ifnum#3=1 \edef\pgfutil@tmpb{\tikz@td@kvlay}\let\tikz@td@kst\tikz@td@kvis
  \else\edef\pgfutil@tmpb{\tikz@td@khlay}\let\tikz@td@kst\tikz@td@khid\fi
  % \pgfonlayer resets the line width, so enter it only to change the layer
  \ifx\pgfutil@tmpb\pgfonlayer@name\def\tikz@td@klayer{\begingroup}\def\tikz@td@kendlayer{\endgroup}%
  \else\def\tikz@td@klayer{\pgfonlayer{\pgfutil@tmpb}}\def\tikz@td@kendlayer{\endpgfonlayer}\fi
  \tikz@td@klayer
  \def\tikz@td@kcls{#3}%
  \iftikz@td@scene\def\pgfutil@tmpa{1}\else
    \ifx\tikz@td@kgA\pgfutil@empty\def\pgfutil@tmpa{0}\else\def\pgfutil@tmpa{1}\fi
  \fi
  \ifnum\pgfutil@tmpa=0
    % one path
    \edef\pgfutil@tmpa{\noexpand\draw[\tikz@td@kst]}%
    \tikz@td@circle@arc{#1}{#2}{#1}{#2}%
  \else
    % sub-arcs: in a scene each is a piece, with the depth of its cell (the cells
    % are \tikz@td@kgs degrees wide, start at the angle 0 and end where the
    % visibility changes), so that pieces of coinciding circles tie; with a color
    % gradient they are at most \tikz@td@kgg degrees long. Only the last one keeps
    % the arrow tips, and it spans at least 10 degrees, so that it can be
    % shortened for them.
    \ifx\tikz@td@kgA\pgfutil@empty\let\tikz@td@kgh\tikz@td@kgs\else\let\tikz@td@kgh\tikz@td@kgg\fi
    \ifdim\fpeval{\tikz@td@kgb-(#2)}pt<0.001pt
      \edef\tikz@td@kge{\fpeval{#2-min(#2-(#1),max(10,\tikz@td@kgh))}}%
    \else\def\tikz@td@kge{#2}\fi
    \def\tikz@td@kta{#1}%
    \tikz@td@circle@sub{#2}%
  \fi
  \tikz@td@kendlayer}%
% the sub-arcs from \tikz@td@kta on: up to \tikz@td@kge at multiples of
% \tikz@td@kgh, then up to #1
\def\tikz@td@circle@sub#1{%
  \ifdim\tikz@td@kta pt<\fpeval{\tikz@td@kge-0.001}pt
    \edef\tikz@td@ktb{\fpeval{min(\tikz@td@kge,(floor(round(\tikz@td@kta/\tikz@td@kgh,6))+1)*\tikz@td@kgh)}}%
  \else
    \edef\tikz@td@ktb{#1}%
  \fi
  \ifx\tikz@td@kgA\pgfutil@empty\else
    \edef\pgfutil@tmpc{\fpeval{round(100*((\tikz@td@kta+\tikz@td@ktb)/2-\tikz@td@kga)/(\tikz@td@kgb-\tikz@td@kga),2)}}%
    \colorlet{tikz@td@kgc}{\tikz@td@kgB!\pgfutil@tmpc!\tikz@td@kgA}%
  \fi
  \edef\pgfutil@tmpa{\noexpand\draw[\tikz@td@kst,/tikz/3d/joint caps
    \ifx\tikz@td@kgA\pgfutil@empty\else,color=tikz@td@kgc\fi
    \ifdim\fpeval{\tikz@td@kgb-\tikz@td@ktb}pt<0.001pt \else,arrows=-\fi]}%
  % the cell of the middle of the sub-arc
  \edef\tikz@td@kcm{\fpeval{(\tikz@td@kta+\tikz@td@ktb)/2}}%
  \edef\tikz@td@kca{\fpeval{floor(\tikz@td@kcm/\tikz@td@kgs)*\tikz@td@kgs}}%
  \edef\tikz@td@kcb{\fpeval{\tikz@td@kca+\tikz@td@kgs}}%
  \ifdim\fpeval{0<\tikz@td@kf<180}pt=1pt
    \ifnum\tikz@td@kcls=1
      \edef\pgfutil@tmpc{\fpeval{round(\tikz@td@kcm/360)*360}}%
      \edef\tikz@td@kca{\fpeval{max(\tikz@td@kca,\pgfutil@tmpc-\tikz@td@kf)}}%
      \edef\tikz@td@kcb{\fpeval{min(\tikz@td@kcb,\pgfutil@tmpc+\tikz@td@kf)}}%
    \else
      \edef\pgfutil@tmpc{\fpeval{floor((\tikz@td@kcm-\tikz@td@kf)/360)*360}}%
      \edef\tikz@td@kca{\fpeval{max(\tikz@td@kca,\pgfutil@tmpc+\tikz@td@kf)}}%
      \edef\tikz@td@kcb{\fpeval{min(\tikz@td@kcb,\pgfutil@tmpc+360-\tikz@td@kf)}}%
    \fi
  \fi
  \tikz@td@circle@arc{\tikz@td@kta}{\tikz@td@ktb}{\tikz@td@kca}{\tikz@td@kcb}%
  \let\tikz@td@kta\tikz@td@ktb
  \ifdim\tikz@td@kta pt<\fpeval{#1-0.001}pt\expandafter\tikz@td@circle@sub\expandafter{#1}\fi}%
% \tikz@td@circle@arc{<a>}{<b>}{<c>}{<d>}: \pgfutil@tmpa (a \draw command)
% applied to the arc from <a> to <b>, as a piece at the depth of the centroid
% of the arc from <c> to <d>. Lines on a sphere cover each other only where
% they cross, at equal depth: the visible ones all go just in front of the
% front shell of the sphere (\tikz@td@kfront), in the order of the code, the
% others just in front of the back shell (\tikz@td@kback), in the reverse
% order (what is on top in front is at the bottom behind).
\def\tikz@td@circle@arc#1#2#3#4{%
  \ifdim\fpeval{#4-(#3)}pt>359.999pt
    \def\pgfutil@tmpc{(0,0,0)}%
  \else
    \edef\pgfutil@tmpc{(\fpeval{\tikz@td@kr*(sind(#4)-sind(#3))/((#4-(#3))*pi/180)},%
      \fpeval{\tikz@td@kr*(cosd(#3)-cosd(#4))/((#4-(#3))*pi/180)},0)}%
  \fi
  \ifx\tikz@td@kfront\relax\else
    \ifnum\tikz@td@kcls=1 \edef\pgfutil@tmpc{\fpeval{\tikz@td@kfront+0.01}}%
    \else\edef\pgfutil@tmpc{\fpeval{\tikz@td@kback+0.01-0.000001*(\the\tikz@td@npieces+1)}}\fi
  \fi
  \ifdim\fpeval{#2-(#1)}pt>359.999pt
    \expandafter\tikz@td@piece\expandafter{\pgfutil@tmpc}{\pgfutil@tmpa (0,0) circle[radius=\tikz@td@kr];}%
  \else
    \expandafter\tikz@td@piece\expandafter{\pgfutil@tmpc}{\pgfutil@tmpa (#1:\tikz@td@kr)
      arc[start angle=#1,end angle=#2,radius=\tikz@td@kr];}%
  \fi}%
\let\tikz@td@kfront\relax
\def\tikz@td@circle@gradient#1{%
  \edef\tikz@td@kgA{\pgfkeysvalueof{#1/start color}}%
  \edef\tikz@td@kgB{\pgfkeysvalueof{#1/end color}}%
  \ifx\tikz@td@kgB\pgfutil@empty\let\tikz@td@kgB\tikz@td@kgA\fi
  \pgfmathsetmacro\tikz@td@kgs{\pgfkeysvalueof{#1/segment angle}}%
  \pgfmathsetmacro\tikz@td@kgg{min(\tikz@td@kgs,\pgfkeysvalueof{#1/gradient angle})}}%
\tikzset{pics/3d/circle on sphere/.style={code={%
 \tikzset{3d/circle on sphere/.cd,#1}%
 \def\pgfutil@tmpo##1{\pgfkeysvalueof{/tikz/3d/circle on sphere/##1}}%
 \pgfmathsetmacro{\tikz@td@kR}{\pgfutil@tmpo{R}}%
 \pgfmathsetmacro{\tikz@td@kC}{TD("\pgfutil@tmpo{C}")}%
 \pgfmathsetmacro{\tikz@td@kP}{TD("\pgfutil@tmpo{P}")}%
 \pgfmathsetmacro{\tikz@td@kn}{TD("\pgfutil@tmpo{n}")}%
 \edef\pgfutil@tmpa{\pgfutil@tmpo{latitude}}%
 \ifx\pgfutil@tmpa\pgfutil@empty\else
  \def\tikz@td@kn{0,0,1}%
  \edef\tikz@td@kP{\csname __tdtools_vcomb:nnnn\endcsname{1}{\tikz@td@kC}%
   {\tikz@td@kR*sind(\pgfutil@tmpa)}{0,0,1}}%
 \fi
 \edef\pgfutil@tmpa{\pgfutil@tmpo{longitude}}%
 \ifx\pgfutil@tmpa\pgfutil@empty\else
  \edef\tikz@td@kn{\fpeval{-sind(\pgfutil@tmpa)},\fpeval{cosd(\pgfutil@tmpa)},0}%
  \let\tikz@td@kP\tikz@td@kC
 \fi
 \edef\pgfutil@tmpb{\csname __tdtools_vcomb:nnnn\endcsname{1}{\tikz@td@kP}{-1}{\tikz@td@kC}}%
 \edef\pgfutil@tmpd{\fpeval{sqrt(\csname __tdtools_dot:nn\endcsname{\pgfutil@tmpb}{\pgfutil@tmpb})}}%
 \ifdim\pgfutil@tmpd pt>0.02pt\relax\let\tikz@td@kn\pgfutil@tmpb\fi
 \ifdim\fpeval{\csname __tdtools_dot:nn\endcsname{\tikz@td@kn}{\tikz@td@kn}}pt<0.0004pt\relax
  \PackageWarning{3dtools}{Chosen normal is too short.}%
 \else\ifdim\pgfutil@tmpd pt<\tikz@td@kR pt\relax
  \pgfmathsetmacro{\tikz@td@kr}{sqrt(\tikz@td@kR*\tikz@td@kR-\pgfutil@tmpd*\pgfutil@tmpd)}%
  \tikz@td@circle@frame{\tikz@td@kn}%
  % the coordinates (ex), (ey), (ez), (nscreen) and (<prefix>P)
  \path[overlay,3d coordinate/.expanded={(nscreen)=(\tikz@td@kN)},
   3d coordinate/.expanded={\pgfkeysvalueof{/tikz/3d/aux keys/ex}=(\tikz@td@kx)},
   3d coordinate/.expanded={\pgfkeysvalueof{/tikz/3d/aux keys/ey}=(\tikz@td@ky)},
   3d coordinate/.expanded={\pgfkeysvalueof{/tikz/3d/aux keys/ez}=(\tikz@td@kz)},
   3d coordinate/.expanded={(\pgfutil@tmpo{auxiliary coordinate prefix}P)=(\tikz@td@kP)}];%
  % visible where the depth exceeds the one of the center of the sphere
  \edef\pgfutil@tmpc{\fpeval{\csname __tdtools_dot:nn\endcsname{\pgfutil@tmpb}{\tikz@td@kN}}}%
  \ifdim\tikz@td@kw pt<0.04pt\relax
   \ifdim\pgfutil@tmpc pt<0pt\def\pgfutil@tmpc{2}\else\def\pgfutil@tmpc{-2}\fi
  \else
   \edef\pgfutil@tmpc{\fpeval{-\pgfutil@tmpc/(\tikz@td@kr*\tikz@td@kw)}}%
  \fi
  \def\tikz@td@kvis{/tikz/3d/visible}\def\tikz@td@khid{/tikz/3d/hidden}%
  \def\tikz@td@knosphere{1}%
  \tikz@td@circle@gradient{/tikz/3d/circle on sphere}%
  % in a scene, the parts lie on the front and back shells of the sphere
  % (3d/sphere)
  \edef\pgfutil@tmpa{(\csname __tdtools_vcomb:nnnn\endcsname{1}{\tikz@td@kC}{\tikz@td@kR}{\tikz@td@kN})}%
  \expandafter\tikz@td@depthof\expandafter{\pgfutil@tmpa}%
  \let\tikz@td@kfront\tikz@td@depthval
  \edef\pgfutil@tmpa{(\csname __tdtools_vcomb:nnnn\endcsname{1}{\tikz@td@kC}{-\tikz@td@kR}{\tikz@td@kN})}%
  \expandafter\tikz@td@depthof\expandafter{\pgfutil@tmpa}%
  \let\tikz@td@kback\tikz@td@depthval
  \def\tikz@td@kfill{/tikz/3d/circle on sphere/fill}%
  \edef\tikz@td@kvlay{\pgfutil@tmpo{fore layer}}\edef\tikz@td@khlay{\pgfutil@tmpo{back layer}}%
  \pgfmathsetmacro{\tikz@td@ku}{TD("\pgfutil@tmpo{u}")}%
  \tikz@td@circle@draw{\tikz@td@kP}{\tikz@td@kr}{\pgfutil@tmpc}%
   {\pgfutil@tmpo{start angle}}{\pgfutil@tmpo{end angle}}{\tikz@td@ku}%
 \else
  \PackageWarning{3dtools}{Circle center is not inside the sphere.}%
 \fi\fi
}},3d/circle on sphere/.cd,C/.initial={(0,0,0)},
	P/.initial={(0,0,0)},R/.initial=1,%
	n/.initial={(0,0,1)},auxiliary coordinate prefix/.initial=tmp,%
	fore layer/.initial=main,back layer/.initial=main,%
	latitude/.initial=,longitude/.initial=,fill/.style={},%
	start angle/.initial=0,end angle/.initial=360,u/.initial={(1,0,0)},%
	start color/.initial=,end color/.initial=,segment angle/.initial=360,gradient angle/.initial=2}
% a circle (or arc) around the pic position, split into the half closer to
% the viewer (3d/circle/front) and the other one (3d/circle/back)
\tikzset{pics/3d/circle/.style={code={%
 \tikzset{3d/circle/.cd,#1}%
 \def\pgfutil@tmpo##1{\pgfkeysvalueof{/tikz/3d/circle/##1}}%
 \pgfmathsetmacro{\tikz@td@kn}{TD("\pgfutil@tmpo{n}")}%
 \pgfmathsetmacro{\tikz@td@ku}{TD("\pgfutil@tmpo{u}")}%
 \tikz@td@circle@frame{\tikz@td@kn}%
 \ifdim\tikz@td@kw pt<0.04pt\relax\def\pgfutil@tmpc{-2}\else\def\pgfutil@tmpc{0}\fi
 \def\tikz@td@kvis{/tikz/3d/circle/front}\def\tikz@td@khid{/tikz/3d/circle/back}%
 \tikz@td@circle@gradient{/tikz/3d/circle}%
 \def\tikz@td@kfill{/tikz/3d/circle/fill}%
 \def\tikz@td@kvlay{main}\def\tikz@td@khlay{main}%
 \tikz@td@circle@draw{0,0,0}{\pgfutil@tmpo{r}}{\pgfutil@tmpc}%
  {\pgfutil@tmpo{start angle}}{\pgfutil@tmpo{end angle}}{\tikz@td@ku}%
}},3d/circle/.cd,r/.initial=1,n/.initial={(0,0,1)},u/.initial={(1,0,0)},%
	start angle/.initial=0,end angle/.initial=360,%
	front/.style={},back/.style={},fill/.style={},%
	start color/.initial=,end color/.initial=,segment angle/.initial=360,gradient angle/.initial=2}
%
% a sphere of radius 3d/sphere/R around the pic position, painted with
% 3d/sphere/shell (on the screen) and outlined with 3d/sphere/outline. In a
% scene the shell is painted twice, as a back and a front shell, so that what
% is inside the sphere lies between them (for a translucent sphere the
% opacity of the shell applies to each one), and the outline right after the
% back shell.
\tikzset{pics/3d/sphere/.style={code={%
 \tikzset{3d/sphere/.cd,#1}%
 \pgfmathsetmacro{\tikz@td@kR}{\pgfkeysvalueof{/tikz/3d/sphere/R}}%
 \let\tikz@td@kext\tikz@td@kR
 \def\tikz@td@knosphere{1}%
 \pgfmathsetmacro{\tikz@td@kN}{TDunit("(nscreenx,nscreeny,nscreenz)")}%
 % in a scene, the sphere places the pieces inside, on and around it:
 % {x}{y}{depth of the center}{radius}{front depth}{back depth}
 \iftikz@td@scene
  \tikz@td@depthof{(0,0,0)}\let\tikz@td@kc\tikz@td@depthpt\let\tikz@td@kcd\tikz@td@depthval
  \edef\pgfutil@tmpa{(\csname __tdtools_vcomb:nnnn\endcsname{\tikz@td@kR}{\tikz@td@kN}{0}{\tikz@td@kN})}%
  \expandafter\tikz@td@depthof\expandafter{\pgfutil@tmpa}\let\tikz@td@kfd\tikz@td@depthval
  \edef\pgfutil@tmpa{(\csname __tdtools_vcomb:nnnn\endcsname{-\tikz@td@kR}{\tikz@td@kN}{0}{\tikz@td@kN})}%
  \expandafter\tikz@td@depthof\expandafter{\pgfutil@tmpa}\let\tikz@td@kbd\tikz@td@depthval
  \expandafter\tikz@td@sphere@gput\expandafter{\expandafter\tikz@td@sphere@data\tikz@td@kc}%
 \fi
 % back shell, then the outline (same depth), so that all the rest covers it
 \edef\pgfutil@tmpa{(\csname __tdtools_vcomb:nnnn\endcsname{-\tikz@td@kR}{\tikz@td@kN}{0}{\tikz@td@kN})}%
 \expandafter\tikz@td@piece\expandafter{\pgfutil@tmpa}{%
  \path[/tikz/3d/sphere/shell,/tikz/3d/screen coords] (0,0) circle[radius=\tikz@td@kR];}%
 \expandafter\tikz@td@piece\expandafter{\pgfutil@tmpa}{%
  \path[/tikz/3d/sphere/outline,/tikz/3d/screen coords] (0,0) circle[radius=\tikz@td@kR];}%
 \iftikz@td@scene
  \edef\pgfutil@tmpa{(\csname __tdtools_vcomb:nnnn\endcsname{\tikz@td@kR}{\tikz@td@kN}{0}{\tikz@td@kN})}%
  \expandafter\tikz@td@piece\expandafter{\pgfutil@tmpa}{%
   \path[/tikz/3d/sphere/shell,/tikz/3d/screen coords] (0,0) circle[radius=\tikz@td@kR];}%
 \fi
}},3d/sphere/.cd,R/.initial=1,shell/.style={ball color=white},outline/.style={draw}}
\def\tikz@td@sphere@data#1#2#3{{#1}{#2}{\tikz@td@kcd}{\fpeval{\tikz@td@kfd-\tikz@td@kcd}}{\tikz@td@kfd}{\tikz@td@kbd}}%
% a point on a sphere: a node with 3d/point on sphere/front (and front text)
% if the point is on the side of the sphere that faces the viewer, else with
% back (and back text)
\tikzset{pics/3d/point on sphere/.style={code={%
 \tikzset{3d/point on sphere/.cd,#1}%
 \def\pgfutil@tmpo##1{\pgfkeysvalueof{/tikz/3d/point on sphere/##1}}%
 \pgfmathsetmacro{\tikz@td@kC}{TD("\pgfutil@tmpo{C}")}%
 \pgfmathsetmacro{\tikz@td@kP}{TD("\pgfutil@tmpo{P}")}%
 \pgfmathsetmacro{\pgfutil@tmpa}{TD("(\tikz@td@kP)-(\tikz@td@kC)o(nscreenx,nscreeny,nscreenz)")}%
 \ifdim\pgfutil@tmpa pt<0pt\def\pgfutil@tmpb{back}\else\def\pgfutil@tmpb{front}\fi
 \path (\tikz@td@kP) node[/tikz/3d/point on sphere/\pgfutil@tmpb]
  {\pgfutil@tmpo{\pgfutil@tmpb\space text}};
}},3d/point on sphere/.cd,C/.initial={(0,0,0)},P/.initial={(0,0,1)},%
	front/.style={},back/.style={},%
	front text/.initial=\textbullet,back text/.initial=\(\circ\)}
% cones and frusta
% \tikz@td@cone@frame{<axis>}{<s>} computes, for the axis direction <axis>
% and a signed length <s> along it (pointing from the base center to the
% apex), the unit axis \tikz@td@ca, \tikz@td@cnz = axis o n_screen,
% \tikz@td@cs = |axis x n_screen| and an orthonormal frame \tikz@td@ce,
% \tikz@td@cf of the base plane. \tikz@td@cf lies in the screen plane, and
% \tikz@td@ce is oriented such that the projection of the apex lies on the
% positive x axis of the frame, so the arc through angle 0 is always the one
% on the apex side.
\def\tikz@td@cone@frame#1#2{%
  \pgfmathsetmacro{\tikz@td@ca}{TDunit("#1")}%
  \pgfmathsetmacro{\tikz@td@cnz}{TD("(\tikz@td@ca)o(nscreenx,nscreeny,nscreenz)")}%
  \pgfmathsetmacro{\tikz@td@cf}{TD("(\tikz@td@ca)x(nscreenx,nscreeny,nscreenz)")}%
  \pgfmathsetmacro{\tikz@td@cs}{sqrt(abs(TD("(\tikz@td@cf)o(\tikz@td@cf)")))}%
  \ifnum\fpeval{\tikz@td@cs<0.001}=1\relax
   % looking along the axis: any vector orthogonal to the axis will do
   \pgfmathsetmacro{\tikz@td@cs}{0}%
   \pgfmathsetmacro{\tikz@td@cf}{normalcandidate(\tikz@td@ca)}%
   \pgfmathsetmacro{\tikz@td@cf}{TD("(\tikz@td@ca)x(\tikz@td@cf)")}%
  \fi
  \pgfmathsetmacro{\tikz@td@cf}{TDunit("(\tikz@td@cf)")}%
  \pgfmathsetmacro{\tikz@td@ce}{TD("(\tikz@td@cf)x(\tikz@td@ca)")}%
  \pgfmathtruncatemacro{\tikz@td@ci}{(#2)*\tikz@td@cnz*TD("(\tikz@td@ce)o(nscreenx,nscreeny,nscreenz)")>0?1:0}%
  \ifnum\tikz@td@ci=1\relax
   \pgfmathsetmacro{\tikz@td@ce}{TD("-1*(\tikz@td@ce)")}%
  \fi
  % screen images of the frame vectors (for the x and y keys)
  \tikz@scan@one@point\pgfutil@firstofone(\tikz@td@ce)\relax
  \edef\tikz@td@cx{\the\pgf@x,\the\pgf@y}%
  \tikz@scan@one@point\pgfutil@firstofone(\tikz@td@cf)\relax
  \edef\tikz@td@cy{\the\pgf@x,\the\pgf@y}}%
% after \tikz@td@cone@frame: \tikz@td@kw is the unit vector in the base plane
% that points to the viewer
\def\tikz@td@cone@viewer{%
  \pgfmathsetmacro{\tikz@td@kq}{TD("(\tikz@td@ce)o(nscreenx,nscreeny,nscreenz)")}%
  \ifdim\tikz@td@kq pt<0pt\pgfmathsetmacro{\tikz@td@kw}{TD("-1*(\tikz@td@ce)")}%
  \else\let\tikz@td@kw\tikz@td@ce\fi}%
% \tikz@td@kdepth\macro{<a>}{<b>}: \macro is the depth of the point
% <a>*\tikz@td@ca + <b>*\tikz@td@kw (unit axis and \tikz@td@kw)
\def\tikz@td@kdepth#1#2#3{%
  \edef\tikz@td@kpt{(\csname __tdtools_vcomb:nnnn\endcsname{#2}{\tikz@td@ca}{#3}{\tikz@td@kw})}%
  \expandafter\tikz@td@depthof\expandafter{\tikz@td@kpt}%
  \let#1\tikz@td@depthval}%
\def\tikz@td@cone@code#1{% #1=1: shaded
  \def\tikz@td@cshaded{#1}%
  \pgfmathsetmacro{\tikz@td@r}{\pgfkeysvalueof{/tikz/3d/r}}%
  \pgfmathsetmacro{\tikz@td@h}{\pgfkeysvalueof{/tikz/3d/h}}%
  \tikz@td@cone@frame{\pgfkeysvalueof{/tikz/3d/cone/axis}}{\tikz@td@h}%
  \pgfmathsetmacro{\tikz@td@cp}{TD("\tikz@td@h*(\tikz@td@ca)")}%
  \path[overlay] (0,0,0) coordinate (-base) (\tikz@td@cp) coordinate (-apex)
   coordinate (-tip);% -tip: old name of the apex
  % screen angle of the direction from the base center to the apex
  \tikz@scan@one@point\pgfutil@firstofone(\tikz@td@cp)\relax
  \edef\pgfutil@tmpa{\the\pgf@x}\edef\pgfutil@tmpb{\the\pgf@y}%
  \pgfmathsetmacro{\tikz@td@cp}{atan2(\pgfutil@tmpb,\pgfutil@tmpa)}%
  \pgfkeyslet{/tikz/3d/cone/screen angle}{\tikz@td@cp}%
  \ifnum#1=1\relax % light model for 3d/cone/outer shade and inner shade
   \tikz@td@cone@light{\tikz@td@ce}{\tikz@td@cf}{\tikz@td@ca}{0,0,0}%
  \fi
  % base face oriented towards the viewer?
  \pgfmathtruncatemacro{\tikz@td@ci}{\tikz@td@h*\tikz@td@cnz<0?1:0}%
  % projection of the apex outside the projection of the base?
  \pgfmathtruncatemacro{\tikz@td@cj}{abs(\tikz@td@h)*\tikz@td@cs>abs(\tikz@td@r*\tikz@td@cnz)?1:0}%
  % depths of the pieces: base center, centroid of the front half of the
  % mantle, middle of the rim behind; the radius breaks ties
  \pgfmathsetmacro\tikz@td@kext{abs(\tikz@td@r)}%
  \tikz@td@cone@viewer
  \tikz@td@kdepth\tikz@td@kdi{0}{0}%
  \tikz@td@kdepth\tikz@td@kdo{\tikz@td@h/3}{4*abs(\tikz@td@r)/(3*pi)}%
  \tikz@td@kdepth\tikz@td@kdh{0}{-2*abs(\tikz@td@r)/pi}%
  \begin{scope}[x/.expanded={(\tikz@td@cx)},y/.expanded={(\tikz@td@cy)}]
   \ifnum\tikz@td@cj=1\relax
    \pgfmathsetmacro{\tikz@td@ac}{acos(abs(\tikz@td@r*\tikz@td@cnz)/(abs(\tikz@td@h)*\tikz@td@cs))}%
    \ifnum\tikz@td@ci=1\relax
     % base (or interior) and the part of the mantle outside the base
     \tikz@td@piece{\tikz@td@kdi}{\path[/tikz/3d/cone/td@region=inner] (0,0) circle[radius=\tikz@td@r];}%
     \tikz@td@piece{\tikz@td@kdo}{\path[/tikz/3d/cone/td@region=outer] (\tikz@td@ac:\tikz@td@r)
      arc[start angle=\tikz@td@ac,end angle=-\tikz@td@ac,radius=\tikz@td@r]
      -- (-apex) -- cycle;}%
    \else
     % mantle, the base rim on the apex side is hidden
     \tikz@td@piece{\tikz@td@kdo}{\path[/tikz/3d/cone/td@region=outer] (\tikz@td@ac:\tikz@td@r)
      arc[start angle=\tikz@td@ac,end angle=360-\tikz@td@ac,radius=\tikz@td@r]
      -- (-apex) -- cycle;}%
     \tikz@td@piece{\tikz@td@kdh}{\path[/tikz/3d/hidden,/tikz/3d/edge only] (\tikz@td@ac:\tikz@td@r)
      arc[start angle=\tikz@td@ac,end angle=-\tikz@td@ac,radius=\tikz@td@r];}%
    \fi
   \else
    \ifnum\tikz@td@ci=1\relax
     \tikz@td@piece{\tikz@td@kdi}{\path[/tikz/3d/cone/td@region=inner] (0,0) circle[radius=\tikz@td@r];}%
    \else
     \tikz@td@piece{\tikz@td@kdo}{\path[/tikz/3d/cone/td@region=outer] (0,0) circle[radius=\tikz@td@r];}%
    \fi
   \fi
  \end{scope}}%
\tikzset{pics/3d/cone/.style={code={%
	\tikzset{3d/.cd,#1}%
	\tikz@td@cone@code{0}%
	}},
	pics/3d/shaded cone/.style={code={%
	\tikzset{3d/.cd,#1}%
	\tikz@td@cone@code{1}%
	}},
	3d/frustum/outer/.style={},3d/frustum/inner/.style={},
	3d/cone/.cd,outer/.style={},inner/.style={},
	% visible region of a cone (outer: mantle, inner: base/interior)
	td@region/.code={\tikzset{/tikz/3d/visible,/tikz/3d/cone/#1}%
	 \ifnum\tikz@td@cshaded=1\relax
	  \tikzset{path picture={\path[/tikz/3d/cone/#1 shade,/tikz/3d/screen coords]
	   (-apex) circle[radius={abs(\tikz@td@r)+abs(\tikz@td@h)}];}}%
	 \fi},
	outer shade/.style={3d/cone/lit side=outer},
	inner shade/.style={3d/cone/lit side=inner},
	axis/.initial={(0,0,1)},
	apex/.code={% apex relative to the base center; sets axis and h
	 \pgfmathsetmacro{\pgfutil@tmpa}{TD("#1")}%
	 \pgfmathsetmacro{\pgfutil@tmpb}{sqrt(abs(TD("(\pgfutil@tmpa)o(\pgfutil@tmpa)")))}%
	 \pgfkeyslet{/tikz/3d/h}{\pgfutil@tmpb}%
	 \pgfkeys{/tikz/3d/cone/axis/.expanded={(\pgfutil@tmpa)}}},
	screen angle/.initial=90}
%%
% \tikz@td@kdepthe\macro{<a>}{<b>}: \macro is the depth of the point
% <a>*\tikz@td@ca + <b>*\tikz@td@ce
\def\tikz@td@kdepthe#1#2#3{%
  \edef\tikz@td@kpt{(\csname __tdtools_vcomb:nnnn\endcsname{#2}{\tikz@td@ca}{#3}{\tikz@td@ce})}%
  \expandafter\tikz@td@depthof\expandafter{\tikz@td@kpt}%
  \let#1\tikz@td@depthval}%
% \tikz@td@frustum@arc{<rim B or S>}{<height>}{<radius>}{<from>}{<to>}{<mean
% cosine>}{<visible or hidden>}: an arc of a rim of the frustum as a piece
\def\tikz@td@frustum@arc#1#2#3#4#5#6#7{%
  \tikz@td@kdepthe\tikz@td@kd{#2}{(#3)*(#6)}%
  \edef\tikz@td@kpC{\noexpand\path[/tikz/3d/#7,/tikz/3d/edge only]}%
  \expandafter\tikz@td@piece\expandafter{\tikz@td@kd}{\tikz@td@kpC
    ($(-tdtmp#1)+(#4:#3)$) arc[start angle=#4,end angle=#5,radius=#3];}}%
\tikzset{pics/3d/frustum/.style={code={%
  \tikzset{3d/.cd,#1}%
  \pgfmathsetmacro{\tikz@td@R}{\pgfkeysvalueof{/tikz/3d/R}}%
  \pgfmathsetmacro{\tikz@td@r}{\pgfkeysvalueof{/tikz/3d/r}}%
  \pgfmathsetmacro{\tikz@td@h}{\pgfkeysvalueof{/tikz/3d/h}}%
  \pgfmathsetmacro{\tikz@td@cp}{TDunit("\pgfkeysvalueof{/tikz/3d/cone/axis}")}%
  \pgfmathsetmacro{\tikz@td@cp}{TD("\tikz@td@h*(\tikz@td@cp)")}%
  \path[overlay] (0,0,0) coordinate (-base) (\tikz@td@cp) coordinate (-upper base);
  % the (virtual) apex; omitted for near-cylinders, where it lies so far
  % away that placing it would exceed TeX's dimension range
  \ifnum\fpeval{abs(\tikz@td@R-\tikz@td@r)>0.001
    && abs(\tikz@td@R*\tikz@td@h/(\tikz@td@R-\tikz@td@r))<=100}=1\relax
   \pgfmathsetmacro{\tikz@td@cp}{TD("\tikz@td@R/(\tikz@td@R-\tikz@td@r)*(\tikz@td@cp)")}%
   \path[overlay] (\tikz@td@cp) coordinate (-tip);
  \fi
  % big rim B (radius \tikz@td@Rb) and small rim S (radius \tikz@td@Rs);
  % \tikz@td@d is the signed distance from B to S along the axis
  \pgfmathtruncatemacro{\tikz@td@ci}{\tikz@td@R<\tikz@td@r?1:0}%
  \ifnum\tikz@td@ci=1\relax
   \let\tikz@td@Rb\tikz@td@r \let\tikz@td@Rs\tikz@td@R
   \pgfmathsetmacro{\tikz@td@d}{-1*\tikz@td@h}%
   \path[overlay] (-upper base) coordinate (-tdtmpB) (-base) coordinate (-tdtmpS);
  \else
   \let\tikz@td@Rb\tikz@td@R \let\tikz@td@Rs\tikz@td@r
   \let\tikz@td@d\tikz@td@h
   \path[overlay] (-base) coordinate (-tdtmpB) (-upper base) coordinate (-tdtmpS);
  \fi
  \tikz@td@cone@frame{\pgfkeysvalueof{/tikz/3d/cone/axis}}{\tikz@td@d}%
  % big face oriented towards the viewer? (otherwise the small one is)
  \pgfmathtruncatemacro{\tikz@td@ci}{\tikz@td@d*\tikz@td@cnz<0?1:0}%
  % tangents exist?
  \pgfmathtruncatemacro{\tikz@td@cj}{abs(\tikz@td@d)*\tikz@td@cs>abs((\tikz@td@Rb-\tikz@td@Rs)*\tikz@td@cnz)?1:0}%
  % pieces, with the depths of their centroids: the two patches of the mantle
  % between the generatrices of the silhouette (A around the angle 0, i.e. on
  % the apex side, and B), one shows its outer side, the other its inner side
  % (styles 3d/frustum/outer and inner); the four arcs of the rims and the two
  % generatrices (3d/visible or 3d/hidden). The ends are open.
  \pgfmathsetmacro\tikz@td@kext{max(abs(\tikz@td@R),abs(\tikz@td@r))}%
  \ifdim\tikz@td@R pt<\tikz@td@r pt\def\tikz@td@kzB{\tikz@td@h}\def\tikz@td@kzS{0}%
  \else\def\tikz@td@kzB{0}\def\tikz@td@kzS{\tikz@td@h}\fi
  \pgfmathsetmacro\tikz@td@krm{(abs(\tikz@td@Rb)+abs(\tikz@td@Rs))/2}%
  % does patch A show its outer side? (outward normal at the angle 0)
  \pgfmathsetmacro{\tikz@td@kq}{TD("(\tikz@td@ce)o(nscreenx,nscreeny,nscreenz)")}%
  \pgfmathsetmacro{\tikz@td@kq}{abs(\tikz@td@d)*\tikz@td@kq+(\tikz@td@d<0?-1:1)*(\tikz@td@Rb-\tikz@td@Rs)*\tikz@td@cnz}%
  \ifdim\tikz@td@kq pt>0pt\def\tikz@td@kA{outer}\def\tikz@td@kB{inner}%
  \else\def\tikz@td@kA{inner}\def\tikz@td@kB{outer}\fi
  \begin{scope}[x/.expanded={(\tikz@td@cx)},y/.expanded={(\tikz@td@cy)}]
   \ifnum\tikz@td@cj=1\relax
    \pgfmathsetmacro{\tikz@td@ac}{acos(abs((\tikz@td@Rb-\tikz@td@Rs)*\tikz@td@cnz)/(abs(\tikz@td@d)*\tikz@td@cs))}%
    % mean cosine of the angle over [-ac,ac] and over [ac,360-ac]
    \pgfmathsetmacro\tikz@td@kmA{sin(\tikz@td@ac)/(max(\tikz@td@ac,0.01)*pi/180)}%
    \pgfmathsetmacro\tikz@td@kmB{-sin(\tikz@td@ac)/(max(180-\tikz@td@ac,0.01)*pi/180)}%
    \tikz@td@kdepthe\tikz@td@kdA{\tikz@td@h/2}{\tikz@td@krm*\tikz@td@kmA}%
    \tikz@td@kdepthe\tikz@td@kdB{\tikz@td@h/2}{\tikz@td@krm*\tikz@td@kmB}%
    \edef\tikz@td@kpA{\noexpand\path[/tikz/3d/frustum/\tikz@td@kA]}%
    \edef\tikz@td@kpB{\noexpand\path[/tikz/3d/frustum/\tikz@td@kB]}%
    \expandafter\tikz@td@piece\expandafter{\tikz@td@kdA}{\tikz@td@kpA
      ($(-tdtmpB)+(-\tikz@td@ac:\tikz@td@Rb)$)
      arc[start angle=-\tikz@td@ac,end angle=\tikz@td@ac,radius=\tikz@td@Rb]
      -- ($(-tdtmpS)+(\tikz@td@ac:\tikz@td@Rs)$)
      arc[start angle=\tikz@td@ac,end angle=-\tikz@td@ac,radius=\tikz@td@Rs] -- cycle;}%
    \expandafter\tikz@td@piece\expandafter{\tikz@td@kdB}{\tikz@td@kpB
      ($(-tdtmpB)+(\tikz@td@ac:\tikz@td@Rb)$)
      arc[start angle=\tikz@td@ac,end angle=360-\tikz@td@ac,radius=\tikz@td@Rb]
      -- ($(-tdtmpS)+(360-\tikz@td@ac:\tikz@td@Rs)$)
      arc[start angle=360-\tikz@td@ac,end angle=\tikz@td@ac,radius=\tikz@td@Rs] -- cycle;}%
    % rims: hidden is the part of S on the side B (if the big face is towards
    % the viewer), else the part of B on the side A
    \tikz@td@frustum@arc{B}{\tikz@td@kzB}{\tikz@td@Rb}{-\tikz@td@ac}{\tikz@td@ac}{\tikz@td@kmA}%
      {\ifnum\tikz@td@ci=1 visible\else hidden\fi}%
    \tikz@td@frustum@arc{B}{\tikz@td@kzB}{\tikz@td@Rb}{\tikz@td@ac}{360-\tikz@td@ac}{\tikz@td@kmB}{visible}%
    \tikz@td@frustum@arc{S}{\tikz@td@kzS}{\tikz@td@Rs}{-\tikz@td@ac}{\tikz@td@ac}{\tikz@td@kmA}{visible}%
    \tikz@td@frustum@arc{S}{\tikz@td@kzS}{\tikz@td@Rs}{\tikz@td@ac}{360-\tikz@td@ac}{\tikz@td@kmB}%
      {\ifnum\tikz@td@ci=1 hidden\else visible\fi}%
    % generatrices of the silhouette
    \tikz@td@kdepthe\tikz@td@kdg{\tikz@td@h/2}{\tikz@td@krm*cos(\tikz@td@ac)}%
    \expandafter\tikz@td@piece\expandafter{\tikz@td@kdg}{\path[/tikz/3d/visible,/tikz/3d/edge only]
      ($(-tdtmpB)+(\tikz@td@ac:\tikz@td@Rb)$) -- ($(-tdtmpS)+(\tikz@td@ac:\tikz@td@Rs)$)
      ($(-tdtmpB)+(-\tikz@td@ac:\tikz@td@Rb)$) -- ($(-tdtmpS)+(-\tikz@td@ac:\tikz@td@Rs)$);}%
   \else
    % seen (almost) along the axis: the whole mantle shows one side
    \ifnum\tikz@td@ci=1\relax\def\tikz@td@kA{inner}\else\def\tikz@td@kA{outer}\fi
    \tikz@td@kdepthe\tikz@td@kdA{\tikz@td@h/2}{0}%
    \edef\tikz@td@kpA{\noexpand\path[even odd rule,/tikz/3d/frustum/\tikz@td@kA]}%
    \expandafter\tikz@td@piece\expandafter{\tikz@td@kdA}{\tikz@td@kpA
      (-tdtmpB) circle[radius=\tikz@td@Rb] (-tdtmpS) circle[radius=\tikz@td@Rs];}%
    \tikz@td@kdepthe\tikz@td@kdB{\tikz@td@kzB}{0}%
    \tikz@td@kdepthe\tikz@td@kdS{\tikz@td@kzS}{0}%
    \edef\tikz@td@kpB{\noexpand\path[/tikz/3d/\ifnum\tikz@td@ci=1 hidden\else visible\fi,/tikz/3d/edge only]}%
    \expandafter\tikz@td@piece\expandafter{\tikz@td@kdS}{\tikz@td@kpB (-tdtmpS) circle[radius=\tikz@td@Rs];}%
    \expandafter\tikz@td@piece\expandafter{\tikz@td@kdB}{\path[/tikz/3d/visible,/tikz/3d/edge only]
      (-tdtmpB) circle[radius=\tikz@td@Rb];}%
   \fi
  \end{scope}
}}}
%		
% intersection of three spheres (see https://stackoverflow.com/a/18654302)
% 
%     temp1 = P2-P1                                        
%     e_x = temp1/norm(temp1)                              
%     temp2 = P3-P1                                        
%     i = dot(e_x,temp2)                                   
%     temp3 = temp2 - i*e_x                                
%     e_y = temp3/norm(temp3)                              
%     e_z = cross(e_x,e_y)                                 
%     d = norm(P2-P1)                                      
%     j = dot(e_y,temp2)                                   
%     x = (r1*r1 - r2*r2 + d*d) / (2*d)                    
%     y = (r1*r1 - r3*r3 -2*i*x + i*i + j*j) / (2*j)       
%     temp4 = r1*r1 - x*x - y*y                            
%     if temp4<0:                                          
%         raise Exception("The three spheres do not intersect!");
%     z = sqrt(temp4)                                      
%     p_12_a = P1 + x*e_x + y*e_y + z*e_z                  
%     p_12_b = P1 + x*e_x + y*e_y - z*e_z                  
%     return p_12_a,p_12_b                       
\tikzset{3d/.cd,intersection of three spheres/.code={%
\tikzset{3d/aux keys/.cd,#1,/tikz/3d/define orthonormal dreibein}%
\ifnum\tikz@td@dreibein@ok=0
\PackageWarning{3dtools}{The sphere centers are (almost) collinear. Cannot intersect the three spheres.}%
\else
\pgfmathsetmacro\pgfutil@tmpA{TD("\pgfkeysvalueof{/tikz/3d/aux keys/A}")}%
\pgfmathsetmacro\pgfutil@tmpB{TD("\pgfkeysvalueof{/tikz/3d/aux keys/B}")}%
\pgfmathsetmacro\pgfutil@tmpC{TD("\pgfkeysvalueof{/tikz/3d/aux keys/C}")}%
\pgfmathsetmacro\pgfutil@tmpex{TD("\pgfkeysvalueof{/tikz/3d/aux keys/ex}")}%
\pgfmathsetmacro\pgfutil@tmpey{TD("\pgfkeysvalueof{/tikz/3d/aux keys/ey}")}%
\pgfmathsetmacro\pgfutil@tmpez{TD("\pgfkeysvalueof{/tikz/3d/aux keys/ez}")}%
\pgfmathsetmacro\pgfutil@tmpAB{TD("(\pgfutil@tmpB)-(\pgfutil@tmpA)")}%
\pgfmathsetmacro\pgfutil@tmpd{sqrt(TD("(\pgfutil@tmpAB)o(\pgfutil@tmpAB)"))}%
\ifnum\fpeval{\pgfutil@tmpd<0.1}=1\relax
\PackageWarning{3dtools}{Points \pgfkeysvalueof{/tikz/3d/aux keys/A} and \pgfkeysvalueof{/tikz/3d/aux keys/B} are too close to each other.}%
\else
\pgfmathsetmacro\pgfutil@tmpAC{TD("(\pgfutil@tmpC)-(\pgfutil@tmpA)")}% temp2
%     i = dot(e_x,temp2)                                   
\pgfmathsetmacro\pgfutil@tmpi{TD("(\pgfutil@tmpex)o(\pgfutil@tmpAC)")}%
%     j = dot(e_y,temp2)                                   
\pgfmathsetmacro\pgfutil@tmpj{TD("(\pgfutil@tmpey)o(\pgfutil@tmpAC)")}%
\edef\pgfutil@tmprA{\pgfkeysvalueof{/tikz/3d/aux keys/rA}}%
\edef\pgfutil@tmprB{\pgfkeysvalueof{/tikz/3d/aux keys/rB}}%
\edef\pgfutil@tmprC{\pgfkeysvalueof{/tikz/3d/aux keys/rC}}%
%     x = (r1*r1 - r2*r2 + d*d) / (2*d)                    
\pgfmathsetmacro\pgfutil@tmpx{(pow(\pgfutil@tmprA,2)-pow(\pgfutil@tmprB,2)%
+pow(\pgfutil@tmpd,2))/(2*(\pgfutil@tmpd))}%
%     y = (r1*r1 - r3*r3 -2*i*x + i*i + j*j) / (2*j)       
\pgfmathsetmacro\pgfutil@tmpy{(\pgfutil@tmprA*\pgfutil@tmprA-\pgfutil@tmprC*\pgfutil@tmprC%
-2*(\pgfutil@tmpi)*(\pgfutil@tmpx)+(\pgfutil@tmpi)*(\pgfutil@tmpi)%
+(\pgfutil@tmpj)*(\pgfutil@tmpj))/(2*(\pgfutil@tmpj))}%
%     temp4 = r1*r1 - x*x - y*y                            
\pgfmathsetmacro\pgfutil@tmpc{\pgfutil@tmprA*\pgfutil@tmprA-pow(\pgfutil@tmpx,2)-pow(\pgfutil@tmpy,2)}%
\ifnum\fpeval{\pgfutil@tmpc<0}=1\relax
\PackageWarning{3dtools}{Not all spheres intersect.}%
\else
\pgfmathsetmacro\pgfutil@tmpz{sqrt(\pgfutil@tmpc)}%
\pgfmathsetmacro\pgfutil@tmpd{TD("(\pgfutil@tmpA)+\pgfutil@tmpx*(\pgfutil@tmpex)%
+\pgfutil@tmpy*(\pgfutil@tmpey)+\pgfutil@tmpz*(\pgfutil@tmpez)")}%
\pgfmathsetmacro\pgfutil@tmpe{TD("(\pgfutil@tmpA)+\pgfutil@tmpx*(\pgfutil@tmpex)%
+\pgfutil@tmpy*(\pgfutil@tmpey)-\pgfutil@tmpz*(\pgfutil@tmpez)")}%
\tikzset{insert path={%
(\pgfutil@tmpd) coordinate(\pgfkeysvalueof{/tikz/3d/aux keys/i1}) 
(\pgfutil@tmpe) coordinate(\pgfkeysvalueof{/tikz/3d/aux keys/i2}) 
}}
\fi
\fi
\fi
}}%
%
\tikzset{3d/insphere center/.code={%
 \tikzset{3d/insphere/.cd,#1}%
 \pgfmathsetmacro{\pgfutil@tmpA}{TD("\pgfkeysvalueof{/tikz/3d/insphere/C}-\pgfkeysvalueof{/tikz/3d/insphere/B}x\pgfkeysvalueof{/tikz/3d/insphere/D}-\pgfkeysvalueof{/tikz/3d/insphere/B}")}%
 \pgfmathsetmacro{\pgfutil@tmpA}{sqrt(TD("(\pgfutil@tmpA)o(\pgfutil@tmpA)"))/2}%
 \pgfmathsetmacro{\pgfutil@tmpB}{TD("\pgfkeysvalueof{/tikz/3d/insphere/C}-\pgfkeysvalueof{/tikz/3d/insphere/A}x\pgfkeysvalueof{/tikz/3d/insphere/D}-\pgfkeysvalueof{/tikz/3d/insphere/A}")}%
 \pgfmathsetmacro{\pgfutil@tmpB}{sqrt(TD("(\pgfutil@tmpB)o(\pgfutil@tmpB)"))/2}%
 \pgfmathsetmacro{\pgfutil@tmpC}{TD("\pgfkeysvalueof{/tikz/3d/insphere/B}-\pgfkeysvalueof{/tikz/3d/insphere/A}x\pgfkeysvalueof{/tikz/3d/insphere/D}-\pgfkeysvalueof{/tikz/3d/insphere/A}")}%
 \pgfmathsetmacro{\pgfutil@tmpC}{sqrt(TD("(\pgfutil@tmpC)o(\pgfutil@tmpC)"))/2}%
 \pgfmathsetmacro{\pgfutil@tmpD}{TD("\pgfkeysvalueof{/tikz/3d/insphere/A}-\pgfkeysvalueof{/tikz/3d/insphere/B}x\pgfkeysvalueof{/tikz/3d/insphere/C}-\pgfkeysvalueof{/tikz/3d/insphere/B}")}%
 \pgfmathsetmacro{\pgfutil@tmpD}{sqrt(TD("(\pgfutil@tmpD)o(\pgfutil@tmpD)"))/2}%
 \pgfmathsetmacro{\pgfutil@tmps}{\pgfutil@tmpA+\pgfutil@tmpB+\pgfutil@tmpC+\pgfutil@tmpD}%
 \pgfmathsetmacro{\pgfutil@tmpA}{\pgfutil@tmpA/\pgfutil@tmps}%
 \pgfmathsetmacro{\pgfutil@tmpB}{\pgfutil@tmpB/\pgfutil@tmps}%
 \pgfmathsetmacro{\pgfutil@tmpC}{\pgfutil@tmpC/\pgfutil@tmps}%
 \pgfmathsetmacro{\pgfutil@tmpD}{\pgfutil@tmpD/\pgfutil@tmps}%
 \pgfmathsetmacro{\pgfutil@tmpI}{TD("\pgfutil@tmpA*\pgfkeysvalueof{/tikz/3d/insphere/A}+\pgfutil@tmpB*\pgfkeysvalueof{/tikz/3d/insphere/B}+\pgfutil@tmpC*\pgfkeysvalueof{/tikz/3d/insphere/C}+\pgfutil@tmpD*\pgfkeysvalueof{/tikz/3d/insphere/D}")}%
 \tikzset{insert path={(\pgfutil@tmpI)}}
},3d/insphere/.cd,A/.initial={(A)},B/.initial={(B)},
C/.initial={(C)},D/.initial={(D)}}%
%
\newif\if@tikz@td@complete@dashes
\@tikz@td@complete@dashesfalse
\tikzset{cheating dash/.code args={on #1 off #2}{% from https://tex.stackexchange.com/a/133357
        % Use csname so catcode of @ doesn't have do be changed.		
        \tikz@addoption{%
            \pgfgetpath\pgfutil@tmpc%
            \pgfprocessround{\pgfutil@tmpc}{\pgfutil@tmpc}%
            \pgf@decorate@parsesoftpath{\pgfutil@tmpc}{\pgfutil@tmpc}%
            \pgfmathsetmacro{\pgfutil@tmpa}{\pgf@decorate@totalpathlength-#1}%\let\rest=\pgfmathresult%
            \pgfmathsetmacro{\pgfutil@tmpb}{#1+#2}%\let\onoff=\pgfmathresult%
            \pgfmathsetmacro{\pgfutil@tmpd}{max(floor(\pgfutil@tmpa/\pgfutil@tmpb), 1)}%\let\nfullonoff=\pgfmathresult%
            \pgfmathsetmacro{\pgfutil@tmpe}{max((\pgfutil@tmpa-\pgfutil@tmpb*\pgfutil@tmpd)/\pgfutil@tmpd+#2,#2)}%\let\offexpand=\pgfmathresult%
            \pgfsetdash{{#1}{\pgfutil@tmpe}}{0pt}}%
    },cheating dash/.default={on 2pt off 2pt}}%
\tikzset{3d/.cd,% 
reset vertex counter/.code={\c@pgf@counta0\relax},
add vertex/.code={\advance\c@pgf@counta by1\relax
\path (#1) coordinate (\the\c@pgf@counta);},
define vertices/.code={\c@pgf@counta0\relax
\tikzset{3d/add vertex/.list={#1}}},%
append vertices/.code={\tikzset{3d/add vertex/.list={#1}}},%
}	
\tikzset{3d/.cd,polyhedron/.cd,
face edges/.style={},
edges have complete dashes/.is if=@tikz@td@complete@dashes,
edges have complete dashes/.default=true,
complete dashes/.code args={on #1 off #2}{\tikzset{%
/tikz/3d/polyhedron/edges have complete dashes=true,
/tikz/3d/polyhedron/face edges/.append style={cheating dash=on #1 off #2}}},
complete dashes/.default={on 2pt off 2pt},
draw face with corners/.code={\begingroup\c@pgf@counta0\relax
\pgfutil@for\pgfutil@tmpa:={#1}\do{%
\pgfmathsetmacro\pgfutil@tmpT{TD("\pgfutil@tmpa")}%
\ifcase\c@pgf@counta 
\pgfmathsetmacro\pgfutil@tmpA{TD("\pgfutil@tmpa")}%
\or
\pgfmathsetmacro\pgfutil@tmpB{TD("\pgfutil@tmpa")}%
\or
\pgfmathsetmacro\pgfutil@tmpC{TD("\pgfutil@tmpa")}%
\fi
\advance\c@pgf@counta by1\relax
\ifnum\c@pgf@counta=1\relax
\edef\pgfutil@tmpl{(\pgfutil@tmpT)}% regular list
\edef\pgfutil@tmplf{(\pgfutil@tmpT)}% first
\edef\pgfutil@tmple{(\pgfutil@tmpT)}% edge list
\else
\edef\pgfutil@tmpl{\pgfutil@tmpl -- (\pgfutil@tmpT)}%
\edef\pgfutil@tmple{\pgfutil@tmple\space edge (\pgfutil@tmpT)
(\pgfutil@tmpT)}%
\fi
}%
\ifnum\c@pgf@counta<3\relax
\PackageWarning{3dtools}{A face needs at least three vertices. However, only \the\c@pgf@counta\space vertices were specified.}%
\else
\edef\pgfutil@tmpO{\pgfkeysvalueof{/tikz/3d/polyhedron/O}}%
\pgfmathsetmacro\pgfutil@tmpO{TD("\pgfutil@tmpO")}%
\edef\pgfutil@tmpL{\pgfkeysvalueof{/tikz/3d/polyhedron/L}}%
\pgfmathsetmacro\pgfutil@tmpL{TD("\pgfutil@tmpL")}%
\pgfmathsetmacro{\pgfutil@tmpb}{TD("(\pgfutil@tmpB)-(\pgfutil@tmpA)x(\pgfutil@tmpC)-(\pgfutil@tmpA)")}% 
\pgfmathsetmacro{\pgfutil@tmpc}{TD("(\pgfutil@tmpA)-(\pgfutil@tmpO)")}% 
\pgfmathtruncatemacro{\pgfutil@tmpd}{sign(TD("(\pgfutil@tmpb)o(\pgfutil@tmpc)"))}%
\ifnum\pgfutil@tmpd=-1\relax
\pgfmathsetmacro{\pgfutil@tmpb}{TD("(\pgfutil@tmpB)-(\pgfutil@tmpA)x(\pgfutil@tmpA)-(\pgfutil@tmpC)")}% 
\fi
\pgfmathsetmacro{\pgfutil@tmpe}{screendepth(\pgfutil@tmpb)}%
\pgfmathsetmacro{\pgfutil@tmpf}{sqrt(TD("(\pgfutil@tmpb)o(\pgfutil@tmpb)"))}%
\pgfmathsetmacro{\pgfutil@tmpg}{sqrt(TD("(\pgfutil@tmpL)o(\pgfutil@tmpL)"))}%
\pgfmathsetmacro{\pgfutil@tmph}{TD("(\pgfutil@tmpL)o(\pgfutil@tmpb)")/\pgfutil@tmpf/\pgfutil@tmpg}%
\pgfmathtruncatemacro{\pgfutil@tmpi}{\pgfkeysvalueof{/tikz/3d/polyhedron/shading function}(\pgfutil@tmph)}%
\ifnum\fpeval{\pgfutil@tmpe<0}=1\relax
\begin{pgfonlayer}{\pgfkeysvalueof{/tikz/3d/polyhedron/back layer}}
\tikzset{3d/polyhedron/before hidden,3d/polyhedron/back}%
\if@tikz@td@complete@dashes
\path[fill=tikz@td@face@color!\pgfutil@tmpi!black,3d/polyhedron/back,draw=none] 
	\pgfutil@tmpl -- cycle;
\edef\pgfutil@tmplg{\pgfutil@tmple\space edge \pgfutil@tmplf}%
\draw[3d/polyhedron/back,every edge/.append style={/tikz/3d/polyhedron/face edges}]
	\pgfutil@tmplg;
\else
\draw[fill=tikz@td@face@color!\pgfutil@tmpi!black,3d/polyhedron/back] \pgfutil@tmpl -- cycle;
\fi
\tikzset{3d/polyhedron/after hidden}%
\end{pgfonlayer}
\else
\begin{pgfonlayer}{\pgfkeysvalueof{/tikz/3d/polyhedron/fore layer}}
\tikzset{3d/polyhedron/before visible,3d/polyhedron/fore}%
\if@tikz@td@complete@dashes
\path[fill=tikz@td@face@color!\pgfutil@tmpi!black,3d/polyhedron/fore,draw=none]
 \pgfutil@tmpl -- cycle;
\edef\pgfutil@tmplg{\pgfutil@tmple\space edge \pgfutil@tmplf}%
\draw[3d/polyhedron/fore,every edge/.append style={/tikz/3d/polyhedron/face edges}]
	\pgfutil@tmplg;
\else
\draw[fill=tikz@td@face@color!\pgfutil@tmpi!black,3d/polyhedron/fore] \pgfutil@tmpl -- cycle;
\fi
\tikzset{3d/polyhedron/after visible}%
\end{pgfonlayer}
\fi
\fi
\endgroup
},
draw open face with corners/.code={\begingroup\c@pgf@counta0\relax
\pgfutil@for\pgfutil@tmpa:={#1}\do{%
\pgfmathsetmacro\pgfutil@tmpT{TD("\pgfutil@tmpa")}%
\ifcase\c@pgf@counta 
\pgfmathsetmacro\pgfutil@tmpA{TD("\pgfutil@tmpa")}%
\or
\pgfmathsetmacro\pgfutil@tmpB{TD("\pgfutil@tmpa")}%
\or
\pgfmathsetmacro\pgfutil@tmpC{TD("\pgfutil@tmpa")}%
\fi
\advance\c@pgf@counta by1\relax
\ifnum\c@pgf@counta=1\relax
\edef\pgfutil@tmpl{(\pgfutil@tmpT)}% regular list
\edef\pgfutil@tmplf{(\pgfutil@tmpT)}% first
\edef\pgfutil@tmple{(\pgfutil@tmpT)}% edge list
\else
\edef\pgfutil@tmpl{\pgfutil@tmpl -- (\pgfutil@tmpT)}%
\edef\pgfutil@tmple{\pgfutil@tmple\space edge (\pgfutil@tmpT)
(\pgfutil@tmpT)}%
\fi
}%
\ifnum\c@pgf@counta<3\relax
\PackageWarning{3dtools}{A face needs at least three vertices. However, only \the\c@pgf@counta\space vertices were specified.}%
\else
\edef\pgfutil@tmpO{\pgfkeysvalueof{/tikz/3d/polyhedron/O}}%
\pgfmathsetmacro\pgfutil@tmpO{TD("\pgfutil@tmpO")}%
\edef\pgfutil@tmpL{\pgfkeysvalueof{/tikz/3d/polyhedron/L}}%
\pgfmathsetmacro\pgfutil@tmpL{TD("\pgfutil@tmpL")}%
\pgfmathsetmacro{\pgfutil@tmpb}{TD("(\pgfutil@tmpB)-(\pgfutil@tmpA)x(\pgfutil@tmpC)-(\pgfutil@tmpA)")}% 
\pgfmathsetmacro{\pgfutil@tmpc}{TD("(\pgfutil@tmpA)-(\pgfutil@tmpO)")}% 
\pgfmathtruncatemacro{\pgfutil@tmpd}{sign(TD("(\pgfutil@tmpb)o(\pgfutil@tmpc)"))}%
\ifnum\pgfutil@tmpd=-1\relax
\pgfmathsetmacro{\pgfutil@tmpb}{TD("(\pgfutil@tmpB)-(\pgfutil@tmpA)x(\pgfutil@tmpA)-(\pgfutil@tmpC)")}% 
\fi
\pgfmathsetmacro{\pgfutil@tmpe}{screendepth(\pgfutil@tmpb)}%
\pgfmathsetmacro{\pgfutil@tmpf}{sqrt(TD("(\pgfutil@tmpb)o(\pgfutil@tmpb)"))}%
\pgfmathsetmacro{\pgfutil@tmpg}{sqrt(TD("(\pgfutil@tmpL)o(\pgfutil@tmpL)"))}%
\pgfmathsetmacro{\pgfutil@tmph}{TD("(\pgfutil@tmpL)o(\pgfutil@tmpb)")/\pgfutil@tmpf/\pgfutil@tmpg}%
\pgfmathtruncatemacro{\pgfutil@tmpi}{\pgfkeysvalueof{/tikz/3d/polyhedron/shading function}(\pgfutil@tmph)}%
\ifnum\fpeval{\pgfutil@tmpe<0}=1\relax
\begin{pgfonlayer}{\pgfkeysvalueof{/tikz/3d/polyhedron/back layer}}
\tikzset{3d/polyhedron/before hidden,3d/polyhedron/back}%
\if@tikz@td@complete@dashes
\path[fill=tikz@td@face@color!\pgfutil@tmpi!black,3d/polyhedron/back,draw=none]
	\pgfutil@tmpl;
\let\pgfutil@tmplg\pgfutil@tmple
\draw[3d/polyhedron/back,every edge/.append style={/tikz/3d/polyhedron/face edges}]
	\pgfutil@tmplg;
\else
\draw[fill=tikz@td@face@color!\pgfutil@tmpi!black,3d/polyhedron/back] \pgfutil@tmpl ;
\fi
\tikzset{3d/polyhedron/after hidden}%
\end{pgfonlayer}
\else
\begin{pgfonlayer}{\pgfkeysvalueof{/tikz/3d/polyhedron/fore layer}}
\tikzset{3d/polyhedron/before visible,3d/polyhedron/fore}%
\if@tikz@td@complete@dashes
\path[fill=tikz@td@face@color!\pgfutil@tmpi!black,3d/polyhedron/fore,draw=none]
 \pgfutil@tmpl ;
\let\pgfutil@tmplg\pgfutil@tmple
\draw[3d/polyhedron/fore,every edge/.append style={/tikz/3d/polyhedron/face edges}]
	\pgfutil@tmplg;
\else
\draw[fill=tikz@td@face@color!\pgfutil@tmpi!black,3d/polyhedron/fore] \pgfutil@tmpl ;
\fi
\tikzset{3d/polyhedron/after visible}%
\end{pgfonlayer}
\fi
\fi
\endgroup
},
color/.code={\colorlet{tikz@td@face@color}{#1}},color=yellow,
O/.initial={(0,0,0)},% point inside the polyhedron
L/.initial={\pgfkeysvalueof{/tikz/3d/light}},% light direction, defaults to 3d/light
fore/.style={draw,solid},fore layer/.initial=main,
back/.style={draw,dashed,fill opacity=0},back layer/.initial=main,
shading function/.initial={tikztdpolyhedronshade},
/tikz/declare function={tikztdpolyhedronshade(\x)=100*tikztdlight(\x);},
before visible/.code={},after visible/.code={},
before hidden/.code={},after hidden/.code={},
create face from vertices/.code={\let\pgfutil@tmpa\relax%
\pgfutil@for\pgfutil@tmpb:={#1}\do{%
\ifx\pgfutil@tmpa\relax
\edef\pgfutil@tmpa{{(\pgfkeysvalueof{/tikz/name prefix}\pgfutil@tmpb)}}%
\else
\edef\pgfutil@tmpa{\pgfutil@tmpa,{(\pgfkeysvalueof{/tikz/name prefix}\pgfutil@tmpb)}}%
\fi
}%
\tikzset{3d/polyhedron/draw face with corners/.expanded=\pgfutil@tmpa}%
},%
create faces from vertex list/.style={/tikz/3d/polyhedron/create face from vertices/.list={#1}}%
}%
% draw named paths in some order
\tikzset{3d/draw named path behind/.code 2 args={% #1 object, #2 list of objects in front
   %\typeout{Draw #1 behind #2.}
   \tikzset{3d/current named path=#1,
   	3d/draw named path clipped against/.list={#2},
	3d/draw named path avoiding={#2}}
  },3d/draw named path clipped against/.code={% #1=clip path
  %\typeout{Clip \pgfkeysvalueof{/tikz/3d/current path} against #1.}
  \begin{scope}
   	\clip[use named path/.expanded=#1];
	\draw[3d/hidden,3d/ordered paths/\pgfkeysvalueof{/tikz/3d/current named path}/.try,
		use named path/.expanded=\pgfkeysvalueof{/tikz/3d/current named path}];
  \end{scope}},
  3d/draw named path avoiding/.code={% #1=list of paths to be avoided
  %\typeout{Draw \pgfkeysvalueof{/tikz/3d/current path} avoiding #1.}
  \begin{scope}
   \tikzset{3d/avoid named path/.list={#1}}%
   \draw[3d/visible,3d/ordered paths/\pgfkeysvalueof{/tikz/3d/current named path}/.try,
   	use named path/.expanded=\pgfkeysvalueof{/tikz/3d/current named path}];
  \end{scope}},3d/avoid named path/.code={%
  	\clip[even odd clip,use named path/.expanded=#1,generous outside path];
  },
  3d/current named path/.initial={},
  3d/draw ordered paths/.code={%
  \tikzset{3d/temporary list 1/.initial={#1},
  	3d/draw ordered path/.list={#1}}%  
  },3d/draw ordered path/.code={%
  % test if remaining list has only one element
  \expanded{\noexpand\pgfutil@in@{,}{\pgfkeysvalueof{/tikz/3d/temporary list 1}}}% 
  \tikzset{3d/temporary list 1/.remove item={#1}}
  \ifpgfutil@in@% 
	\tikzset{3d/draw named path behind/.expanded={#1}{\pgfkeysvalueof{/tikz/3d/temporary list 1}}} 
  \else
   \begin{scope}
	\draw[3d/visible,3d/ordered paths/#1/.try,
	use named path/.expanded=#1];
   \end{scope}	
  \fi  
  },3d/ordered paths/.unknown/.code={}}
% y cylinder
\tikzset{pics/ycylinder/.style={code={%
% old version
% \tikzset{3d/.cd,#1}
% \def\pv##1{\pgfkeysvalueof{/tikz/3d/##1}}
% \pgfmathsetmacro{\vmin}{atan2(sin(\pgfkeysvalueof{/tikz/3d/theta})*sin(\pgfkeysvalueof{/tikz/3d/phi}),%
% min(cos(\pgfkeysvalueof{/tikz/3d/phi}),-1/sqrt(2))*cos(\pgfkeysvalueof{/tikz/3d/theta}))}
% \pgfmathsetmacro{\vmax}{\vmin-180}
% \path[3d/cylinder/mantle]
%   let \p1=($(0,1,0)-(0,0,0)$),\n1={atan2(\y1,\x1)+180} in
%   [shading angle=\n1]
%   plot[variable=\t,domain=\vmin:\vmax,smooth]
%     ({\pv{r}*cos(\t)},0,{\pv{r}*sin(\t)})
%     -- 
%     plot[variable=\t,domain=\vmax:\vmin,smooth]
%     ({\pv{r}*cos(\t)},\pv{h},{\pv{r}*sin(\t)})
%     --cycle;
% \pgfmathtruncatemacro{\itest}{sign(cos(\pgfkeysvalueof{/tikz/3d/phi}))}
% \ifnum\itest=-1
%     \path[3d/cylinder/top] plot[variable=\t,domain=0:360,smooth cycle]
%     ({\pv{r}*cos(\t)},\pv{h},{\pv{r}*sin(\t)}) ;
% \fi
% \ifnum\itest=1
%     \path[3d/cylinder/top] plot[variable=\t,domain=0:360,smooth cycle]
%     ({\pv{r}*cos(\t)},0,{\pv{r}*sin(\t)}) ;
% \fi
% new version
\tikzset{3d/.cd,#1}%
%\def\pv##1{\pgfkeysvalueof{/tikz/3d/##1}}
\pgfmathsetmacro{\tikz@td@vmin}{-atan2(sin(\pgfkeysvalueof{/tikz/3d/phi})*sin(\pgfkeysvalueof{/tikz/3d/theta}),cos(\pgfkeysvalueof{/tikz/3d/theta}))}%
\pgfmathsetmacro{\tikz@td@vmax}{\tikz@td@vmin+180}%
\def\tikz@td@mantle{plot[variable=\tikz@td@t,domain=\tikz@td@vmin:\tikz@td@vmax,smooth]
    ({\pgfkeysvalueof{/tikz/3d/r}*cos(\tikz@td@t)},0,{\pgfkeysvalueof{/tikz/3d/r}*sin(\tikz@td@t)})
    -- 
    plot[variable=\tikz@td@t,domain=\tikz@td@vmax:\tikz@td@vmin,smooth]
    ({\pgfkeysvalueof{/tikz/3d/r}*cos(\tikz@td@t)},\pgfkeysvalueof{/tikz/3d/h},{\pgfkeysvalueof{/tikz/3d/r}*sin(\tikz@td@t)})
    --cycle}%
\pgfmathsetmacro{\tikz@td@ex}{TDunit("(0,1,0)x(nscreenx,nscreeny,nscreenz)")}%
\tikz@td@cylinder@light{\tikz@td@ex}{0,1,0}%
\path[3d/cylinder/lighting] \tikz@td@mantle;
\path[3d/cylinder/mantle]
  let \p1=($(0,1,0)-(0,0,0)$),\n1={atan2(\y1,\x1)+180} in
  [3d/cylinder/shading angle=\n1]
  \tikz@td@mantle;
\tikz@td@cylinder@cap{0,1,0}%
\pgfmathtruncatemacro{\tikz@td@i}{sign(cos(\pgfkeysvalueof{/tikz/3d/phi}))}%
\ifnum\tikz@td@i=-1\relax
\pgfmathtruncatemacro{\tikz@td@i}{sign(cos(\pgfkeysvalueof{/tikz/3d/phi}))}%
    \def\tikz@td@mantle{plot[variable=\tikz@td@t,domain=0:360,smooth cycle]
    ({\pgfkeysvalueof{/tikz/3d/r}*cos(\tikz@td@t)},\pgfkeysvalueof{/tikz/3d/h},{\pgfkeysvalueof{/tikz/3d/r}*sin(\tikz@td@t)})}%
    \tikz@td@cylinder@paintcap
    \path[3d/cylinder/top] \tikz@td@mantle;
\fi
\ifnum\tikz@td@i=1\relax
    \def\tikz@td@mantle{plot[variable=\tikz@td@t,domain=0:360,smooth cycle]
    ({\pgfkeysvalueof{/tikz/3d/r}*cos(\tikz@td@t)},0,{\pgfkeysvalueof{/tikz/3d/r}*sin(\tikz@td@t)})}%
    \tikz@td@cylinder@paintcap
    \path[3d/cylinder/top] \tikz@td@mantle;
\fi
}},pics/cylinder/.style={code={%
% arbitrary cylinder
\tikzset{3d/.cd,#1}%
\pgfmathsetmacro{\tikz@td@v}{TD("\pgfkeysvalueof{/tikz/3d/cylinder/axis}x(nscreenx,nscreeny,nscreenz)")}%
\pgfmathtruncatemacro{\tikz@td@i}{abs(TD("(\tikz@td@v)o(\tikz@td@v)"))>0.01?1:0}%
\pgfmathsetmacro{\tikz@td@ez}{TDunit("\pgfkeysvalueof{/tikz/3d/cylinder/axis}")}%
\tikz@td@cylinder@cap{\tikz@td@ez}%
\ifnum\tikz@td@i=0\relax % seen along the axis: only the cap facing the viewer
  % with the hole and open ends as in the general view
  \edef\tikz@td@mantle{(0,0) circle[radius=\pgfkeysvalueof{/tikz/3d/r}]
    \ifdim\pgfkeysvalueof{/tikz/3d/cylinder/hole}pt>0pt
      (0,0) circle[radius=\pgfkeysvalueof{/tikz/3d/cylinder/hole}]\fi}%
  \ifcsname tikz@td@open@\pgfkeysvalueof{/tikz/3d/cylinder/open}\endcsname
    \path[3d/cylinder/top,3d/edge only] \tikz@td@mantle;
  \else
    \tikz@td@iffills{/tikz/3d/cylinder/top}{}{\tikz@td@cylinder@paintcap}%
    \path[even odd rule,3d/cylinder/top] \tikz@td@mantle;
  \fi
\else 
  \pgfmathsetmacro{\tikz@td@ex}{TDunit("\pgfkeysvalueof{/tikz/3d/cylinder/axis}x(nscreenx,nscreeny,nscreenz)")}%
  \pgfmathsetmacro{\tikz@td@ey}{TDunit("(\tikz@td@ex)x\pgfkeysvalueof{/tikz/3d/cylinder/axis}")}%
  % ex x axis points towards the viewer; the frame uses ey = -(ex x axis) so
  % that the mantle (local y<0) is the half facing the viewer from either side
  \pgfmathsetmacro{\tikz@td@ey}{TD("-1*(\tikz@td@ey)")}%
  \pgfmathsetmacro{\tikz@td@ez}{TDunit("\pgfkeysvalueof{/tikz/3d/cylinder/axis}")}%
  \tikz@td@cylinder@light{\tikz@td@ex}{\tikz@td@ez}%
  \expanded{\noexpand\path[overlay,
	3d coordinate={\pgfkeysvalueof{/tikz/3d/aux keys/ex}=(\tikz@td@ex)},
	3d coordinate={\pgfkeysvalueof{/tikz/3d/aux keys/ey}=(\tikz@td@ey)},
	3d coordinate={\pgfkeysvalueof{/tikz/3d/aux keys/ez}=(\tikz@td@ez)}];}%
  \pgfmathtruncatemacro{\tikz@td@i}{TD("\pgfkeysvalueof{/tikz/3d/cylinder/axis}o(nscreenx,nscreeny,nscreenz)")>0?1:0}%
  \begin{scope}[x/.expanded={\pgfkeysvalueof{/tikz/3d/aux keys/ex}},
    y/.expanded={\pgfkeysvalueof{/tikz/3d/aux keys/ey}},
    z/.expanded={\pgfkeysvalueof{/tikz/3d/aux keys/ez}}]       
    % pieces: the rim behind, the mantle (centroid of its front half), the cap;
    % the radius breaks ties with coplanar caps
    \pgfmathsetmacro\tikz@td@kext{abs(\pgfkeysvalueof{/tikz/3d/r})}%
    \ifnum\tikz@td@i=1\relax 
      \tikz@td@piece{(0,{2*abs(\pgfkeysvalueof{/tikz/3d/r})/pi},0)}{%
      \draw[3d/hidden,3d/edge only] (\pgfkeysvalueof{/tikz/3d/r},0,0) arc[start angle=0,end angle=180,radius=\pgfkeysvalueof{/tikz/3d/r}];}%
    \else 
      \tikz@td@piece{(0,{2*abs(\pgfkeysvalueof{/tikz/3d/r})/pi},\pgfkeysvalueof{/tikz/3d/h})}{%
      \draw[3d/hidden,3d/edge only] (\pgfkeysvalueof{/tikz/3d/r},0,\pgfkeysvalueof{/tikz/3d/h}) arc[start angle=0,end angle=180,radius=\pgfkeysvalueof{/tikz/3d/r}];}%
    \fi
    \def\tikz@td@mantle{(\pgfkeysvalueof{/tikz/3d/r},0,0) arc[start angle=0,end angle=-180,radius=\pgfkeysvalueof{/tikz/3d/r}]
        -- 
      ({-1*(\pgfkeysvalueof{/tikz/3d/r})},0,\pgfkeysvalueof{/tikz/3d/h}) arc[start angle=-180,end angle=0,radius=\pgfkeysvalueof{/tikz/3d/r}]
      -- cycle}%
    \tikz@td@piece{(0,{-2*abs(\pgfkeysvalueof{/tikz/3d/r})/pi},{\pgfkeysvalueof{/tikz/3d/h}/2})}{%
    \tikz@td@iffills{/tikz/3d/cylinder/mantle}{}{\path[3d/cylinder/lighting] \tikz@td@mantle;}%
    % the outline turns back at the top corners: no miter spikes there
    \path[line join=round,3d/cylinder/mantle]
      let \p1=($(0,1,0)-(0,0,0)$),\n1={atan2(\y1,\x1)+180} in
      [3d/cylinder/shading angle=\n1]
      \tikz@td@mantle;}%
    \ifnum\tikz@td@i=1\relax 
      \def\pgfutil@tmpa{(0,0,\pgfkeysvalueof{/tikz/3d/h})}%
    \else 
      \def\pgfutil@tmpa{(0,0,0)}%
    \fi
    % the cap, with a hole of radius 3d/cylinder/hole (even odd rule)
    \edef\tikz@td@mantle{\pgfutil@tmpa\space circle[radius=\pgfkeysvalueof{/tikz/3d/r}]
      \ifdim\pgfkeysvalueof{/tikz/3d/cylinder/hole}pt>0pt
        \pgfutil@tmpa\space circle[radius=\pgfkeysvalueof{/tikz/3d/cylinder/hole}]\fi}%
    \expandafter\tikz@td@piece\expandafter{\pgfutil@tmpa}{%
    \ifcsname tikz@td@open@\pgfkeysvalueof{/tikz/3d/cylinder/open}\endcsname
      % open ends: only the rim
      \path[3d/cylinder/top,3d/edge only] \tikz@td@mantle;
    \else
      \tikz@td@iffills{/tikz/3d/cylinder/top}{}{\tikz@td@cylinder@paintcap}%
      \path[even odd rule,3d/cylinder/top] \tikz@td@mantle;
    \fi}%
  \end{scope}  
\fi
}},%
3d/.cd,cylinder/.cd,
mantle/.style={draw},top/.style={draw}}
\tikzset{3d/.cd,visible/.style={draw,solid},
	hidden/.style={draw,very thin,cheating dash},
	r/.initial=1,h/.initial=1,R/.initial=2,
  cylinder/axis/.initial={(0,0,1)},cylinder/hole/.initial=0,cylinder/open/.initial=false}%
\expandafter\let\csname tikz@td@open@true\endcsname\relax
%%
%% decorations
%% 
%% path restoration
\newcounter{tikz@td@recorded@point}
\pgfmathdeclarefunction{tikztdindexoflastcoordinate}{0}{%
\edef\pgfmathresult{\number\value{tikz@td@recorded@point}}%
}%

\tikzset{stored path/.cd,coordinate prefix/.initial={stored-},
	reset coordinate index/.code={\setcounter{tikz@td@recorded@point}{0}},
	restore path/.style={insert path/.expanded={\csname tikz@td@recorded@path@f@#1\endcsname}},%	
	restore reversed path/.style={insert path/.expanded={\csname tikz@td@recorded@path@r@#1\endcsname}},%	
	append path/.style={insert path/.expanded={-- \csname tikz@td@recorded@path@f@#1\endcsname}},%	
	append reversed path/.style={insert path/.expanded={-- \csname tikz@td@recorded@path@r@#1\endcsname}},%	
	first coordinate of/.style={insert path/.expanded={(\pgfkeysvalueof{/tikz/stored path/coordinate prefix}#1-0)}},%	
	last coordinate of/.style={insert path/.expanded={(\pgfkeysvalueof{/tikz/stored path/coordinate prefix}#1-\csname tikz@td@recorded@path@last@#1\endcsname)}},%
	coordinate with index/.style args={#1 of #2}{insert path/.expanded={(\pgfkeysvalueof{/tikz/stored path/coordinate prefix}#2-#1)}},%	
/tikz/extend stored path/.style={%
	%/utils/exec=\typeout{\arabic{tikz@td@recorded@point}},
    decorate,decoration={show path construction,
    moveto code={\def\tikz@td@prefix{\pgfkeysvalueof{/tikz/stored path/coordinate prefix}#1-}%
	% do not overwrite the last coordinate of the previous subpath
	\ifnum\value{tikz@td@recorded@point}>0 \stepcounter{tikz@td@recorded@point}\fi
	\path (\tikzinputsegmentfirst) 
		coordinate (\tikz@td@prefix\the\numexpr\value{tikz@td@recorded@point});
    \ifnum\value{tikz@td@recorded@point}=0
	 \expandafter\xdef\csname tikz@td@recorded@path@f@#1\endcsname{% forward
	 	(\tikz@td@prefix\the\numexpr\value{tikz@td@recorded@point})}%
	 \expandafter\xdef\csname tikz@td@recorded@path@r@#1\endcsname{% reverse
	 	(\tikz@td@prefix\the\numexpr\value{tikz@td@recorded@point})}%
	\else
	 \expandafter\xdef\csname tikz@td@recorded@path@f@#1\endcsname{% forward
	 	\csname tikz@td@recorded@path@f@#1\endcsname
		(\tikz@td@prefix\the\numexpr\value{tikz@td@recorded@point})}%
	 \expandafter\xdef\csname tikz@td@recorded@path@r@#1\endcsname{% reverse
	 	(\tikz@td@prefix\the\numexpr\value{tikz@td@recorded@point}) 
		\csname tikz@td@recorded@path@r@#1\endcsname}%
		\fi
	\expandafter\xdef\csname tikz@td@recorded@path@last@#1\endcsname{\number\value{tikz@td@recorded@point}}%
	\stepcounter{tikz@td@recorded@point}%
    },
    lineto code={\def\tikz@td@prefix{\pgfkeysvalueof{/tikz/stored path/coordinate prefix}#1-}%
	\path (\tikzinputsegmentfirst) 
		coordinate (\tikz@td@prefix\the\numexpr\value{tikz@td@recorded@point})
		 (\tikzinputsegmentlast) 
		coordinate (\tikz@td@prefix\the\numexpr\value{tikz@td@recorded@point}+1);
    \ifnum\value{tikz@td@recorded@point}=0
	 \expandafter\xdef\csname tikz@td@recorded@path@f@#1\endcsname{% forward
	 	(\tikz@td@prefix\the\numexpr\value{tikz@td@recorded@point})
		--(\tikz@td@prefix\the\numexpr\value{tikz@td@recorded@point}+1)}%
	 \expandafter\xdef\csname tikz@td@recorded@path@r@#1\endcsname{% reverse
	 	(\tikz@td@prefix\the\numexpr\value{tikz@td@recorded@point}+1)
		--(\tikz@td@prefix\the\numexpr\value{tikz@td@recorded@point})}%
	\else
	 \expandafter\xdef\csname tikz@td@recorded@path@f@#1\endcsname{% forward
	 	\csname tikz@td@recorded@path@f@#1\endcsname
		--(\tikz@td@prefix\the\numexpr\value{tikz@td@recorded@point}+1)}%
	 \expandafter\xdef\csname tikz@td@recorded@path@r@#1\endcsname{% reverse
	 	(\tikz@td@prefix\the\numexpr\value{tikz@td@recorded@point}+1) --
		\csname tikz@td@recorded@path@r@#1\endcsname}%
		\fi
	\stepcounter{tikz@td@recorded@point}%
	\expandafter\xdef\csname tikz@td@recorded@path@last@#1\endcsname{\number\value{tikz@td@recorded@point}}%
    },
    curveto code={\def\tikz@td@prefix{\pgfkeysvalueof{/tikz/stored path/coordinate prefix}#1-}%
      \path (\tikzinputsegmentfirst)
	    coordinate (\tikz@td@prefix\the\numexpr\value{tikz@td@recorded@point})
        (\tikzinputsegmentsupporta) 
		coordinate (\tikz@td@prefix\the\numexpr\value{tikz@td@recorded@point}+1)
		(\tikzinputsegmentsupportb)
		coordinate (\tikz@td@prefix\the\numexpr\value{tikz@td@recorded@point}+2)
        (\tikzinputsegmentlast)
		coordinate (\tikz@td@prefix\the\numexpr\value{tikz@td@recorded@point}+3);
    \ifnum\value{tikz@td@recorded@point}=0
	 \expandafter\xdef\csname tikz@td@recorded@path@f@#1\endcsname{% forward
	 	(\tikz@td@prefix\the\numexpr\value{tikz@td@recorded@point})
		.. controls (\tikz@td@prefix\the\numexpr\value{tikz@td@recorded@point}+1)
		and (\tikz@td@prefix\the\numexpr\value{tikz@td@recorded@point}+2)
		.. (\tikz@td@prefix\the\numexpr\value{tikz@td@recorded@point}+3)}%
	 \expandafter\xdef\csname tikz@td@recorded@path@r@#1\endcsname{% reverse
	 	(\tikz@td@prefix\the\numexpr\value{tikz@td@recorded@point}+3)
		.. controls (\tikz@td@prefix\the\numexpr\value{tikz@td@recorded@point}+2)
		and (\tikz@td@prefix\the\numexpr\value{tikz@td@recorded@point}+1)
		.. (\tikz@td@prefix\the\numexpr\value{tikz@td@recorded@point})}%
	\else
	 \expandafter\xdef\csname tikz@td@recorded@path@f@#1\endcsname{% forward
	 	\csname tikz@td@recorded@path@f@#1\endcsname
		.. controls (\tikz@td@prefix\the\numexpr\value{tikz@td@recorded@point}+1)
		and (\tikz@td@prefix\the\numexpr\value{tikz@td@recorded@point}+2)
		.. (\tikz@td@prefix\the\numexpr\value{tikz@td@recorded@point}+3)}%
	 \expandafter\xdef\csname tikz@td@recorded@path@r@#1\endcsname{% reverse
	 	(\tikz@td@prefix\the\numexpr\value{tikz@td@recorded@point}+3)
		.. controls (\tikz@td@prefix\the\numexpr\value{tikz@td@recorded@point}+2)
		and (\tikz@td@prefix\the\numexpr\value{tikz@td@recorded@point}+1) ..
		\csname tikz@td@recorded@path@r@#1\endcsname}%
	\fi
	\stepcounter{tikz@td@recorded@point}%
	\stepcounter{tikz@td@recorded@point}%
	\stepcounter{tikz@td@recorded@point}%
	\expandafter\xdef\csname tikz@td@recorded@path@last@#1\endcsname{\number\value{tikz@td@recorded@point}}%
    },
    closepath code={\def\tikz@td@prefix{\pgfkeysvalueof{/tikz/stored path/coordinate prefix}#1-}%
	\path (\tikzinputsegmentfirst) 
		coordinate (\tikz@td@prefix\the\numexpr\value{tikz@td@recorded@point})
		 (\tikzinputsegmentlast) 
		coordinate (\tikz@td@prefix\the\numexpr\value{tikz@td@recorded@point}+1);
    \ifnum\value{tikz@td@recorded@point}=0
	 \expandafter\xdef\csname tikz@td@recorded@path@f@#1\endcsname{% forward
	 	(\tikz@td@prefix\the\numexpr\value{tikz@td@recorded@point})
		--(\tikz@td@prefix\the\numexpr\value{tikz@td@recorded@point}+1)}%
	 \expandafter\xdef\csname tikz@td@recorded@path@r@#1\endcsname{% reverse
	 	(\tikz@td@prefix\the\numexpr\value{tikz@td@recorded@point}+1)
		--(\tikz@td@prefix\the\numexpr\value{tikz@td@recorded@point})}%
	\else
	 \expandafter\xdef\csname tikz@td@recorded@path@f@#1\endcsname{% forward
	 	\csname tikz@td@recorded@path@f@#1\endcsname
		--(\tikz@td@prefix\the\numexpr\value{tikz@td@recorded@point}+1)}%
	 \expandafter\xdef\csname tikz@td@recorded@path@r@#1\endcsname{% reverse
		 	(\tikz@td@prefix\the\numexpr\value{tikz@td@recorded@point}+1) --
		\csname tikz@td@recorded@path@r@#1\endcsname}%
	\fi
	\stepcounter{tikz@td@recorded@point}%
	\expandafter\xdef\csname tikz@td@recorded@path@last@#1\endcsname{\number\value{tikz@td@recorded@point}}%
	  } }},
/tikz/store path/.style={%	  
	  /tikz/stored path/reset coordinate index,%
	  /tikz/extend stored path=#1}
	  }
%% bounding box for animations
\def\tikz@td@bbox@store@aux#1#2{%
	\immediate\write\@mainaux{\string\expandafter\xdef\noexpand\csname pgfk@/tikz/3d/bbox/stored/#1\string\endcsname{#2}}%
	\expandafter\xdef\csname pgfk@/tikz/3d/bbox/stored/#1\endcsname{#2}}%
\def\tikz@td@bbox@@get@min@from@aux#1#2{%
	\ifcsname pgfk@/tikz/3d/bbox/stored/#1min\endcsname
	\edef#2{\csname pgfk@/tikz/3d/bbox/stored/#1min\endcsname}%
	\else
	\edef#2{16383.99999pt}%
	\expandafter\edef\csname pgfk@/tikz/3d/bbox/stored/#1min\endcsname{16383.99999pt}%
	\fi}
\def\tikz@td@bbox@@get@max@from@aux#1#2{%
	\ifcsname pgfk@/tikz/3d/bbox/stored/#1max\endcsname
	\edef#2{\csname pgfk@/tikz/3d/bbox/stored/#1max\endcsname}%
	\else
	\edef#2{-16383.99999pt}%
	\expandafter\edef\csname pgfk@/tikz/3d/bbox/stored/#1max\endcsname{-16383.99999pt}%
	\fi}
\tikzset{same bounding box/.style={%
execute at end picture={%
\tikz@td@bbox@@get@min@from@aux{#1x}\pgfutil@tmpa%
\tikz@td@bbox@@get@min@from@aux{#1y}\pgfutil@tmpb%
\tikz@td@bbox@@get@max@from@aux{#1x}\pgfutil@tmpc%
\tikz@td@bbox@@get@max@from@aux{#1y}\pgfutil@tmpd%
%\typeout{(\pgfutil@tmpa,\pgfutil@tmpb) (\pgfutil@tmpc,\pgfutil@tmpd)}%
\pgfmathsetmacro\pgfutil@tmpa{min(\pgfutil@tmpa,\pgf@picminx,\pgfkeysvalueof{/tikz/3d/bbox/stored/#1xmin})}%
\pgfmathsetmacro\pgfutil@tmpb{min(\pgfutil@tmpb,\pgf@picminy,\pgfkeysvalueof{/tikz/3d/bbox/stored/#1ymin})}%
\pgfmathsetmacro\pgfutil@tmpc{max(\pgfutil@tmpc,\pgf@picmaxx,\pgfkeysvalueof{/tikz/3d/bbox/stored/#1xmax})}%
\pgfmathsetmacro\pgfutil@tmpd{max(\pgfutil@tmpd,\pgf@picmaxy,\pgfkeysvalueof{/tikz/3d/bbox/stored/#1ymax})}%
%\typeout{(\pgfutil@tmpa,\pgfutil@tmpb) (\pgfutil@tmpc,\pgfutil@tmpd)}%
\tikz@td@bbox@store@aux{#1xmin}{\pgfutil@tmpa}%
\tikz@td@bbox@store@aux{#1ymin}{\pgfutil@tmpb}%
\tikz@td@bbox@store@aux{#1xmax}{\pgfutil@tmpc}%
\tikz@td@bbox@store@aux{#1ymax}{\pgfutil@tmpd}%
\ifdim\pgfutil@tmpa pt<16383.99999pt\relax
\path[reset cm] (\pgfutil@tmpa pt,\pgfutil@tmpb pt) (\pgfutil@tmpc pt,\pgfutil@tmpd pt);
\fi
}},same bounding box/.default=A} 	  
%% list management
\newif\if@pgf@list@first@item
\pgfkeys{/handlers/.add item/.code={\edef\pgfutil@tmpa{\pgfkeysvalueof{\pgfkeyscurrentpath}}%
\ifx\pgfutil@tmpa\pgfutil@empty\pgfkeyssetvalue{\pgfkeyscurrentpath}{#1}%
\else\pgfkeys{\pgfkeyscurrentpath/.append={,#1}}\fi},
/handlers/.remove item/.code={%
\edef\pgfutil@tmpa{\pgfkeysvalueof{\pgfkeyscurrentpath}}%
\edef\pgfutil@tmpb{#1}%
\edef\pgfutil@tmpc{\noexpand\@removeelement{\pgfutil@tmpb}{\pgfutil@tmpa}{\noexpand\pgfutil@tmpa}}%
\pgfutil@tmpc
\pgfkeyslet{\pgfkeyscurrentpath}{\pgfutil@tmpa}},% maybe test if the list is now empty?
/handlers/.count/.code={\begingroup%
\edef\pgfutil@tmpl{\pgfkeysvalueof{\pgfkeyscurrentpath}}%
\c@pgf@counta0\relax
\pgfkeys{/list management/step count/.list/.expanded={\pgfutil@tmpl}}%
\edef#1{\the\c@pgf@counta}%
\pgfmathsmuggle#1\endgroup},%
/handlers/.print list/.code={\begingroup%
\@pgf@list@first@itemtrue
\edef\pgfutil@tmpl{\pgfkeysvalueof{\pgfkeyscurrentpath}}%
\pgfkeys{/list management/print item in list/.list/.expanded={\pgfutil@tmpl}}%
\endgroup},
/handlers/.is array/.code={\pgfkeysifdefined{\pgfkeyscurrentpath}{%
\typeout{The key `\pgfkeyscurrentpath` is already in use, and won't get overwritten.}}{%
\pgfkeysifdefined{\pgfkeyscurrentpath/dim}{%
\typeout{The key `\pgfkeyscurrentpath` is already in use, and won't get overwritten.}}{%
\c@pgf@counta=0\relax
\pgfkeyssetevalue{/list management/current array}{\pgfkeyscurrentpath}%
\pgfkeys{/list management/store array entry/.list/.expanded={#1}}%
\pgfkeys{\pgfkeysvalueof{/list management/current array}/dim/.initial=\the\c@pgf@counta}%
\pgfkeys{\pgfkeysvalueof{/list management/current array}/content/.initial={#1}}%
}}},
/handlers/.sort numeric list/.code 2 args={\tikz@td@sort@numeric#1#2},
/list management/.cd,print item/.code={#1},
print item in list/.code={%
\if@pgf@list@first@item
 \pgfkeys{/list management/print item={#1}}%
 \@pgf@list@first@itemfalse
\else
 \pgfkeys{/list management/print item={%
  \pgfkeysvalueof{/list management/list separator}#1}}%
\fi},
list separator/.initial={,\space},
step count/.code={\advance\c@pgf@counta by1\relax},
sort comparison/.initial={>=},
sort map/.code={\pgfmathparse{#1}},
current array/.initial={},
store array entry/.code={\advance\c@pgf@counta by1\relax
\pgfkeyssetevalue{\pgfkeysvalueof{/list management/current array}/\the\c@pgf@counta}{#1}}
}%
%% coil
\newif\ifcoil@closed%
\pgfkeys{%%
/pgf/decoration/.cd,%
3d coil color/.initial=black,%
3d coil width/.initial=0.4pt,%
3d coil dist/.initial=0.6pt,%
3d coil opacity/.initial=1,%
3d coil closed/.code=\coil@closedtrue%
}%
% https://tex.stackexchange.com/a/219088/121799%
% \tikzset{get stroke color/.code={%%
%     \expandafter\global% Jump over, now we have \global%
%     \expandafter\let% Jump over now we have \global\let%
%     \expandafter\pgfsavedstrokecolor% Jump we have \global\let\pgf...%
%     \csname\string\color@pgfstrokecolor\endcsname% Finally expand this and put it at the end %
% 	\typeout{\pgfsavedstrokecolor}%
% 	\pgfutil@namelet{\string\color@pgfstrokecolor}{pgf@strokecolor@global}%
%     },                                           % \global\let\pgf...{} in expanded form %
%     restore stroke color/.code={\pgf@setstrokecolor#1},%
% }%
\def\pgfpoint@onthreedcoil#1#2#3{%%
  \pgf@x=#1\pgfdecorationsegmentamplitude%%
  \pgf@x=\pgfdecorationsegmentaspect\pgf@x%%
  \pgf@y=#2\pgfdecorationsegmentamplitude%%
  \pgf@xa=0.083333333333\pgfdecorationsegmentlength%%
  \advance\pgf@x by#3\pgf@xa%%
  \advance\pgf@x by-\generaloffset pt%%
}%
% coil decoration%
%%
% Parameters: \pgfdecorationsegmentamplitude, \pgfdecorationsegmentlength,%
\pgfdeclaredecoration{3d complete coil}{initial}%
{ %
    \state{initial}[width=0.5*\pgfdecorationsegmentlength,%
    next state=coil, persistent precomputation={% from https://tex.stackexchange.com/a/25689/121799%
    \pgfmathsetmacro\matchinglength{\pgfdecoratedinputsegmentlength / max(1,int(\pgfdecoratedinputsegmentlength/\pgfdecorationsegmentlength))}% at least one winding
    \setlength{\pgfdecorationsegmentlength}{\matchinglength pt}%
    %\tikzset{get stroke color}%
    \pgfmathsetmacro{\generaloffset}{\pgfdecorationsegmentlength}%
    \pgfmathsetmacro{\initialoffset}{1.5*\pgfdecorationsegmentlength}%
    \pgfmathsetmacro{\auxoffset}{2.5*\pgfdecorationsegmentlength}%
  }]    { %
    % line in the back%
    %%
	\edef\TDCoilColor{\pgfkeysvalueof{/pgf/decoration/3d coil color}}%
	\edef\TDCoilOpacity{\pgfkeysvalueof{/pgf/decoration/3d coil opacity}}%
	\edef\TDCoilWidth{\pgfkeysvalueof{/pgf/decoration/3d coil width}}%
	\edef\TDCoilDist{\pgfkeysvalueof{/pgf/decoration/3d coil dist}}%
    \pgfsetstrokecolor{\pgfkeysvalueof{/pgf/decoration/3d coil color}}%
    \pgfsetfillcolor{\pgfkeysvalueof{/pgf/decoration/3d coil color}}%
    \pgfsetstrokeopacity{\pgfkeysvalueof{/pgf/decoration/3d coil opacity}}%
    \pgfsetlinewidth{\pgfkeysvalueof{/pgf/decoration/3d coil width}}     %
    \ifcoil@closed%
     \begingroup%
      \def\generaloffset{\auxoffset}%
      \pgfpathmoveto{\pgfpoint@onthreedcoil{1    }{ 1    }{15}}%
      \pgfpathcurveto%
      {\pgfpoint@onthreedcoil{1.555}{ 1    }{16}}%
      {\pgfpoint@onthreedcoil{2    }{ 0.555}{17}}%
      {\pgfpoint@onthreedcoil{2    }{ 0    }{18}}%
      \pgfcoordinate{TD@coilast}{\pgfpoint@onthreedcoil{2    }{ 0    }{18}}%
      \pgfcoordinate{TD@coilfirst}{\pgfpoint@onthreedcoil{1    }{ 1    }{15}}%
      \pgfusepath{stroke} %
      \pgfsetstrokecolor{\pgfkeysvalueof{/pgf/decoration/3d coil color}}%
     \endgroup%
    \fi%
    \begingroup %%
    \def\generaloffset{\initialoffset}%
    \ifcoil@closed%
     \pgfpathmoveto{\pgfpointanchor{TD@coilast}{center}}%
    \else%
     \pgfpathmoveto{\pgfpointorigin}%
    \fi%
    \pgfpathcurveto%
    {\pgfpoint@onthreedcoil{2    }{-0.555}{7}}%
    {\pgfpoint@onthreedcoil{1.555}{-1    }{8}}%
    {\pgfpoint@onthreedcoil{1    }{-1    }{9}}%
    \pgfusepath{stroke} %
    %%
    % white background for front thick part%
    %%
    \pgfsetstrokeopacity{1}%
    \pgfsetstrokecolor{white}%
    \pgfsetfillcolor{white}%
    \pgfsetlinewidth{1.5*\pgfkeysvalueof{/pgf/decoration/3d coil width}+1.5*\pgfkeysvalueof{/pgf/decoration/3d coil dist}}%
    \pgfpathmoveto{\pgfpoint@onthreedcoil{1    }{-1    }{9}}%
    % draw forward%
    \pgfpathcurveto%
    {\pgfpoint@onthreedcoil{0.445}{-1    }{10}}%
    {\pgfpoint@onthreedcoil{0    }{-0.555}{11.25}}%
    {\pgfpoint@onthreedcoil{0    }{ 0    }{12.5}}%
    \pgfpathcurveto%
    {\pgfpoint@onthreedcoil{0    }{ 0.555}{13.25}}%
    {\pgfpoint@onthreedcoil{0.445}{ 1    }{14.25}}%
    {\pgfpoint@onthreedcoil{1    }{ 1    }{15}}%
    % draw the curve back%
    \pgfpathcurveto%
    {\pgfpoint@onthreedcoil{0.445}{ 1    }{14}}%
    {\pgfpoint@onthreedcoil{0    }{ 0.555}{12.75}}%
    {\pgfpoint@onthreedcoil{0    }{ 0    }{11.5}}%
    \pgfpathcurveto%
    {\pgfpoint@onthreedcoil{0    }{-0.555}{10.75}}%
    {\pgfpoint@onthreedcoil{0.445}{-1    }{10}}%
    {\pgfpoint@onthreedcoil{1    }{-1    }{9}}%
    \pgfusepath{stroke,fill} %
    % %
    % draw the thick foreground path%
    %%
    \pgfsetstrokecolor{\pgfkeysvalueof{/pgf/decoration/3d coil color}}%
    \pgfsetfillcolor{\pgfkeysvalueof{/pgf/decoration/3d coil color}}%
    \pgfsetstrokeopacity{\pgfkeysvalueof{/pgf/decoration/3d coil opacity}}%
    \pgfpathmoveto{\pgfpoint@onthreedcoil{1    }{ 1    }{3}}%
    \pgfsetlinewidth{\pgfkeysvalueof{/pgf/decoration/3d coil width}} %
    % forward shifted +%
    \pgfpathmoveto{\pgfpoint@onthreedcoil{1    }{-1    }{9}}%
    \pgfpathcurveto%
    {\pgfpoint@onthreedcoil{0.445}{-1    }{10}}%
    {\pgfpoint@onthreedcoil{0    }{-0.555}{11.25}}%
    {\pgfpoint@onthreedcoil{0    }{ 0    }{12.5}}%
    \pgfpathcurveto%
    {\pgfpoint@onthreedcoil{0    }{ 0.555}{13.25}}%
    {\pgfpoint@onthreedcoil{0.445}{ 1    }{14.25}}%
    {\pgfpoint@onthreedcoil{1    }{ 1    }{15}}%
    % draw the curve back shfted -%
    \pgfpathcurveto%
    {\pgfpoint@onthreedcoil{0.445}{ 1    }{14}}%
    {\pgfpoint@onthreedcoil{0    }{ 0.555}{12.75}}%
    {\pgfpoint@onthreedcoil{0    }{ 0    }{11.5}}%
    \pgfpathcurveto%
    {\pgfpoint@onthreedcoil{0    }{-0.555}{10.75}}%
    {\pgfpoint@onthreedcoil{0.445}{-1    }{10}}%
    {\pgfpoint@onthreedcoil{1    }{-1    }{9}}%
    \pgfusepath{stroke,fill} %
    \pgfpathmoveto{\pgfpoint@onthreedcoil{1    }{ 1    }{15}}%
    \pgfpathcurveto%
    {\pgfpoint@onthreedcoil{1.555}{ 1    }{16}}%
    {\pgfpoint@onthreedcoil{2    }{ 0.555}{17}}%
    {\pgfpoint@onthreedcoil{2    }{ 0    }{18}}%
    \pgfcoordinate{TD@coilast}{\pgfpoint@onthreedcoil{2    }{ 0    }{18}} %
    \pgfusepath{stroke}     %
    \endgroup%
  }%
  \state{coil}[switch if less than=%%
    1.9*\pgfdecorationsegmentlength  to last,%
               width=+\pgfdecorationsegmentlength]%
    { % line in the back%
    %%
    \pgfsetstrokecolor{\pgfkeysvalueof{/pgf/decoration/3d coil color}}%
    \pgfsetfillcolor{\pgfkeysvalueof{/pgf/decoration/3d coil color}}%
    \pgfsetstrokeopacity{\pgfkeysvalueof{/pgf/decoration/3d coil opacity}}%
    \pgfpathmoveto{\pgfpointanchor{TD@coilast}{center}}%
    \pgfsetlinewidth{\pgfkeysvalueof{/pgf/decoration/3d coil width}} %
    \pgfpathcurveto%
    {\pgfpoint@onthreedcoil{2    }{-0.555}{7}}%
    {\pgfpoint@onthreedcoil{1.555}{-1    }{8}}%
    {\pgfpoint@onthreedcoil{1    }{-1    }{9}}%
    \pgfusepath{stroke} %
    %%
    % white background for front thick part%
    %%
    \pgfsetstrokeopacity{1}%
    \pgfsetstrokecolor{white}%
    \pgfsetfillcolor{white}%
    \pgfsetlinewidth{1.5*\pgfkeysvalueof{/pgf/decoration/3d coil width}+1.5*\pgfkeysvalueof{/pgf/decoration/3d coil dist}}%
    \pgfpathmoveto{\pgfpoint@onthreedcoil{1    }{ 1    }{3}}%
    \pgfpathmoveto{\pgfpoint@onthreedcoil{1    }{-1    }{9}}%
    % draw forward%
    \pgfpathcurveto%
    {\pgfpoint@onthreedcoil{0.445}{-1    }{10}}%
    {\pgfpoint@onthreedcoil{0    }{-0.555}{11.25}}%
    {\pgfpoint@onthreedcoil{0    }{ 0    }{12.5}}%
    \pgfpathcurveto%
    {\pgfpoint@onthreedcoil{0    }{ 0.555}{13.25}}%
    {\pgfpoint@onthreedcoil{0.445}{ 1    }{14.25}}%
    {\pgfpoint@onthreedcoil{1    }{ 1    }{15}}%
    % draw the curve back%
    \pgfpathcurveto%
    {\pgfpoint@onthreedcoil{0.445}{ 1    }{14}}%
    {\pgfpoint@onthreedcoil{0    }{ 0.555}{12.75}}%
    {\pgfpoint@onthreedcoil{0    }{ 0    }{11.5}}%
    \pgfpathcurveto%
    {\pgfpoint@onthreedcoil{0    }{-0.555}{10.75}}%
    {\pgfpoint@onthreedcoil{0.445}{-1    }{10}}%
    {\pgfpoint@onthreedcoil{1    }{-1    }{9}}%
    \pgfusepath{stroke,fill} %
    % %
    % draw the thick foreground path%
    %%
    \pgfsetstrokecolor{\pgfkeysvalueof{/pgf/decoration/3d coil color}}%
    \pgfsetfillcolor{\pgfkeysvalueof{/pgf/decoration/3d coil color}}%
    \pgfsetstrokeopacity{\pgfkeysvalueof{/pgf/decoration/3d coil opacity}}%
    \pgfpathmoveto{\pgfpoint@onthreedcoil{1    }{ 1    }{3}}%
    \pgfsetlinewidth{\pgfkeysvalueof{/pgf/decoration/3d coil width}} %
    % forward shifted +%
    \pgfpathmoveto{\pgfpoint@onthreedcoil{1    }{-1    }{9}}%
    \pgfpathcurveto%
    {\pgfpoint@onthreedcoil{0.445}{-1    }{10}}%
    {\pgfpoint@onthreedcoil{0    }{-0.555}{11.25}}%
    {\pgfpoint@onthreedcoil{0    }{ 0    }{12.5}}%
    \pgfpathcurveto%
    {\pgfpoint@onthreedcoil{0    }{ 0.555}{13.25}}%
    {\pgfpoint@onthreedcoil{0.445}{ 1    }{14.25}}%
    {\pgfpoint@onthreedcoil{1    }{ 1    }{15}}%
    % draw the curve back shfted -%
    \pgfpathcurveto%
    {\pgfpoint@onthreedcoil{0.445}{ 1    }{14}}%
    {\pgfpoint@onthreedcoil{0    }{ 0.555}{12.75}}%
    {\pgfpoint@onthreedcoil{0    }{ 0    }{11.5}}%
    \pgfpathcurveto%
    {\pgfpoint@onthreedcoil{0    }{-0.555}{10.75}}%
    {\pgfpoint@onthreedcoil{0.445}{-1    }{10}}%
    {\pgfpoint@onthreedcoil{1    }{-1    }{9}}%
    \pgfusepath{stroke,fill} %
    \pgfpathmoveto{\pgfpoint@onthreedcoil{1    }{ 1    }{15}}%
    \pgfpathcurveto%
    {\pgfpoint@onthreedcoil{1.555}{ 1    }{16}}%
    {\pgfpoint@onthreedcoil{2    }{ 0.555}{17}}%
    {\pgfpoint@onthreedcoil{2    }{ 0    }{18}}%
    \pgfusepath{stroke} %
    \pgfcoordinate{TD@coilast}{\pgfpoint@onthreedcoil{2    }{ 0    }{18}} %
  }%
  \state{last}[next state=final]%
    { % line in the back%
    %%
    \pgfsetstrokecolor{\pgfkeysvalueof{/pgf/decoration/3d coil color}}%
    \pgfsetfillcolor{\pgfkeysvalueof{/pgf/decoration/3d coil color}}%
    \pgfsetstrokeopacity{\pgfkeysvalueof{/pgf/decoration/3d coil opacity}}%
    \pgfpathmoveto{\pgfpointanchor{TD@coilast}{center}}%
    \pgfsetlinewidth{\pgfkeysvalueof{/pgf/decoration/3d coil width}} %
    \pgfpathcurveto%
    {\pgfpoint@onthreedcoil{2    }{-0.555}{7}}%
    {\pgfpoint@onthreedcoil{1.555}{-1    }{8}}%
    {\pgfpoint@onthreedcoil{1    }{-1    }{9}}%
    \pgfusepath{stroke} %
    % %
    % draw the thick foreground path%
    %%
    \ifcoil@closed %\pgfpointanchor{TD@coilfirst}{center}%
     %%
     % white background for front thick part%
     %%
     \pgfsetstrokeopacity{1}%
     \pgfsetstrokecolor{white}%
     \pgfsetfillcolor{white}%
     \pgfsetlinewidth{1.5*\pgfkeysvalueof{/pgf/decoration/3d coil width}+1.5*\pgfkeysvalueof{/pgf/decoration/3d coil dist}}%
     \pgfpathmoveto{\pgfpoint@onthreedcoil{1    }{ 1    }{3}}%
     \pgfpathmoveto{\pgfpoint@onthreedcoil{1    }{-1    }{9}}%
     % draw forward%
     \pgfpathcurveto%
     {\pgfpoint@onthreedcoil{0.445}{-1    }{10}}%
     {\pgfpoint@onthreedcoil{0    }{-0.555}{11.25}}%
     {\pgfpoint@onthreedcoil{0    }{ 0    }{12.5}}%
     \pgfpathcurveto%
     {\pgfpoint@onthreedcoil{0    }{ 0.555}{13.25}}%
     {\pgfpoint@onthreedcoil{0.445}{ 1    }{14.25}}%
     {\pgfpointanchor{TD@coilfirst}{center}}%
     % draw the curve back%
     \pgfpathcurveto%
     {\pgfpoint@onthreedcoil{0.445}{ 1    }{14}}%
     {\pgfpoint@onthreedcoil{0    }{ 0.555}{12.75}}%
     {\pgfpoint@onthreedcoil{0    }{ 0    }{11.5}}%
     \pgfpathcurveto%
     {\pgfpoint@onthreedcoil{0    }{-0.555}{10.75}}%
     {\pgfpoint@onthreedcoil{0.445}{-1    }{10}}%
     {\pgfpoint@onthreedcoil{1    }{-1    }{9}}%
     \pgfusepath{stroke,fill} %
     \pgfsetstrokecolor{\pgfkeysvalueof{/pgf/decoration/3d coil color}}%
     \pgfsetfillcolor{\pgfkeysvalueof{/pgf/decoration/3d coil color}}%
     \pgfsetstrokeopacity{\pgfkeysvalueof{/pgf/decoration/3d coil opacity}}%
     \pgfpathmoveto{\pgfpoint@onthreedcoil{1    }{ 1    }{3}}%
     \pgfsetlinewidth{\pgfkeysvalueof{/pgf/decoration/3d coil width}} %
     % forward shifted +%
     \pgfpathmoveto{\pgfpoint@onthreedcoil{1    }{-1    }{9}}%
     \pgfpathcurveto%
     {\pgfpoint@onthreedcoil{0.445}{-1    }{10}}%
     {\pgfpoint@onthreedcoil{0    }{-0.555}{11.25}}%
     {\pgfpoint@onthreedcoil{0    }{ 0    }{12.5}}%
     \pgfpathcurveto%
     {\pgfpoint@onthreedcoil{0    }{ 0.555}{13.25}}%
     {\pgfpoint@onthreedcoil{0.445}{ 1    }{14.25}}%
     {\pgfpointanchor{TD@coilfirst}{center}}%
     % draw the curve back shifted %
     \pgfpathcurveto%
     {\pgfpoint@onthreedcoil{0.445}{ 1    }{14}}%
     {\pgfpoint@onthreedcoil{0    }{ 0.555}{12.75}}%
     {\pgfpoint@onthreedcoil{0    }{ 0    }{11.5}}%
     \pgfpathcurveto%
     {\pgfpoint@onthreedcoil{0    }{-0.555}{10.75}}%
     {\pgfpoint@onthreedcoil{0.445}{-1    }{10}}%
     {\pgfpoint@onthreedcoil{1    }{-1    }{9}}%
     \pgfusepath{stroke,fill} %
    \else%
     %%
     % white background for front thick part%
     %%
     \pgfsetstrokeopacity{1}%
     \pgfsetstrokecolor{white}%
     \pgfsetfillcolor{white}%
     \pgfsetlinewidth{1.5*\pgfkeysvalueof{/pgf/decoration/3d coil width}+1.5*\pgfkeysvalueof{/pgf/decoration/3d coil dist}}%
     \pgfpathmoveto{\pgfpoint@onthreedcoil{1}{1}{3}}%
     \pgfpathmoveto{\pgfpoint@onthreedcoil{1}{-1}{9}}%
     % draw forward%
     \pgfpathcurveto%
     {\pgfpoint@onthreedcoil{0.445}{-1    }{10}}%
     {\pgfpoint@onthreedcoil{0    }{-0.555}{11.25}}%
     {\pgfpoint@onthreedcoil{0    }{ 0    }{12.5}}%
     \pgfpathcurveto%
     {\pgfpoint@onthreedcoil{0    }{ 0.555}{13.25}}%
     {\pgfpoint@onthreedcoil{0.445}{ 1    }{14.25}}%
     {\pgfpoint@onthreedcoil{1    }{ 1    }{15}}%
     % draw the curve back%
     \pgfpathcurveto%
     {\pgfpoint@onthreedcoil{0.445}{ 1    }{14}}%
     {\pgfpoint@onthreedcoil{0    }{ 0.555}{12.75}}%
     {\pgfpoint@onthreedcoil{0    }{ 0    }{11.5}}%
     \pgfpathcurveto%
     {\pgfpoint@onthreedcoil{0    }{-0.555}{10.75}}%
     {\pgfpoint@onthreedcoil{0.445}{-1    }{10}}%
     {\pgfpoint@onthreedcoil{1    }{-1    }{9}}%
     \pgfusepath{stroke,fill} %
     \pgfsetstrokecolor{\pgfkeysvalueof{/pgf/decoration/3d coil color}}%
     \pgfsetfillcolor{\pgfkeysvalueof{/pgf/decoration/3d coil color}}%
     \pgfsetstrokeopacity{\pgfkeysvalueof{/pgf/decoration/3d coil opacity}}%
     \pgfpathmoveto{\pgfpoint@onthreedcoil{1}{1}{3}}%
     \pgfsetlinewidth{\pgfkeysvalueof{/pgf/decoration/3d coil width}} %
     % forward shifted +%
     \pgfpathmoveto{\pgfpoint@onthreedcoil{1}{-1}{9}}%
     \pgfpathcurveto%
     {\pgfpoint@onthreedcoil{0.445}{-1    }{10}}%
     {\pgfpoint@onthreedcoil{0    }{-0.555}{11.25}}%
     {\pgfpoint@onthreedcoil{0    }{ 0    }{12.5}}%
     \pgfpathcurveto%
     {\pgfpoint@onthreedcoil{0    }{ 0.555}{13.25}}%
     {\pgfpoint@onthreedcoil{0.445}{ 1    }{14.25}}%
     {\pgfpoint@onthreedcoil{1    }{ 1    }{15}}%
     % draw the curve back shifted %
     \pgfpathcurveto%
     {\pgfpoint@onthreedcoil{0.445}{ 1    }{14}}%
     {\pgfpoint@onthreedcoil{0    }{ 0.555}{12.75}}%
     {\pgfpoint@onthreedcoil{0    }{ 0    }{11.5}}%
     \pgfpathcurveto%
     {\pgfpoint@onthreedcoil{0    }{-0.555}{10.75}}%
     {\pgfpoint@onthreedcoil{0.445}{-1    }{10}}%
     {\pgfpoint@onthreedcoil{1    }{-1    }{9}}%
     \pgfusepath{stroke,fill} %
    \fi%
    \pgfpathmoveto{\pgfpoint@onthreedcoil{1    }{ 1    }{15}}%
    \ifcoil@closed %TD@coilfirst%
    \else%
     \pgfpathcurveto%
     {\pgfpoint@onthreedcoil{1.555}{ 1    }{16}}%
     {\pgfpoint@onthreedcoil{2    }{ 0.555}{17}}%
     {\pgfpoint@onthreedcoil{2    }{ 0    }{18}}%
    \fi%
    \pgfusepath{stroke} %
    %\pgfcoordinate{TD@coilast}{\pgfpoint@onthreedcoil{2    }{ 0    }{18}} %
  }%
  \state{final}%
  {%
    \pgfpathmoveto{\pgfpointdecoratedpathlast}%
	\pgfutil@namelet{\string\color@pgfstrokecolor}{pgf@strokecolor@global}%
    %\tikzset{restore stroke color/.expand once=\pgfsavedstrokecolor}%
  }%
}%
% cone shading
% based on https://tex.stackexchange.com/a/333409/238301
\newcounter{tikz@td@cone@shading}
% x and y run from -25 to 25, alpha is an angle
\tikzset{3d/cone/shading/.cd,x/.initial=0,y/.initial=0,alpha/.initial=0,
brightness/.initial=0.5,
create shading/.code={%
\stepcounter{tikz@td@cone@shading}%
\tikzset{3d/cone/shading/.cd,#1}%
\pgfdeclarefunctionalshading{simple cone \number\value{tikz@td@cone@shading}}{\pgfpoint{-25bp}{-25bp}}{\pgfpoint{25bp}{25bp}}{}{
    % x y
	\pgfkeysvalueof{/tikz/3d/cone/shading/y} -1 mul add % shift y by the offset
	exch % swap x and y
	\pgfkeysvalueof{/tikz/3d/cone/shading/x} -1 mul add % shift x by the offset
	exch % swap y and x (reverse above swap)
    atan   % theta
	\pgfkeysvalueof{/tikz/3d/cone/shading/alpha} add % 
    sin         % cos(theta)
	1 add       % 1 + sin(theta)
    \pgfkeysvalueof{/tikz/3d/cone/shading/brightness} mul     % brightness*(1 + sin(theta))
    dup mul     % [f*(1 + sin(theta))]^2
	2 mul
	dup 1 gt {pop 1} if
	dup dup		% important so that it gets three color components
	\ifpgfshadingmodelcmyk\pgffuncshadingrgbtocmyk\fi
	\ifpgfshadingmodelgray\pgffuncshadingrgbtogray\fi
}}}
\tikzset{3d/cone shading/.style={3d/cone/shading/create shading={#1},
 shading=simple cone \number\value{tikz@td@cone@shading}}}
%% lighting
%% One directional light. 3d/light is the direction pointing towards the light,
%% given in the 3d frame in which coordinates are specified (the frame of
%% install view), so it stays fixed in the scene when the view changes.
%% Intensity of a surface element with unit outward normal n (Lambert):
%%   I = min(1, ambient + diffuse*max(0, n.L)),  colour = <color>!100I!black
\newcounter{tikz@td@lit}%
\newcount\tikz@td@k
\tikzset{3d/light/.initial={(1,1,1)},%
  3d/light/ambient/.initial=0.5,3d/light/diffuse/.initial=0.5,%
  3d/cylinder/color/.initial=white,3d/cone/color/.initial=white,%
  % painted below 3d/cylinder/mantle, so a fill or shading there wins
  3d/cylinder/lighting/.style={path picture={\tikz@td@cyl@paint}},%
  % orients a user shading of the mantle across the axis without forcing one
  3d/cylinder/shading angle/.code={\edef\tikz@shade@angle{#1}},%
  % cone hook 3d/cone/light shading: #1 = canvas angle (degrees, around the
  % tip) of the brightest visible generator; lit/max and lit/min are the
  % intensities of the brightest and darkest visible generators
  3d/cone/lit/max/.initial=1,3d/cone/lit/min/.initial=0.5,%
  3d/cone/lit side/.code={%
    \pgfkeyssetvalue{/tikz/3d/cone/lit/max}{\csname tikz@td@cone@#1@max\endcsname}%
    \pgfkeyssetvalue{/tikz/3d/cone/lit/min}{\csname tikz@td@cone@#1@min\endcsname}%
    \expandafter\let\expandafter\tikz@td@cone@ps\csname tikz@td@cone@#1@ps\endcsname
    \pgfkeysalso{/tikz/3d/cone/light shading/.expanded=\csname tikz@td@cone@#1@angle\endcsname}},%
  3d/cone/light shading/.code={\tikz@td@cone@shade{#1}}}%
\pgfmathdeclarefunction{tikztdlight}{1}{%
  \begingroup
  \pgfmathparse{max(0,min(1,\pgfkeysvalueof{/tikz/3d/light/ambient}%
    +\pgfkeysvalueof{/tikz/3d/light/diffuse}*max(0,#1)))}%
  \pgfmathsmuggle\pgfmathresult\endgroup}%
\def\tikz@td@splitvec#1,#2,#3\relax#4{%
  \expandafter\edef\csname tikz@td@l@#4x\endcsname{#1}%
  \expandafter\edef\csname tikz@td@l@#4y\endcsname{#2}%
  \expandafter\edef\csname tikz@td@l@#4z\endcsname{#3}}%
% \tikz@td@l@L{x,y,z}, \tikz@td@l@N{x,y,z}: unit light direction, unit screen normal
\def\tikz@td@getlight{%
  \edef\pgfutil@tmpa{\pgfkeysvalueof{/tikz/3d/light}}%
  \pgfmathsetmacro\pgfutil@tmpa{TDunit("\pgfutil@tmpa")}%
  \expandafter\tikz@td@splitvec\pgfutil@tmpa\relax{L}%
  \pgfmathsetmacro\pgfutil@tmpa{TDunit("(nscreenx,nscreeny,nscreenz)")}%
  \expandafter\tikz@td@splitvec\pgfutil@tmpa\relax{N}}%
% \tikz@td@canvas{<x>}{<y>}{<z>}{<name>}: canvas point -> \tikz@td@<name>x/y (pt)
\def\tikz@td@canvas#1#2#3#4{%
  \pgfpointtransformed{\pgfpointxyz{#1}{#2}{#3}}%
  \expandafter\edef\csname tikz@td@l@#4x\endcsname{\the\pgf@x}%
  \expandafter\edef\csname tikz@td@l@#4y\endcsname{\the\pgf@y}}%
% Cylinder of radius 3d/r and height 3d/h, base centre at the origin.
% #1 = unit vector across the silhouette, #2 = unit axis (x,y,z, current frame).
% In parallel projection the mantle brightness only depends on s in [-1,1],
% the position across the silhouette; the visible normal is s*e + sqrt(1-s^2)*w
% (w: normal facing the viewer), so an axial shading is exact.
% Defines \tikz@td@cyl@paint (used by 3d/cylinder/lighting).
\def\tikz@td@cylinder@light#1#2{%
  \let\tikz@td@cyl@paint\pgfutil@empty
  \tikz@td@getlight
  \edef\pgfutil@tmpa{#1}\expandafter\tikz@td@splitvec\pgfutil@tmpa\relax{e}%
  \edef\pgfutil@tmpa{#2}\expandafter\tikz@td@splitvec\pgfutil@tmpa\relax{a}%
  \pgfmathsetmacro\tikz@td@l@r{abs(\pgfkeysvalueof{/tikz/3d/r})}%
  \pgfmathsetmacro\tikz@td@l@h{\pgfkeysvalueof{/tikz/3d/h}}%
  \pgfmathsetmacro\pgfutil@tmpb{\tikz@td@l@Nx*\tikz@td@l@ax+\tikz@td@l@Ny*\tikz@td@l@ay+\tikz@td@l@Nz*\tikz@td@l@az}%
  \pgfmathsetmacro\tikz@td@l@wx{\tikz@td@l@Nx-\pgfutil@tmpb*\tikz@td@l@ax}%
  \pgfmathsetmacro\tikz@td@l@wy{\tikz@td@l@Ny-\pgfutil@tmpb*\tikz@td@l@ay}%
  \pgfmathsetmacro\tikz@td@l@wz{\tikz@td@l@Nz-\pgfutil@tmpb*\tikz@td@l@az}%
  \pgfmathsetmacro\pgfutil@tmpb{sqrt(\tikz@td@l@wx*\tikz@td@l@wx+\tikz@td@l@wy*\tikz@td@l@wy+\tikz@td@l@wz*\tikz@td@l@wz)}%
  \pgfmathsetmacro\pgfutil@tmpc{sqrt(\tikz@td@l@ex*\tikz@td@l@ex+\tikz@td@l@ey*\tikz@td@l@ey+\tikz@td@l@ez*\tikz@td@l@ez)}%
  \ifnum\fpeval{\pgfutil@tmpb>0.01}=1\relax\ifnum\fpeval{\pgfutil@tmpc>0.5}=1\relax
  \pgfmathsetmacro\tikz@td@l@Le{\tikz@td@l@Lx*\tikz@td@l@ex+\tikz@td@l@Ly*\tikz@td@l@ey+\tikz@td@l@Lz*\tikz@td@l@ez}%
  \pgfmathsetmacro\tikz@td@l@Lw{(\tikz@td@l@Lx*\tikz@td@l@wx+\tikz@td@l@Ly*\tikz@td@l@wy+\tikz@td@l@Lz*\tikz@td@l@wz)/\pgfutil@tmpb}%
  % canvas: centre c, c + r e, c + (h/2) a
  \pgfmathsetmacro\pgfutil@tmpa{\tikz@td@l@h/2}%
  \tikz@td@canvas{\pgfutil@tmpa*\tikz@td@l@ax}{\pgfutil@tmpa*\tikz@td@l@ay}{\pgfutil@tmpa*\tikz@td@l@az}{c}%
  \tikz@td@canvas{\pgfutil@tmpa*\tikz@td@l@ax+\tikz@td@l@r*\tikz@td@l@ex}%
    {\pgfutil@tmpa*\tikz@td@l@ay+\tikz@td@l@r*\tikz@td@l@ey}%
    {\pgfutil@tmpa*\tikz@td@l@az+\tikz@td@l@r*\tikz@td@l@ez}{u}%
  \tikz@td@canvas{2*\pgfutil@tmpa*\tikz@td@l@ax}{2*\pgfutil@tmpa*\tikz@td@l@ay}{2*\pgfutil@tmpa*\tikz@td@l@az}{t}%
  \pgfmathsetmacro\tikz@td@l@angle{atan2(\tikz@td@l@uy-\tikz@td@l@cy,\tikz@td@l@ux-\tikz@td@l@cx)}%
  \pgfmathsetmacro\tikz@td@l@rc{veclen(\tikz@td@l@uy-\tikz@td@l@cy,\tikz@td@l@ux-\tikz@td@l@cx)}%
  \pgfmathsetmacro\tikz@td@l@D{2*\tikz@td@l@rc+2*veclen(\tikz@td@l@ty-\tikz@td@l@cy,\tikz@td@l@tx-\tikz@td@l@cx)+1}%
  % pgf maps the central 50bp of the shading onto the 2D x 2D square below
  \pgfmathsetmacro\tikz@td@l@q{25*\tikz@td@l@rc/\tikz@td@l@D}%
  \colorlet{tikz@td@base}{\pgfkeysvalueof{/tikz/3d/cylinder/color}}%
  \tikz@td@k=0
  \loop
    \pgfmathsetmacro\tikz@td@l@s{-cos(11.25*\tikz@td@k)}%
    \pgfmathsetmacro\tikz@td@l@I{100*tikztdlight(\tikz@td@l@s*\tikz@td@l@Le+sin(11.25*\tikz@td@k)*\tikz@td@l@Lw)}%
    \pgfmathsetmacro\tikz@td@l@p{50+\tikz@td@l@q*\tikz@td@l@s}%
    \ifnum\tikz@td@k=0
      \edef\tikz@td@l@spec{color(0bp)=(tikz@td@base!\tikz@td@l@I!black);}%
    \fi
    \edef\tikz@td@l@spec{\tikz@td@l@spec color(\tikz@td@l@p bp)=(tikz@td@base!\tikz@td@l@I!black);}%
    \advance\tikz@td@k by1
  \ifnum\tikz@td@k<17 \repeat
  \stepcounter{tikz@td@lit}%
  \edef\tikz@td@l@litname{tikz@td@lit@\number\value{tikz@td@lit}}%
  \edef\pgfutil@tmpa{\noexpand\pgfdeclarehorizontalshading{\tikz@td@l@litname}{100bp}%
    {\tikz@td@l@spec color(100bp)=(tikz@td@base!\tikz@td@l@I!black)}}%
  \pgfutil@tmpa
  \pgfmathsetmacro\pgfutil@tmpa{\tikz@td@l@cx-\tikz@td@l@D}%
  \pgfmathsetmacro\pgfutil@tmpb{\tikz@td@l@cy-\tikz@td@l@D}%
  \pgfmathsetmacro\pgfutil@tmpc{\tikz@td@l@cx+\tikz@td@l@D}%
  \pgfmathsetmacro\pgfutil@tmpd{\tikz@td@l@cy+\tikz@td@l@D}%
  \edef\tikz@td@cyl@paint{\noexpand\pgftransformreset
    \noexpand\pgfpathrectanglecorners{\noexpand\pgfqpoint{\pgfutil@tmpa pt}{\pgfutil@tmpb pt}}%
      {\noexpand\pgfqpoint{\pgfutil@tmpc pt}{\pgfutil@tmpd pt}}%
    \noexpand\pgfshadepath{\tikz@td@l@litname}{\tikz@td@l@angle}%
    \noexpand\pgfusepath{discard}}%
  \fi\fi}%
% End caps are flat: the visible cap (outward normal n = +-#1, the sign that
% faces the viewer) gets the constant intensity tikztdlight(n.L).
% #1 = unit axis (x,y,z, current frame). Defines the colour tikz@td@l@cap.
\def\tikz@td@cylinder@cap#1{%
  \tikz@td@getlight
  \edef\pgfutil@tmpa{#1}\expandafter\tikz@td@splitvec\pgfutil@tmpa\relax{a}%
  \pgfmathsetmacro\pgfutil@tmpb{\tikz@td@l@Nx*\tikz@td@l@ax+\tikz@td@l@Ny*\tikz@td@l@ay+\tikz@td@l@Nz*\tikz@td@l@az<0?-1:1}%
  \pgfmathsetmacro\pgfutil@tmpb{100*tikztdlight(\pgfutil@tmpb*(\tikz@td@l@Lx*\tikz@td@l@ax+\tikz@td@l@Ly*\tikz@td@l@ay+\tikz@td@l@Lz*\tikz@td@l@az))}%
  \colorlet{tikz@td@l@cap}{\pgfkeysvalueof{/tikz/3d/cylinder/color}}%
  \colorlet{tikz@td@l@cap}{tikz@td@l@cap!\pgfutil@tmpb!black}}%
% paints the cap path \tikz@td@mantle through 3d/cylinder/lighting (below
% 3d/cylinder/top, so a fill there wins)
% \tikz@td@iffills{<style>}{<true>}{<false>}: does <style> fill or shade a path?
% (then the lighting below it would only show through a translucent fill)
\def\tikz@td@iffills#1{%
  \begingroup
    \let\tikz@mode\pgfutil@empty\tikz@mode@fillfalse\tikz@mode@shadefalse
    \tikzset{#1}\tikz@mode
    \iftikz@mode@fill\aftergroup\pgfutil@firstoftwo
    \else\iftikz@mode@shade\aftergroup\pgfutil@firstoftwo
    \else\aftergroup\pgfutil@secondoftwo\fi\fi
  \endgroup}%
\def\tikz@td@cylinder@paintcap{%
  \begingroup
  \def\tikz@td@cyl@paint{%
    \pgfpathrectanglecorners{\pgfpointanchor{path picture bounding box}{south west}}%
      {\pgfpointanchor{path picture bounding box}{north east}}%
    \pgfsetfillcolor{tikz@td@l@cap}\pgfusepath{fill}}%
  \path[even odd rule,3d/cylinder/lighting] \tikz@td@mantle;
  \endgroup}%
% Cone with base radius 3d/r, base centre #4, tip #4 + 3d/h * #3, base spanned by
% #1, #2 (orthonormal frame, x,y,z strings, current frame). Outward normal at
% azimuth p: (|h| cos p, |h| sin p, sgn(h) r)/sqrt(h^2+r^2) in that frame.
% For each side (outer: n.N>0, inner: n.N<0, normal flipped) stores
% \tikz@td@cone@<side>@angle/max/min/ps for 3d/cone/lit side.
\def\tikz@td@cone@light#1#2#3#4{%
  \tikz@td@getlight
  \edef\pgfutil@tmpa{#1}\expandafter\tikz@td@splitvec\pgfutil@tmpa\relax{f}%
  \edef\pgfutil@tmpa{#2}\expandafter\tikz@td@splitvec\pgfutil@tmpa\relax{g}%
  \edef\pgfutil@tmpa{#3}\expandafter\tikz@td@splitvec\pgfutil@tmpa\relax{a}%
  \edef\pgfutil@tmpa{#4}\expandafter\tikz@td@splitvec\pgfutil@tmpa\relax{o}%
  \pgfmathsetmacro\tikz@td@l@r{abs(\pgfkeysvalueof{/tikz/3d/r})}%
  \pgfmathsetmacro\tikz@td@l@h{\pgfkeysvalueof{/tikz/3d/h}}%
  \pgfmathsetmacro\tikz@td@l@sg{\tikz@td@l@h<0?-1:1}%
  \pgfmathsetmacro\tikz@td@l@hh{abs(\tikz@td@l@h)}%
  \pgfmathsetmacro\tikz@td@l@S{max(0.0001,veclen(\tikz@td@l@h,\tikz@td@l@r))}%
  \def\tikz@td@l@dot##1##2{\csname tikz@td@l@##1x\endcsname*\csname tikz@td@l@##2x\endcsname
    +\csname tikz@td@l@##1y\endcsname*\csname tikz@td@l@##2y\endcsname
    +\csname tikz@td@l@##1z\endcsname*\csname tikz@td@l@##2z\endcsname}%
  \pgfmathsetmacro\tikz@td@l@Lf{\tikz@td@l@dot{L}{f}}%
  \pgfmathsetmacro\tikz@td@l@Lg{\tikz@td@l@dot{L}{g}}%
  \pgfmathsetmacro\tikz@td@l@La{\tikz@td@l@dot{L}{a}}%
  \pgfmathsetmacro\tikz@td@l@Nf{\tikz@td@l@dot{N}{f}}%
  \pgfmathsetmacro\tikz@td@l@Ng{\tikz@td@l@dot{N}{g}}%
  \pgfmathsetmacro\tikz@td@l@Na{\tikz@td@l@dot{N}{a}}%
  \pgfmathsetmacro\tikz@td@l@phiL{atan2(\tikz@td@l@Lg,\tikz@td@l@Lf)}%
  \pgfmathsetmacro\tikz@td@l@phiN{atan2(\tikz@td@l@Ng,\tikz@td@l@Nf)}%
  % outer generators are visible for cos(p-phiN) > kappa
  \pgfmathsetmacro\pgfutil@tmpb{\tikz@td@l@hh*veclen(\tikz@td@l@Nf,\tikz@td@l@Ng)}%
  % pgfmath evaluates both branches of ?:, so keep |kappa|>1 (axis towards
  % the viewer, where kappa may also exceed the dimension limit) out of acos
  \pgfmathsetmacro\pgfutil@tmpc{\tikz@td@l@r*abs(\tikz@td@l@Na)}%
  \ifdim\pgfutil@tmpc pt<\pgfutil@tmpb pt\relax
    \pgfmathsetmacro\tikz@td@l@beta{acos(-\tikz@td@l@sg*\tikz@td@l@r*\tikz@td@l@Na/\pgfutil@tmpb)}%
  \else
    \pgfmathsetmacro\tikz@td@l@beta{-\tikz@td@l@sg*\tikz@td@l@Na>0?0:180}%
  \fi
  \tikz@td@canvas{\tikz@td@l@ox+\tikz@td@l@h*\tikz@td@l@ax}{\tikz@td@l@oy+\tikz@td@l@h*\tikz@td@l@ay}%
    {\tikz@td@l@oz+\tikz@td@l@h*\tikz@td@l@az}{tip}%
  \tikz@td@cone@side{outer}{1}{\tikz@td@l@phiN}{\tikz@td@l@beta}%
  \tikz@td@cone@side{inner}{-1}{\tikz@td@l@phiN+180}{180-\tikz@td@l@beta}}%
% #1 = side, #2 = +1/-1 (normal orientation), #3 = centre, #4 = half width of
% the visible azimuth arc. Brightness is constant along each generator, so on
% the canvas it is a function of the angle psi around the tip only. It is
% sampled at 25 generators of the visible arc and stored as Type 4 code that
% maps psi (on the stack) to the intensity: \tikz@td@cone@<side>@ps.
\def\tikz@td@cone@side#1#2#3#4{%
  \pgfmathsetmacro\tikz@td@l@m{#3}%
  \pgfmathsetmacro\tikz@td@l@b{#4}%
  \def\tikz@td@l@sig{#2}%
  \let\tikz@td@l@cmax\relax
  \pgfmathsetmacro\tikz@td@l@pL{#2>0?\tikz@td@l@phiL:\tikz@td@l@phiL+180}%
  \tikz@td@cone@try{\tikz@td@l@m-\tikz@td@l@b}%
  \tikz@td@cone@try{\tikz@td@l@m+\tikz@td@l@b}%
  \tikz@td@cone@tryin{\tikz@td@l@pL}%
  \tikz@td@cone@tryin{\tikz@td@l@pL+180}%
  \tikz@td@cone@psi{\tikz@td@l@pmax}%
  \expandafter\let\csname tikz@td@cone@#1@angle\endcsname\pgfmathresult
  \pgfmathsetmacro\pgfutil@tmpa{tikztdlight(\tikz@td@l@cmax)}%
  \expandafter\let\csname tikz@td@cone@#1@max\endcsname\pgfutil@tmpa
  \pgfmathsetmacro\pgfutil@tmpa{tikztdlight(\tikz@td@l@cmin)}%
  \expandafter\let\csname tikz@td@cone@#1@min\endcsname\pgfutil@tmpa
  % profile: t_j = psi_j - psi_0 (orientation fixed by \tikz@td@l@dir), unwrapped
  \tikz@td@k=0
  \loop
    \pgfmathsetmacro\tikz@td@l@p{\tikz@td@l@m-\tikz@td@l@b+\tikz@td@l@b*\tikz@td@k/12}%
    \tikz@td@cone@try{\tikz@td@l@p}% sets \tikz@td@l@c
    \pgfmathsetmacro\tikz@td@l@I{tikztdlight(\tikz@td@l@c)}%
    \tikz@td@cone@psi{\tikz@td@l@p}%
    \ifnum\tikz@td@k=0
      \let\tikz@td@l@psiz\pgfmathresult
      \def\tikz@td@l@t{0}\def\tikz@td@l@dir{1}%
      \let\tikz@td@l@Iz\tikz@td@l@I
      \let\tikz@td@l@ps\pgfutil@empty
    \else
      \pgfmathsetmacro\pgfutil@tmpa{\pgfmathresult-\tikz@td@l@psiz}%
      \ifnum\tikz@td@k=1
        \pgfmathsetmacro\tikz@td@l@dir{mod(\pgfutil@tmpa+720,360)>180?-1:1}%
      \fi
      \pgfmathsetmacro\pgfutil@tmpa{mod(\tikz@td@l@dir*\pgfutil@tmpa+720,360)}%
      \pgfmathsetmacro\pgfutil@tmpa{\pgfutil@tmpa<\tikz@td@l@t-0.001?\pgfutil@tmpa+360:\pgfutil@tmpa}%
      \pgfmathsetmacro\pgfutil@tmpb{\pgfutil@tmpa-\tikz@td@l@t}%
      \ifnum\fpeval{\pgfutil@tmpb>0.0001}=1\relax
        \pgfmathsetmacro\pgfutil@tmpc{(\tikz@td@l@I-\tikz@td@l@Ilast)/\pgfutil@tmpb}%
        % acc += slope*clip(t - t_{j-1}, 0, dt); stack: t acc
        \edef\tikz@td@l@ps{\tikz@td@l@ps\space 1 index \tikz@td@l@t\space sub
          dup 0 lt {pop 0} if dup \pgfutil@tmpb\space gt {pop \pgfutil@tmpb} if
          \pgfutil@tmpc\space mul add}%
        \let\tikz@td@l@t\pgfutil@tmpa
      \fi
    \fi
    \let\tikz@td@l@Ilast\tikz@td@l@I
    \advance\tikz@td@k by1
  \ifnum\tikz@td@k<25 \repeat
  % centre the unused angular gap around the arc so that pixels just outside
  % the silhouette clamp to the nearest end
  \pgfmathsetmacro\pgfutil@tmpa{max(0,(360-\tikz@td@l@t)/2)}%
  \pgfmathsetmacro\pgfutil@tmpb{mod(\tikz@td@l@psiz-\tikz@td@l@dir*\pgfutil@tmpa+720,360)}%
  \expandafter\edef\csname tikz@td@cone@#1@ps\endcsname{%
    \pgfutil@tmpb\space sub \ifdim\tikz@td@l@dir pt<0pt neg \fi
    dup 0 lt {360 add} if
    \pgfutil@tmpa\space sub \tikz@td@l@Iz\space\tikz@td@l@ps\space exch pop}}%
\def\tikz@td@cone@try#1{%
  \pgfmathsetmacro\tikz@td@l@p{#1}%
  \pgfmathsetmacro\tikz@td@l@c{\tikz@td@l@sig*(\tikz@td@l@hh*(cos(\tikz@td@l@p)*\tikz@td@l@Lf
    +sin(\tikz@td@l@p)*\tikz@td@l@Lg)+\tikz@td@l@sg*\tikz@td@l@r*\tikz@td@l@La)/\tikz@td@l@S}%
  \ifx\tikz@td@l@cmax\relax
    \let\tikz@td@l@cmax\tikz@td@l@c\let\tikz@td@l@pmax\tikz@td@l@p
    \let\tikz@td@l@cmin\tikz@td@l@c\let\tikz@td@l@pmin\tikz@td@l@p
  \else
    \ifdim\tikz@td@l@c pt>\tikz@td@l@cmax pt\relax
      \let\tikz@td@l@cmax\tikz@td@l@c\let\tikz@td@l@pmax\tikz@td@l@p\fi
    \ifdim\tikz@td@l@c pt<\tikz@td@l@cmin pt\relax
      \let\tikz@td@l@cmin\tikz@td@l@c\let\tikz@td@l@pmin\tikz@td@l@p\fi
  \fi}%
\def\tikz@td@cone@tryin#1{%
  \pgfmathsetmacro\pgfutil@tmpa{#1-\tikz@td@l@m}%
  \pgfmathsetmacro\pgfutil@tmpa{abs(\pgfutil@tmpa-360*round(\pgfutil@tmpa/360))}%
  \ifdim\pgfutil@tmpa pt>\tikz@td@l@b pt\else\tikz@td@cone@try{#1}\fi}%
% canvas angle of the generator from the tip to the base point at azimuth #1
\def\tikz@td@cone@psi#1{%
  \pgfmathsetmacro\pgfutil@tmpb{\tikz@td@l@r*cos(#1)}%
  \pgfmathsetmacro\pgfutil@tmpc{\tikz@td@l@r*sin(#1)}%
  \tikz@td@canvas{\tikz@td@l@ox+\pgfutil@tmpb*\tikz@td@l@fx+\pgfutil@tmpc*\tikz@td@l@gx}%
    {\tikz@td@l@oy+\pgfutil@tmpb*\tikz@td@l@fy+\pgfutil@tmpc*\tikz@td@l@gy}%
    {\tikz@td@l@oz+\pgfutil@tmpb*\tikz@td@l@fz+\pgfutil@tmpc*\tikz@td@l@gz}{bp}%
  \pgfmathparse{atan2(\tikz@td@l@bpy-\tikz@td@l@tipy,\tikz@td@l@bpx-\tikz@td@l@tipx)}}%
% default 3d/cone/light shading: functional shading in the angle around the tip
% using the exact profile of the current side (#1 is not needed here)
\def\tikz@td@cone@shade#1{%
  \pgfshadecolortorgb{\pgfkeysvalueof{/tikz/3d/cone/color}}{\tikz@td@l@rgb}%
  \stepcounter{tikz@td@lit}%
  \edef\tikz@td@l@litname{tikz@td@lit@\number\value{tikz@td@lit}}%
  \pgfdeclarefunctionalshading{\tikz@td@l@litname}{\pgfpoint{-25bp}{-25bp}}%
   {\pgfpoint{25bp}{25bp}}{}{%
    exch 0.00001 add atan \tikz@td@cone@ps\space
    dup 0 lt {pop 0} if dup 1 gt {pop 1} if
    dup \tikz@td@l@rgbred mul exch dup \tikz@td@l@rgbgreen mul exch \tikz@td@l@rgbblue mul
    \ifpgfshadingmodelcmyk\pgffuncshadingrgbtocmyk\fi
    \ifpgfshadingmodelgray\pgffuncshadingrgbtogray\fi}%
  \pgfkeysalso{/tikz/shading/.expanded=\tikz@td@l@litname}}%
%%
%% scenes: automatic occlusion
%% Inside a scope with 3d/scene every path is a piece. TikZ builds it as usual
%% (bounding box, names and nodes are unaffected), but its output goes to a box
%% of its own. At the end of the scope the boxes are emitted from back to front
%% (painter's algorithm; the sort is stable). The depth of a path is the mean
%% depth of the points it visits; the pics of the library split themselves into
%% pieces with explicit depths (\tikz@td@piece). A piece also records a point
%% (canvas position and depth) and a canvas box around it, which serve for
%% - ties: pieces of equal depth (e.g. coplanar caps) are painted larger first
%%   (extent \tikz@td@kext);
%% - straight paths whose depth varies by more than 3d/split: the box is emitted
%%   in runs, clipped to bands across the path, cut at the depths of the pieces
%%   near it, so the path can pass through them;
%% - spheres (3d/sphere): what is inside stays between the back and the front
%%   shell, what is on the surface goes just in front of the shell it lies on,
%%   what is outside goes behind or in front of the whole sphere.
%% A piece goes to the layer that was active when it began. The graphic state
%% set by scopes inside the scene (line width, dash, colors, opacity, cap,
%% join) is replayed in every piece; clips and transparency groups of such
%% scopes are not.
\newif\iftikz@td@scene % local: inside 3d/scene
\newif\iftikz@td@explicit % local: inside an explicit piece, no automatic ones
\newif\iftikz@td@pfill % global: the current path fills or shades
\newcount\tikz@td@npieces
\let\tikz@td@piece@idx\relax
\let\tikz@td@scene@state\pgfutil@empty
\def\tikz@td@kext{0}%
\def\tikz@td@knosphere{0}% 1: the piece is placed with respect to a sphere already
\let\tikz@td@orig@command@path\tikz@@command@path
\def\tikz@td@scene@path{%
  \iftikz@td@explicit\else\ifx\pgfpictureid\tikz@td@scene@pic % not in a nested picture
    \global\tikz@td@pfillfalse
    \tikz@td@piece@open{}%
    \def\tikz@path@do@at@end{\tikz@td@piece@close@path\tikz@path@do@at@end}%
  \fi\fi
  \tikz@td@orig@command@path}%
\let\tikz@td@orig@finish\tikz@finish
\def\tikz@finish{%
  \iftikz@td@scene\begingroup\tikz@mode
    \iftikz@mode@fill\global\tikz@td@pfilltrue\fi
    \iftikz@mode@shade\global\tikz@td@pfilltrue\fi
  \endgroup\fi
  \tikz@td@orig@finish}%
% \tikz@td@depthof{<spec>}: \tikz@td@depthval is empty for an empty <spec>,
% front, back, or the absolute depth (pt, larger is closer to the viewer) of
% the point <spec> = (<coordinate>) in the current coordinate system; then
% \tikz@td@depthpt is {canvas x}{canvas y}{depth} of the point (else empty)
\def\tikz@td@depthof#1{%
  \pgfutil@in@({#1}%
  \ifpgfutil@in@
    \begingroup
    \tikz@scan@one@point\pgfutil@firstofone#1\relax
    \xdef\tikz@td@gtmpa{\fpeval{round(\pgf@sys@tonumber\tikz@td@z
      +\pgf@sys@tonumber\tikz@td@sd,3)}}%
    \pgfpointtransformed{\pgfqpoint{\pgf@x}{\pgf@y}}%
    \xdef\tikz@td@gtmpb{{\pgf@sys@tonumber\pgf@x}{\pgf@sys@tonumber\pgf@y}{\tikz@td@gtmpa}}%
    \endgroup
    \let\tikz@td@depthval\tikz@td@gtmpa\let\tikz@td@depthpt\tikz@td@gtmpb
  \else
    \edef\tikz@td@depthval{#1}\let\tikz@td@depthpt\pgfutil@empty
  \fi}%
% \tikz@td@piece@open{<spec>} ... \tikz@td@piece@close: everything in between
% goes to a new piece; an empty <spec> means the depth is accumulated
\def\tikz@td@piece@open#1{%
  \tikz@td@depthof{#1}%
  \ifx\tikz@td@depthval\pgfutil@empty
    \ifdefined\tikz@td@depthkey\let\tikz@td@depthval\tikz@td@depthkey\fi
  \fi
  \global\advance\tikz@td@npieces by1\relax
  \edef\tikz@td@tmpidx{\the\tikz@td@npieces}%
  \ifcsname tikz@td@pbox@\tikz@td@tmpidx\endcsname\else
    \expandafter\newbox\csname tikz@td@pbox@\tikz@td@tmpidx\endcsname
  \fi
  \expandafter\xdef\csname tikz@td@pacc@\tikz@td@tmpidx\endcsname{0,0}%
  \expandafter\gdef\csname tikz@td@ppts@\tikz@td@tmpidx\endcsname{}%
  \expandafter\xdef\csname tikz@td@prep@\tikz@td@tmpidx\endcsname{\tikz@td@depthpt}%
  \expandafter\xdef\csname tikz@td@pext@\tikz@td@tmpidx\endcsname{\tikz@td@kext}%
  \expandafter\xdef\csname tikz@td@pns@\tikz@td@tmpidx\endcsname{\tikz@td@knosphere}%
  \expandafter\xdef\csname tikz@td@pfl@\tikz@td@tmpidx\endcsname{1}%
  \expandafter\xdef\csname tikz@td@pdep@\tikz@td@tmpidx\endcsname{\tikz@td@depthval}%
  \expandafter\xdef\csname tikz@td@play@\tikz@td@tmpidx\endcsname{\pgfonlayer@name}%
  \global\expandafter\setbox\csname tikz@td@pbox@\tikz@td@tmpidx\endcsname=\hbox to 0pt\bgroup
    \let\tikz@td@piece@idx\tikz@td@tmpidx
    \pgfsys@beginscope\tikz@td@scene@state
    \begingroup}%
\def\tikz@td@piece@close{%
    \endgroup
    \pgfsys@endscope
    \hss
  \expandafter\egroup\expandafter\tikz@td@piece@register\expandafter{\tikz@td@piece@idx}}%
% an automatic piece also records whether its path fills or shades
\def\tikz@td@piece@close@path{%
    \endgroup
    \pgfsys@endscope
    \hss
  \expandafter\egroup\expandafter\tikz@td@piece@register@path\expandafter{\tikz@td@piece@idx}}%
\def\tikz@td@piece@register@path#1{%
  \expandafter\xdef\csname tikz@td@pfl@#1\endcsname{\iftikz@td@pfill1\else0\fi}%
  \tikz@td@piece@register{#1}}%
% \tikz@td@piece{<spec>}{<code>}: <code> as one piece (paths in it do not
% become pieces of their own); outside a scene just <code>
\long\def\tikz@td@piece#1#2{%
  \iftikz@td@scene
    \tikz@td@piece@open{#1}\tikz@td@explicittrue#2\tikz@td@piece@close
  \else#2\fi}%
% the points of a path accumulate their depths in the current piece, and
% their canvas positions and depths in its list of points
\let\tikz@td@noscene@makelast\tikz@make@last@position
\def\tikz@make@last@position#1{\tikz@td@noscene@makelast{#1}%
  \ifx\tikz@td@piece@idx\relax\else\iftikz@td@zok
    \expandafter\expandafter\expandafter\tikz@td@piece@acc
    \csname tikz@td@pacc@\tikz@td@piece@idx\endcsname\relax
  \fi\fi}%
% With x and y along the canvas axes (3d/screen coords) the drawing is flat,
% and its literal points lie in the screen plane through the origin; named
% coordinates keep their depth.
\def\tikz@td@piece@acc#1,#2\relax{%
  \begingroup
  \dimen@=\tikz@td@sd
  \advance\dimen@\tikz@td@z
  \iftikz@td@simple
    \ifdim\pgf@xx=1cm\ifdim\pgf@yy=1cm\ifdim\pgf@xy=0pt\ifdim\pgf@yx=0pt
      \dimen@=\tikz@td@sd
    \fi\fi\fi\fi
  \fi
  \expandafter\xdef\csname tikz@td@pacc@\tikz@td@piece@idx\endcsname{%
    \fpeval{#1+\pgf@sys@tonumber\dimen@},\the\numexpr#2+1\relax}%
  \pgfpointtransformed{\pgfqpoint{\tikz@lastx}{\tikz@lasty}}%
  \expandafter\xdef\csname tikz@td@ppts@\tikz@td@piece@idx\endcsname{%
    \csname tikz@td@ppts@\tikz@td@piece@idx\endcsname
    {{\pgf@sys@tonumber\pgf@x}{\pgf@sys@tonumber\pgf@y}{\pgf@sys@tonumber\dimen@}}}%
  \endgroup}%
% pieces without a depth (e.g. a lone node at a named coordinate) go in front
\def\tikz@td@piece@register#1{%
  \edef\tikz@td@pdepth{\csname tikz@td@pdep@#1\endcsname}%
  \ifx\tikz@td@pdepth\pgfutil@empty
    \expandafter\expandafter\expandafter\tikz@td@piece@mean
    \csname tikz@td@pacc@#1\endcsname\relax
  \fi
  \tikz@td@piece@gput{\tikz@td@pdepth}{#1}{\csname tikz@td@pext@#1\endcsname}}%
\def\tikz@td@piece@mean#1,#2\relax{%
  \ifnum#2=0 \def\tikz@td@pdepth{front}\else
  \edef\tikz@td@pdepth{\fpeval{round((#1)/#2,3)}}\fi}%
% graphic state set by scopes inside the scene
\def\tikz@td@scene@record#1{\iftikz@td@scene
  \expandafter\def\expandafter\tikz@td@scene@state\expandafter{\tikz@td@scene@state#1}\fi}%
\let\tikz@td@orig@pgfsetlinewidth\pgfsetlinewidth
\def\pgfsetlinewidth#1{\tikz@td@orig@pgfsetlinewidth{#1}%
  \expandafter\tikz@td@scene@record\expandafter{\expandafter
  \pgfsys@setlinewidth\expandafter{\the\pgflinewidth}}\ignorespaces}%
\let\tikz@td@orig@pgfsetdash\pgfsetdash
\def\pgfsetdash#1#2{\tikz@td@orig@pgfsetdash{#1}{#2}%
  \tikz@td@scene@record{\pgfsetdash{#1}{#2}}\ignorespaces}%
% \tikz@td@wrap{<name>}{<parameters>}{<arguments>}: \<name> also records itself
\def\tikz@td@wrap#1#2#3{%
  \expandafter\let\csname tikz@td@orig@#1\expandafter\endcsname\csname#1\endcsname
  \expandafter\edef\csname#1\endcsname#2{%
    \expandafter\noexpand\csname tikz@td@orig@#1\endcsname#3%
    \noexpand\tikz@td@scene@record{\expandafter\noexpand\csname tikz@td@orig@#1\endcsname#3}}}%
\tikz@td@wrap{pgf@setcolor}{#1#2#3#4#5}{{#1}{#2}{#3}{#4}{#5}}%
\tikz@td@wrap{pgf@setstrokecolor}{#1#2#3#4#5}{{#1}{#2}{#3}{#4}{#5}}%
\tikz@td@wrap{pgf@setfillcolor}{#1#2#3#4#5}{{#1}{#2}{#3}{#4}{#5}}%
% a bare color option of a scope uses the color stack of the driver, which is
% popped before the pieces are emitted
\let\tikz@td@orig@compat@color@set\tikz@compat@color@set
\def\tikz@compat@color@set#1{\tikz@td@orig@compat@color@set{#1}%
  \expandafter\tikz@td@scene@record\expandafter{\expandafter
  \tikz@td@orig@pgf@setcolor\pgf@temp}}%
\tikz@td@wrap{pgfsetstrokeopacity}{#1}{{#1}}%
\tikz@td@wrap{pgfsetfillopacity}{#1}{{#1}}%
\tikz@td@wrap{pgfsetbuttcap}{}{}\tikz@td@wrap{pgfsetroundcap}{}{}%
\tikz@td@wrap{pgfsetrectcap}{}{}\tikz@td@wrap{pgfsetmiterjoin}{}{}%
\tikz@td@wrap{pgfsetbeveljoin}{}{}\tikz@td@wrap{pgfsetroundjoin}{}{}%
\ExplSyntaxOn
\cs_generate_variant:Nn \tl_gset:Nn { ce }
\tl_new:N \l__tdtools_line_tl
\seq_new:N \g__tdtools_pieces_seq
\seq_new:N \g__tdtools_spheres_seq
\seq_new:N \l__tdtools_out_seq
\seq_new:N \l__tdtools_lines_seq
\seq_new:N \l__tdtools_base_seq
\seq_new:N \l__tdtools_lcuts_seq
\seq_new:N \l__tdtools_lpts_seq
\tl_new:N \l__tdtools_key_tl
\tl_new:N \l__tdtools_lay_tl
\cs_new_protected:Npn \tikz@td@piece@gput #1#2#3
  { \seq_gput_right:Ne \g__tdtools_pieces_seq { {#1} {#2} {#3} } }
\cs_new_protected:Npn \tikz@td@sphere@gput #1
  { \seq_gput_right:Ne \g__tdtools_spheres_seq { #1 } }
\cs_new:Npn \__tdtools_key:n #1
  { \str_case:nnF {#1} { { front } { 16000 } { back } { -16000 } } {#1} }
\cs_new:Npn \__tdtools_item_key:nnnnnn #1#2#3#4#5#6 { \__tdtools_key:n {#1} }
\cs_new:Npn \__tdtools_item_ext:nnnnnn #1#2#3#4#5#6 {#3}
\cs_new:Npn \__tdtools_item_idx:nnnnnn #1#2#3#4#5#6 {#2}
% --- geometry of a piece: tikz@td@pr@<i> = {x}{y}{depth}, tikz@td@pb@<i> =
% {x0}{y0}{x1}{y1} (canvas, pt), both empty if unknown; tikz@td@pl@<i> =
% {x0}{y0}{d0}{x1}{y1}{d1} for a straight path to be cut into runs
\fp_new:N \l__tdtools_gx_fp \fp_new:N \l__tdtools_gy_fp \fp_new:N \l__tdtools_gd_fp
\fp_new:N \l__tdtools_gxa_fp \fp_new:N \l__tdtools_gya_fp
\fp_new:N \l__tdtools_gxb_fp \fp_new:N \l__tdtools_gyb_fp
\int_new:N \l__tdtools_gn_int
\cs_new_protected:Npn \__tdtools_geom:n #1
  {
    \tl_gclear:c { tikz@td@pl@#1 }
    \tl_set:Ne \l__tdtools_tmp_tl { \use:c { tikz@td@prep@#1 } }
    \tl_if_empty:NTF \l__tdtools_tmp_tl
      {
        \fp_zero:N \l__tdtools_gx_fp \fp_zero:N \l__tdtools_gy_fp \fp_zero:N \l__tdtools_gd_fp
        \fp_set:Nn \l__tdtools_gxa_fp { 16000 } \fp_set:Nn \l__tdtools_gya_fp { 16000 }
        \fp_set:Nn \l__tdtools_gxb_fp { -16000 } \fp_set:Nn \l__tdtools_gyb_fp { -16000 }
        \int_zero:N \l__tdtools_gn_int
        \tl_map_inline:cn { tikz@td@ppts@#1 } { \__tdtools_geom_pt:nnn ##1 }
        \int_compare:nNnTF { \l__tdtools_gn_int } = { 0 }
          { \tl_gclear:c { tikz@td@pr@#1 } \tl_gclear:c { tikz@td@pb@#1 } }
          {
            \tl_gset:ce { tikz@td@pr@#1 }
              {
                { \fp_eval:n { \l__tdtools_gx_fp / \l__tdtools_gn_int } }
                { \fp_eval:n { \l__tdtools_gy_fp / \l__tdtools_gn_int } }
                { \fp_eval:n { \l__tdtools_gd_fp / \l__tdtools_gn_int } }
              }
            \tl_gset:ce { tikz@td@pb@#1 }
              {
                { \fp_eval:n { \l__tdtools_gxa_fp - 2 } } { \fp_eval:n { \l__tdtools_gya_fp - 2 } }
                { \fp_eval:n { \l__tdtools_gxb_fp + 2 } } { \fp_eval:n { \l__tdtools_gyb_fp + 2 } }
              }
            \__tdtools_lineckeck:n {#1}
          }
      }
      {
        \tl_gset_eq:cN { tikz@td@pr@#1 } \l__tdtools_tmp_tl
        \exp_last_unbraced:NV \__tdtools_geom_box:nnnn \l__tdtools_tmp_tl {#1}
      }
  }
\cs_new_protected:Npn \__tdtools_geom_box:nnnn #1#2#3#4
  {
    \fp_set:Nn \l_tmpa_fp { \use:c { tikz@td@pext@#4 } * 28.45274 + 2 }
    \tl_gset:ce { tikz@td@pb@#4 }
      {
        { \fp_eval:n { #1 - \l_tmpa_fp } } { \fp_eval:n { #2 - \l_tmpa_fp } }
        { \fp_eval:n { #1 + \l_tmpa_fp } } { \fp_eval:n { #2 + \l_tmpa_fp } }
      }
  }
\cs_new_protected:Npn \__tdtools_geom_pt:nnn #1#2#3
  {
    \int_incr:N \l__tdtools_gn_int
    \fp_add:Nn \l__tdtools_gx_fp {#1} \fp_add:Nn \l__tdtools_gy_fp {#2} \fp_add:Nn \l__tdtools_gd_fp {#3}
    \fp_set:Nn \l__tdtools_gxa_fp { min ( \l__tdtools_gxa_fp , #1 ) }
    \fp_set:Nn \l__tdtools_gya_fp { min ( \l__tdtools_gya_fp , #2 ) }
    \fp_set:Nn \l__tdtools_gxb_fp { max ( \l__tdtools_gxb_fp , #1 ) }
    \fp_set:Nn \l__tdtools_gyb_fp { max ( \l__tdtools_gyb_fp , #2 ) }
  }
% a path to be cut into runs: automatic, no fill, all points on the 3d line
% through the first and the last one (within 0.5pt), depth varying by more
% than 3d/split
\fp_new:N \l__tdtools_sx_fp
\fp_new:N \l__tdtools_sy_fp
\fp_new:N \l__tdtools_sd_fp
\cs_new_protected:Npn \__tdtools_lineckeck:n #1
  {
    \tl_set:Ne \l__tdtools_key_tl { \use:c { tikz@td@pdep@#1 } }
    \bool_lazy_or:nnF
      { \str_if_eq_p:ee { \use:c { tikz@td@pfl@#1 } } { 1 } }
      { ! \tl_if_empty_p:N \l__tdtools_key_tl }
      {
        \seq_clear:N \l__tdtools_lpts_seq
        \tl_map_inline:cn { tikz@td@ppts@#1 } { \seq_put_right:Nn \l__tdtools_lpts_seq {##1} }
        \int_compare:nNnT { \seq_count:N \l__tdtools_lpts_seq } > { 1 }
          {
            \__tdtools_linecheck:een
              { \seq_item:Nn \l__tdtools_lpts_seq { 1 } }
              { \seq_item:Nn \l__tdtools_lpts_seq { -1 } } {#1}
          }
      }
  }
\cs_new_protected:Npn \__tdtools_linecheck:nnn #1#2#3
  { \__tdtools_linecheck:nnnnnnn #1 #2 {#3} }
\cs_generate_variant:Nn \__tdtools_linecheck:nnn { een }
\cs_new_protected:Npn \__tdtools_linecheck:nnnnnnn #1#2#3#4#5#6#7
  {
    \fp_set:Nn \l__tdtools_sx_fp { #4 - #1 }
    \fp_set:Nn \l__tdtools_sy_fp { #5 - #2 }
    \fp_set:Nn \l__tdtools_sd_fp { #6 - #3 }
    \bool_set_true:N \l_tmpa_bool
    \seq_map_inline:Nn \l__tdtools_lpts_seq
      { \__tdtools_linept:nnnnnn {#1} {#2} {#3} ##1 }
    \bool_lazy_all:nT
      {
        { \l_tmpa_bool }
        { \fp_compare_p:nNn { sqrt ( \l__tdtools_sx_fp ^ 2 + \l__tdtools_sy_fp ^ 2 ) } > { 1 } }
        { \fp_compare_p:nNn { abs ( \l__tdtools_sd_fp ) } >
            { \dim_to_fp:n { \pgfkeysvalueof { /tikz/3d/split } } } }
        { ! \fp_compare_p:nNn { \dim_to_fp:n { \pgfkeysvalueof { /tikz/3d/split } } } = { 0 } }
      }
      { \tl_gset:cn { tikz@td@pl@#7 } { {#1} {#2} {#3} {#4} {#5} {#6} } }
  }
\cs_new_protected:Npn \__tdtools_linept:nnnnnn #1#2#3#4#5#6
  {
    \fp_compare:nNnT
      {
        sqrt (
          ( (#5-#2) * \l__tdtools_sd_fp - (#6-#3) * \l__tdtools_sy_fp ) ^ 2
          + ( (#6-#3) * \l__tdtools_sx_fp - (#4-#1) * \l__tdtools_sd_fp ) ^ 2
          + ( (#4-#1) * \l__tdtools_sy_fp - (#5-#2) * \l__tdtools_sx_fp ) ^ 2 )
      } >
      { 0.5 * sqrt ( \l__tdtools_sx_fp ^ 2 + \l__tdtools_sy_fp ^ 2 + \l__tdtools_sd_fp ^ 2 ) }
      { \bool_set_false:N \l_tmpa_bool }
  }
% --- spheres: {x}{y}{center depth}{radius}{front depth}{back depth}
% \__tdtools_sphere:nnn {<key>} {<point>} : new key in \l__tdtools_key_tl
\cs_new_protected:Npn \__tdtools_place:nn #1#2
  {
    \tl_set:Ne \l__tdtools_key_tl {#1}
    \tl_if_empty:nF {#2}
      {
        \seq_map_inline:Nn \g__tdtools_spheres_seq
          { \__tdtools_place:nnnnnnnnn ##1 #2 }
      }
  }
% (\l__tdtools_pidx_tl: the piece; \l__tdtools_seg_tl: empty, or the
% canvas segment {xa}{ya}{xb}{yb} of a run, then the run meets the disk of
% the sphere if the segment does)
\tl_new:N \l__tdtools_pidx_tl
\tl_new:N \l__tdtools_seg_tl
\cs_new_protected:Npn \__tdtools_place:nnnnnnnnn #1#2#3#4#5#6#7#8#9
  {
    \tl_if_empty:NTF \l__tdtools_seg_tl
      { \fp_set:Nn \l_tmpa_fp { sqrt ( (#7 - #1) ^ 2 + (#8 - #2) ^ 2 ) } }
      { \exp_last_unbraced:NV \__tdtools_segdist:nnnnnn \l__tdtools_seg_tl {#1} {#2} }
    \fp_compare:nNnT { \l_tmpa_fp } < { 1.02 * #4 }
      {
        \fp_set:Nn \l_tmpa_fp { sqrt ( (#7 - #1) ^ 2 + (#8 - #2) ^ 2 ) }
        \fp_set:Nn \l_tmpb_fp { sqrt ( \l_tmpa_fp ^ 2 + (#9 - #3) ^ 2 ) }
        % runs of straight paths are cut at the surface: inside or outside
        \fp_compare:nNnF { \l_tmpb_fp }
          < { \tl_if_empty:NTF \l__tdtools_seg_tl { 0.98 } { 1 } * #4 }
          {
            \fp_compare:nNnTF { \l_tmpb_fp }
              > { \tl_if_empty:NTF \l__tdtools_seg_tl { 1.02 } { 1 } * #4 }
              {
                % outside
                \fp_compare:nNnTF {#9} < {#3}
                  { \tl_set:Ne \l__tdtools_key_tl { \fp_eval:n { min ( \l__tdtools_key_tl , #6 - 0.01 ) } } }
                  { \tl_set:Ne \l__tdtools_key_tl { \fp_eval:n { max ( \l__tdtools_key_tl , #5 + 0.02 ) } } }
              }
              {
                % on the surface (behind: in the reverse order)
                \fp_compare:nNnTF {#9} < {#3}
                  { \tl_set:Ne \l__tdtools_key_tl { \fp_eval:n { #6 + 0.01 - 0.000001 * \l__tdtools_pidx_tl } } }
                  { \tl_set:Ne \l__tdtools_key_tl { \fp_eval:n { #5 + 0.01 } } }
              }
          }
      }
  }
% \l_tmpa_fp = canvas distance of the point {#5}{#6} from the segment #1..#4
\cs_new_protected:Npn \__tdtools_segdist:nnnnnn #1#2#3#4#5#6
  {
    \fp_set:Nn \l_tmpb_fp { (#3 - #1) ^ 2 + (#4 - #2) ^ 2 }
    \fp_set:Nn \l_tmpb_fp
      {
        \l_tmpb_fp < 0.0001 ? 0 :
        max ( 0 , min ( 1 , ( (#5 - #1) * (#3 - #1) + (#6 - #2) * (#4 - #2) ) / \l_tmpb_fp ) )
      }
    \fp_set:Nn \l_tmpa_fp
      {
        sqrt ( ( #5 - #1 - \l_tmpb_fp * (#3 - #1) ) ^ 2
          + ( #6 - #2 - \l_tmpb_fp * (#4 - #2) ) ^ 2 )
      }
  }
% --- the flush
\cs_new_protected:Npn \tikz@td@scene@flush
  {
    \seq_clear:N \l__tdtools_out_seq
    \seq_clear:N \l__tdtools_lines_seq
    \seq_map_inline:Nn \g__tdtools_pieces_seq { \__tdtools_geom:n { \use_ii:nnn ##1 } }
    % pieces that are not cut: placed with respect to the spheres
    \seq_map_inline:Nn \g__tdtools_pieces_seq { \__tdtools_flush_piece:nnn ##1 }
    % straight paths: runs, cut at the depths of the other pieces near them
    \seq_set_eq:NN \l__tdtools_base_seq \l__tdtools_out_seq
    \seq_map_inline:Nn \l__tdtools_lines_seq { \__tdtools_flush_line:nnn ##1 }
    \seq_sort:Nn \l__tdtools_out_seq
      {
        \fp_compare:nNnTF
          { \__tdtools_item_key:nnnnnn ##1 } > { \__tdtools_item_key:nnnnnn ##2 }
          { \sort_return_swapped: }
          {
            % equal depths: larger first, then in the order of the code
            \fp_compare:nNnTF
              { \__tdtools_item_key:nnnnnn ##1 } = { \__tdtools_item_key:nnnnnn ##2 }
              {
                \fp_compare:nNnTF
                  { \__tdtools_item_ext:nnnnnn ##1 } < { \__tdtools_item_ext:nnnnnn ##2 }
                  { \sort_return_swapped: }
                  {
                    \bool_lazy_and:nnTF
                      { \fp_compare_p:nNn { \__tdtools_item_ext:nnnnnn ##1 } = { \__tdtools_item_ext:nnnnnn ##2 } }
                      { \int_compare_p:nNn { \__tdtools_item_idx:nnnnnn ##1 } > { \__tdtools_item_idx:nnnnnn ##2 } }
                      { \sort_return_swapped: } { \sort_return_same: }
                  }
              }
              { \sort_return_same: }
          }
      }
    \seq_map_inline:Nn \l__tdtools_out_seq { \__tdtools_emit:nnnnnn ##1 }
    \seq_map_inline:Nn \l__tdtools_lines_seq
      { \tex_global:D \tex_setbox:D \use:c { tikz@td@pbox@ \use_ii:nnn ##1 } = \tex_box:D \c_empty_box }
    \seq_gclear:N \g__tdtools_pieces_seq
    \seq_gclear:N \g__tdtools_spheres_seq
    \global \tikz@td@npieces = 0 \scan_stop:
  }
% entries of \l__tdtools_out_seq: {key}{piece}{extent}{kind}{a}{b}; kind n
% is the box, kind r a run of a straight path from a to b (fractions of it)
\cs_new_protected:Npn \__tdtools_flush_piece:nnn #1#2#3
  {
    \tl_if_empty:cTF { tikz@td@pl@#2 }
      {
        \tl_set:Ne \l__tdtools_key_tl { \__tdtools_key:n {#1} }
        \tl_set:Nn \l__tdtools_pidx_tl {#2}
        \tl_clear:N \l__tdtools_seg_tl
        \str_if_eq:eeF { \use:c { tikz@td@pns@#2 } } { 1 }
          {
            \str_case:nnF {#1} { { front } { } { back } { } }
              { \exp_args:Nee \__tdtools_place:nn { \l__tdtools_key_tl } { \use:c { tikz@td@pr@#2 } } }
          }
        \seq_put_right:Ne \l__tdtools_out_seq { { \l__tdtools_key_tl } {#2} {#3} { n } { } { } }
      }
      { \seq_put_right:Nn \l__tdtools_lines_seq { {#1} {#2} {#3} } }
  }
\cs_new_protected:Npn \__tdtools_flush_line:nnn #1#2#3
  { \exp_last_unbraced:Nv \__tdtools_flush_line:nnnnnnnn { tikz@td@pl@#2 } {#2} {#3} }
\cs_new_protected:Npn \__tdtools_flush_line:nnnnnnnn #1#2#3#4#5#6#7#8
  {
    % cuts: keys of the pieces near the path, strictly inside its depth range,
    % as fractions of the path
    \seq_clear:N \l__tdtools_lcuts_seq
    \tl_set:Nn \l__tdtools_line_tl { {#1} {#2} {#3} {#4} {#5} {#6} }
    \seq_map_inline:Nn \l__tdtools_base_seq { \__tdtools_cut:nnnnnn ##1 }
    \seq_map_inline:Nn \g__tdtools_spheres_seq { \__tdtools_cut_sphere:nnnnnn ##1 }
    \seq_sort:Nn \l__tdtools_lcuts_seq
      { \fp_compare:nNnTF {##1} > {##2} { \sort_return_swapped: } { \sort_return_same: } }
    \seq_remove_duplicates:N \l__tdtools_lcuts_seq
    \seq_put_left:Nn \l__tdtools_lcuts_seq { -1000 }
    \seq_put_right:Nn \l__tdtools_lcuts_seq { 1000 }
    \tl_clear:N \l__tdtools_prev_tl
    \seq_map_inline:Nn \l__tdtools_lcuts_seq
      {
        \tl_if_empty:NF \l__tdtools_prev_tl
          {
            % the run from prev to ##1: depth and point of its middle (within the path)
            \fp_set:Nn \l_tmpa_fp { ( max ( \l__tdtools_prev_tl , 0 ) + min ( ##1 , 1 ) ) / 2 }
            \tl_set:Ne \l__tdtools_key_tl { \fp_eval:n { #3 + \l_tmpa_fp * (#6 - #3) } }
            \tl_set:Nn \l__tdtools_pidx_tl {#7}
            \tl_set:Ne \l__tdtools_seg_tl
              {
                { \fp_eval:n { #1 + max ( \l__tdtools_prev_tl , 0 ) * (#4 - #1) } }
                { \fp_eval:n { #2 + max ( \l__tdtools_prev_tl , 0 ) * (#5 - #2) } }
                { \fp_eval:n { #1 + min ( ##1 , 1 ) * (#4 - #1) } }
                { \fp_eval:n { #2 + min ( ##1 , 1 ) * (#5 - #2) } }
              }
            \exp_args:Nee \__tdtools_place:nn { \l__tdtools_key_tl }
              {
                { \fp_eval:n { #1 + \l_tmpa_fp * (#4 - #1) } }
                { \fp_eval:n { #2 + \l_tmpa_fp * (#5 - #2) } }
                { \l__tdtools_key_tl }
              }
            \seq_put_right:Ne \l__tdtools_out_seq
              {
                { \l__tdtools_key_tl } {#7} {#8}
                { \int_compare:nNnTF { \seq_count:N \l__tdtools_lcuts_seq } = { 2 } { n } { r } }
                { \l__tdtools_prev_tl } {##1}
              }
          }
        \tl_set:Nn \l__tdtools_prev_tl {##1}
      }
  }
% an entry {key}{piece}{extent}{kind}{a}{b} of the list against the path
% \l__tdtools_line_tl
\cs_new_protected:Npn \__tdtools_cut:nnnnnn #1#2#3#4#5#6
  { \exp_last_unbraced:NV \__tdtools_cut_aux:nnnnnnnn \l__tdtools_line_tl {#1} {#2} }
\cs_new_protected:Npn \__tdtools_cut_aux:nnnnnnnn #1#2#3#4#5#6#7#8
  {
    \fp_compare:nNnT
      { ( #7 - #3 ) * ( #7 - #6 ) } < { -0.0001 }
      {
        \tl_set:Ne \l__tdtools_tmp_tl { \use:c { tikz@td@pb@#8 } }
        \bool_if:nT
          {
            \tl_if_empty_p:N \l__tdtools_tmp_tl
            || \exp_last_unbraced:NV \__tdtools_boxhit_p:nnnnnnnn \l__tdtools_tmp_tl {#1} {#2} {#4} {#5}
          }
          { \seq_put_right:Ne \l__tdtools_lcuts_seq { \fp_eval:n { ( #7 - #3 ) / ( #6 - #3 ) } } }
      }
  }
% a sphere {x}{y}{center depth}{radius}{front}{back} against the path: the
% fractions where the path passes its surface are cuts
\fp_new:N \l__tdtools_scx_fp
\fp_new:N \l__tdtools_scy_fp
\fp_new:N \l__tdtools_scd_fp
\fp_new:N \l__tdtools_scr_fp
\cs_new_protected:Npn \__tdtools_cut_sphere:nnnnnn #1#2#3#4#5#6
  {
    \fp_set:Nn \l__tdtools_scx_fp {#1}
    \fp_set:Nn \l__tdtools_scy_fp {#2}
    \fp_set:Nn \l__tdtools_scd_fp {#3}
    \fp_set:Nn \l__tdtools_scr_fp {#4}
    \exp_last_unbraced:NV \__tdtools_cut_sphere_aux:nnnnnn \l__tdtools_line_tl
  }
\cs_new_protected:Npn \__tdtools_cut_sphere_aux:nnnnnn #1#2#3#4#5#6
  {
    % |P0 - C + t (P1 - P0)|^2 = R^2, i.e. a t^2 + 2 b t + c = 0
    \fp_set:Nn \l__tdtools_gx_fp { (#4 - #1) ^ 2 + (#5 - #2) ^ 2 + (#6 - #3) ^ 2 }
    \fp_set:Nn \l__tdtools_gy_fp
      {
        (#1 - \l__tdtools_scx_fp) * (#4 - #1) + (#2 - \l__tdtools_scy_fp) * (#5 - #2)
        + (#3 - \l__tdtools_scd_fp) * (#6 - #3)
      }
    \fp_set:Nn \l__tdtools_gd_fp
      {
        (#1 - \l__tdtools_scx_fp) ^ 2 + (#2 - \l__tdtools_scy_fp) ^ 2
        + (#3 - \l__tdtools_scd_fp) ^ 2 - \l__tdtools_scr_fp ^ 2
      }
    \fp_set:Nn \l_tmpa_fp { \l__tdtools_gy_fp ^ 2 - \l__tdtools_gx_fp * \l__tdtools_gd_fp }
    \fp_compare:nNnT { \l_tmpa_fp } > { 0 }
      {
        \clist_map_inline:nn { -1 , 1 }
          {
            \fp_set:Nn \l_tmpb_fp
              { ( - \l__tdtools_gy_fp + ##1 * sqrt ( \l_tmpa_fp ) ) / \l__tdtools_gx_fp }
            \fp_compare:nNnT { 0.001 < \l_tmpb_fp < 0.999 } = { 1 }
              { \seq_put_right:Ne \l__tdtools_lcuts_seq { \fp_use:N \l_tmpb_fp } }
          }
      }
  }
% does the box {x0}{y0}{x1}{y1} meet the bounding box of the segment?
\prg_new_conditional:Npnn \__tdtools_boxhit:nnnnnnnn #1#2#3#4#5#6#7#8 { p }
  {
    \bool_lazy_all:nTF
      {
        { \fp_compare_p:nNn {#1} < { max (#5, #7) + 1 } }
        { \fp_compare_p:nNn {#3} > { min (#5, #7) - 1 } }
        { \fp_compare_p:nNn {#2} < { max (#6, #8) + 1 } }
        { \fp_compare_p:nNn {#4} > { min (#6, #8) - 1 } }
      }
      { \prg_return_true: } { \prg_return_false: }
  }
\cs_new_protected:Npn \__tdtools_emit:nnnnnn #1#2#3#4#5#6
  {
    \tl_set:Ne \l__tdtools_lay_tl { \use:c { tikz@td@play@#2 } }
    \str_if_eq:eeTF { \l__tdtools_lay_tl } { \pgfonlayer@name }
      { \__tdtools_emit_box:nnnn {#2} {#4} {#5} {#6} }
      {
        \exp_args:NV \pgfonlayer \l__tdtools_lay_tl
          \__tdtools_emit_box:nnnn {#2} {#4} {#5} {#6}
        \endpgfonlayer
      }
  }
\cs_new_protected:Npn \__tdtools_emit_box:nnnn #1#2#3#4
  {
    \str_if_eq:nnTF {#2} { n }
      { \tex_copy:D \use:c { tikz@td@pbox@#1 } }
      {
        \pgfscope
          \pgftransformreset
          \group_begin:
            \pgf@relevantforpicturesizefalse
            \exp_last_unbraced:Nv \__tdtools_runpath:nnnnnnnn { tikz@td@pl@#1 } {#3} {#4}
            \pgfusepath { clip }
          \group_end:
          \tex_copy:D \use:c { tikz@td@pbox@#1 }
        \endpgfscope
      }
  }
% the band across the path from the fraction #7 to #8 (overlapping its
% neighbours by 0.25pt; the outer runs reach far beyond the ends)
\cs_new_protected:Npn \__tdtools_runpath:nnnnnnnn #1#2#3#4#5#6#7#8
  {
    \fp_set:Nn \l__tdtools_sx_fp { #4 - #1 }
    \fp_set:Nn \l__tdtools_sy_fp { #5 - #2 }
    \fp_set:Nn \l_tmpb_fp { sqrt ( \l__tdtools_sx_fp ^ 2 + \l__tdtools_sy_fp ^ 2 ) }
    \fp_set:Nn \l_tmpa_fp
      { \fp_compare:nNnTF {#7} < { -1 } { -1000 } { #7 * \l_tmpb_fp - 0.25 } }
    \fp_set:Nn \l__tdtools_sd_fp
      { \fp_compare:nNnTF {#8} > { 2 } { \l_tmpb_fp + 1000 } { #8 * \l_tmpb_fp + 0.25 } }
    \__tdtools_runpt:nnnnn {#1} {#2} { \l_tmpa_fp } { 1000 } { moveto }
    \__tdtools_runpt:nnnnn {#1} {#2} { \l__tdtools_sd_fp } { 1000 } { lineto }
    \__tdtools_runpt:nnnnn {#1} {#2} { \l__tdtools_sd_fp } { -1000 } { lineto }
    \__tdtools_runpt:nnnnn {#1} {#2} { \l_tmpa_fp } { -1000 } { lineto }
    \pgfpathclose
  }
\cs_new_protected:Npn \__tdtools_runpt:nnnnn #1#2#3#4#5
  {
    \use:c { pgfpath#5 }
      {
        \pgfqpoint
          { \fp_eval:n { #1 + ( #3 * \l__tdtools_sx_fp - #4 * \l__tdtools_sy_fp ) / \l_tmpb_fp } pt }
          { \fp_eval:n { #2 + ( #3 * \l__tdtools_sy_fp + #4 * \l__tdtools_sx_fp ) / \l_tmpb_fp } pt }
      }
  }
\ExplSyntaxOff
\tikzset{3d/scene/.code={%
  \iftikz@td@scene\else
    \tikz@td@scenetrue
    \let\tikz@td@scene@pic\pgfpictureid
    \let\tikz@td@scene@state\pgfutil@empty
    \let\tikz@@command@path\tikz@td@scene@path
    % occlusion decides what is seen of hidden edges: draw them like visible ones
    \tikzset{3d/hidden/.style={draw,solid},execute at end scope={\tikz@td@scene@flush}}%
  \fi},%
  % straight paths whose depth varies by more than this get cut where they
  % pass other pieces
  3d/split/.initial=0.2mm,%
  % caps of sub-arcs: butt caps leave a notch where a thick stroke turns
  % sharply (a circle seen almost edge-on), round caps of translucent strokes
  % overlap; round caps from 1pt on
  3d/joint caps/.code={\tikz@addoption{\ifdim\pgflinewidth<1pt\else\pgfsetroundcap\fi}},%
  % a path style without its fill, shading and path picture
  3d/edge only/.code={\tikz@addmode{\tikz@mode@fillfalse\tikz@mode@shadefalse}%
    \let\tikz@path@picture\pgfutil@empty},%
  % depth of the pieces of the current path or scope: front, back or a point
  3d/depth/.code={%
    \tikz@td@depthof{#1}\let\tikz@td@depthkey\tikz@td@depthval
    \ifx\tikz@td@piece@idx\relax\else
      \expandafter\xdef\csname tikz@td@pdep@\tikz@td@piece@idx\endcsname{\tikz@td@depthval}%
    \fi}}%
\makeatother
\endinput
